% 2026-10-25 DRAFT (note E14): the model confidence set, the test of superior predictive ability (SPA) and StepM for
% the pooled intraday forecasts and their sign strategies.
% Generator: writeup/make_rehaul_mcs.py (REHAUL_FINAL_RUN=1; before computing anything it reproduces Table 1's every-bar
% QLIKE (results/oos_final_2014start/table1_w3000/metrics.csv to 1e-12, and the printed cells) and its DM statistics,
% the closing-half-hour QLIKE and DM numbers of Section 6 (\rhrQ<M>Pool<X>, \rhrOlsQ, \rhrAf<M>PoolDm), the sign-strategy
% Sharpe ratios of the 12 pooled forecasts at mid-quotes and bid/ask (\rhrShEnetPoolFull = 1.94 among them) and of the
% short straddle, and the paper's D6 bootstrap p-values of the six Full-minus-HAR Sharpe differences; it checks its MCS
% against src/evaluation/model_confidence_set.py, its SPA against experiments/close_kfull_testlib.spa_rc and its StepM
% against experiments/feature_mcs_spa.stepm).
% Files:   generated/rehaul_mcs_numbers.tex  (\rhmcs... macros; macro-only; no name collides with a macro of another
%                                             writeup/generated/*.tex file, checked by the generator)
%          generated/rehaul_mcs_table.tex    (Tables tab:rh_mcs and tab:rh_mcs_spa, the two new labels used below)
%          results/rehaul_mcs_2026-10-25/    (every model, elimination step, family and block length, as CSV)
% To use:  % AUTO-GENERATED by writeup/make_rehaul_mcs.py -- do not edit; re-run the script.
% \rhmcs... macros: model confidence set (Hansen, Lunde and Nason 2011), SPA (Hansen 2005) and StepM (Romano and
% Wolf 2005) on the 12 pooled forecasts of the final rerun (the three linear HAR forecasts as one) and the
% closing-interval HAR benchmark. Samples: Bar (every bar, 2014-2024, daily-averaged losses), All (16:00, 2014-2024),
% Oos (16:00, T_oos), Ch (16:00, the benchmark's span). Losses: Qlike, Mse. MCS statistics: Tmax, Trange (T_R).
% SPA / StepM families: <Sample><Q|M>Af (Full vs own HAR), <Sample><Q|M>Bm (Full vs the closing-interval HAR),
% <Mid|Ba><Short|Twin><Mean|Sr> (sign strategies on the 614 post-warmup days, vs the short straddle or the HAR twin,
% mean return or the paper's pooled Sharpe ratio, D19). p-values of StepM and SPA are one-sided (H_k: no better than the benchmark).
\providecommand{\rhmcsB}{}\renewcommand{\rhmcsB}{4{,}999}  % bootstrap replications, every MCS / SPA / StepM (accuracy and strategies)
\providecommand{\rhmcsBlock}{}\renewcommand{\rhmcsBlock}{10}  % block length in trading days (circular moving-block bootstrap), main
\providecommand{\rhmcsBlockLo}{}\renewcommand{\rhmcsBlockLo}{5}  % block length, sensitivity (shorter)
\providecommand{\rhmcsBlockHi}{}\renewcommand{\rhmcsBlockHi}{20}  % block length, sensitivity (longer)
\providecommand{\rhmcsSeed}{}\renewcommand{\rhmcsSeed}{20261025}  % random seed of the accuracy bootstrap (MCS, SPA, StepM on losses)
\providecommand{\rhmcsSeedEcon}{}\renewcommand{\rhmcsSeedEcon}{20261022}  % random seed of the strategy bootstrap (= the paper's D6 bootstrap seed)
\providecommand{\rhmcsNModels}{}\renewcommand{\rhmcsNModels}{10}  % distinct pooled forecasts in the MCS (the three linear HAR forecasts as one)
\providecommand{\rhmcsNModelsWord}{}\renewcommand{\rhmcsNModelsWord}{ten}  % rhmcsNModels as a word
\providecommand{\rhmcsNModelsBench}{}\renewcommand{\rhmcsNModelsBench}{11}  % ... with the closing-interval HAR benchmark
\providecommand{\rhmcsNModelsBenchWord}{}\renewcommand{\rhmcsNModelsBenchWord}{eleven}  % rhmcsNModelsBench as a word
\providecommand{\rhmcsNPool}{}\renewcommand{\rhmcsNPool}{12}  % pooled forecasts (6 classes x HAR, Full)
\providecommand{\rhmcsNAlt}{}\renewcommand{\rhmcsNAlt}{6}  % alternatives in every SPA / StepM family (the six full-set forecasts or strategies)
\providecommand{\rhmcsNAltWord}{}\renewcommand{\rhmcsNAltWord}{six}  % rhmcsNAlt as a word
\providecommand{\rhmcsLinGap}{}\renewcommand{\rhmcsLinGap}{5 \times 10^{-6}}  % largest relative gap between the daily every-bar QLIKE of a linear HAR forecast and the pooled OLS (math mode)
\providecommand{\rhmcsDaysAll}{}\renewcommand{\rhmcsDaysAll}{2{,}579}  % trading days 2014-2024 (Table 1's span)
\providecommand{\rhmcsDaysOos}{}\renewcommand{\rhmcsDaysOos}{1{,}082}  % trading days of T_oos
\providecommand{\rhmcsDaysCh}{}\renewcommand{\rhmcsDaysCh}{2{,}579}  % trading days of the closing-interval benchmark's span within 2014-2024
\providecommand{\rhmcsChStart}{}\renewcommand{\rhmcsChStart}{2 January 2014}  % first day of the closing-interval HAR benchmark's forecasts
\providecommand{\rhmcsNBars}{}\renewcommand{\rhmcsNBars}{33{,}527}  % bars of the every-bar sample (13 a day)
\providecommand{\rhmcsNPw}{}\renewcommand{\rhmcsNPw}{614}  % post-warmup option days of the strategy tests
\providecommand{\rhmcsLinSameTrades}{}\renewcommand{\rhmcsLinSameTrades}{yes}  % the three linear HAR sign strategies have identical daily returns (mid and bid/ask)
\providecommand{\rhmcsBarQlikeTmaxPLinHar}{}\renewcommand{\rhmcsBarQlikeTmaxPLinHar}{\ensuremath{<}0.001}  % MCS p-value, every bar, 2014--2024, Qlike, Tmax: Linear, HAR
\providecommand{\rhmcsBarQlikeTmaxPRidgeFull}{}\renewcommand{\rhmcsBarQlikeTmaxPRidgeFull}{\ensuremath{<}0.001}  % MCS p-value, every bar, 2014--2024, Qlike, Tmax: Ridge, Full
\providecommand{\rhmcsBarQlikeTmaxPLassoFull}{}\renewcommand{\rhmcsBarQlikeTmaxPLassoFull}{\ensuremath{<}0.001}  % MCS p-value, every bar, 2014--2024, Qlike, Tmax: Lasso, Full
\providecommand{\rhmcsBarQlikeTmaxPEnetFull}{}\renewcommand{\rhmcsBarQlikeTmaxPEnetFull}{\ensuremath{<}0.001}  % MCS p-value, every bar, 2014--2024, Qlike, Tmax: Elastic net, Full
\providecommand{\rhmcsBarQlikeTmaxPLgbmHar}{}\renewcommand{\rhmcsBarQlikeTmaxPLgbmHar}{\ensuremath{<}0.001}  % MCS p-value, every bar, 2014--2024, Qlike, Tmax: LightGBM, HAR
\providecommand{\rhmcsBarQlikeTmaxPLgbmFull}{}\renewcommand{\rhmcsBarQlikeTmaxPLgbmFull}{1.00}  % MCS p-value, every bar, 2014--2024, Qlike, Tmax: LightGBM, Full
\providecommand{\rhmcsBarQlikeTmaxPXgbHar}{}\renewcommand{\rhmcsBarQlikeTmaxPXgbHar}{\ensuremath{<}0.001}  % MCS p-value, every bar, 2014--2024, Qlike, Tmax: XGBoost, HAR
\providecommand{\rhmcsBarQlikeTmaxPXgbFull}{}\renewcommand{\rhmcsBarQlikeTmaxPXgbFull}{0.69}  % MCS p-value, every bar, 2014--2024, Qlike, Tmax: XGBoost, Full
\providecommand{\rhmcsBarQlikeTmaxPRfHar}{}\renewcommand{\rhmcsBarQlikeTmaxPRfHar}{\ensuremath{<}0.001}  % MCS p-value, every bar, 2014--2024, Qlike, Tmax: Random forest, HAR
\providecommand{\rhmcsBarQlikeTmaxPRfFull}{}\renewcommand{\rhmcsBarQlikeTmaxPRfFull}{\ensuremath{<}0.001}  % MCS p-value, every bar, 2014--2024, Qlike, Tmax: Random forest, Full
\providecommand{\rhmcsBarQlikeTmaxNNinety}{}\renewcommand{\rhmcsBarQlikeTmaxNNinety}{2}  % models in the 90 % MCS (of 10)
\providecommand{\rhmcsBarQlikeTmaxSetNinety}{}\renewcommand{\rhmcsBarQlikeTmaxSetNinety}{LightGBM on the full predictor set and XGBoost on the full predictor set}  % the 90 % MCS, best (last eliminated) first
\providecommand{\rhmcsBarQlikeTmaxStableNinety}{}\renewcommand{\rhmcsBarQlikeTmaxStableNinety}{yes}  % the 90 % MCS is the same at block lengths (5, 20)
\providecommand{\rhmcsBarQlikeTmaxPMaxOutNinety}{}\renewcommand{\rhmcsBarQlikeTmaxPMaxOutNinety}{\ensuremath{<}0.001}  % largest MCS p-value outside the 90 % set
\providecommand{\rhmcsBarQlikeTmaxNNinetyLo}{}\renewcommand{\rhmcsBarQlikeTmaxNNinetyLo}{2}  % ... at block length 5
\providecommand{\rhmcsBarQlikeTmaxSetNinetyLo}{}\renewcommand{\rhmcsBarQlikeTmaxSetNinetyLo}{LightGBM on the full predictor set and XGBoost on the full predictor set}  % ... at block length 5
\providecommand{\rhmcsBarQlikeTmaxNNinetyHi}{}\renewcommand{\rhmcsBarQlikeTmaxNNinetyHi}{2}  % ... at block length 20
\providecommand{\rhmcsBarQlikeTmaxSetNinetyHi}{}\renewcommand{\rhmcsBarQlikeTmaxSetNinetyHi}{LightGBM on the full predictor set and XGBoost on the full predictor set}  % ... at block length 20
\providecommand{\rhmcsBarQlikeTmaxNSevenFive}{}\renewcommand{\rhmcsBarQlikeTmaxNSevenFive}{2}  % models in the 75 % MCS (of 10)
\providecommand{\rhmcsBarQlikeTmaxSetSevenFive}{}\renewcommand{\rhmcsBarQlikeTmaxSetSevenFive}{LightGBM on the full predictor set and XGBoost on the full predictor set}  % the 75 % MCS, best (last eliminated) first
\providecommand{\rhmcsBarQlikeTmaxStableSevenFive}{}\renewcommand{\rhmcsBarQlikeTmaxStableSevenFive}{yes}  % the 75 % MCS is the same at block lengths (5, 20)
\providecommand{\rhmcsBarQlikeTmaxPMaxOutSevenFive}{}\renewcommand{\rhmcsBarQlikeTmaxPMaxOutSevenFive}{\ensuremath{<}0.001}  % largest MCS p-value outside the 75 % set
\providecommand{\rhmcsBarQlikeTmaxNSevenFiveLo}{}\renewcommand{\rhmcsBarQlikeTmaxNSevenFiveLo}{2}  % ... at block length 5
\providecommand{\rhmcsBarQlikeTmaxSetSevenFiveLo}{}\renewcommand{\rhmcsBarQlikeTmaxSetSevenFiveLo}{LightGBM on the full predictor set and XGBoost on the full predictor set}  % ... at block length 5
\providecommand{\rhmcsBarQlikeTmaxNSevenFiveHi}{}\renewcommand{\rhmcsBarQlikeTmaxNSevenFiveHi}{2}  % ... at block length 20
\providecommand{\rhmcsBarQlikeTmaxSetSevenFiveHi}{}\renewcommand{\rhmcsBarQlikeTmaxSetSevenFiveHi}{LightGBM on the full predictor set and XGBoost on the full predictor set}  % ... at block length 20
\providecommand{\rhmcsBarQlikeTmaxBest}{}\renewcommand{\rhmcsBarQlikeTmaxBest}{LightGBM on the full predictor set}  % the last model left (MCS p-value 1)
\providecommand{\rhmcsBarQlikeTrangePLinHar}{}\renewcommand{\rhmcsBarQlikeTrangePLinHar}{\ensuremath{<}0.001}  % MCS p-value, every bar, 2014--2024, Qlike, Trange: Linear, HAR
\providecommand{\rhmcsBarQlikeTrangePRidgeFull}{}\renewcommand{\rhmcsBarQlikeTrangePRidgeFull}{\ensuremath{<}0.001}  % MCS p-value, every bar, 2014--2024, Qlike, Trange: Ridge, Full
\providecommand{\rhmcsBarQlikeTrangePLassoFull}{}\renewcommand{\rhmcsBarQlikeTrangePLassoFull}{\ensuremath{<}0.001}  % MCS p-value, every bar, 2014--2024, Qlike, Trange: Lasso, Full
\providecommand{\rhmcsBarQlikeTrangePEnetFull}{}\renewcommand{\rhmcsBarQlikeTrangePEnetFull}{\ensuremath{<}0.001}  % MCS p-value, every bar, 2014--2024, Qlike, Trange: Elastic net, Full
\providecommand{\rhmcsBarQlikeTrangePLgbmHar}{}\renewcommand{\rhmcsBarQlikeTrangePLgbmHar}{\ensuremath{<}0.001}  % MCS p-value, every bar, 2014--2024, Qlike, Trange: LightGBM, HAR
\providecommand{\rhmcsBarQlikeTrangePLgbmFull}{}\renewcommand{\rhmcsBarQlikeTrangePLgbmFull}{1.00}  % MCS p-value, every bar, 2014--2024, Qlike, Trange: LightGBM, Full
\providecommand{\rhmcsBarQlikeTrangePXgbHar}{}\renewcommand{\rhmcsBarQlikeTrangePXgbHar}{\ensuremath{<}0.001}  % MCS p-value, every bar, 2014--2024, Qlike, Trange: XGBoost, HAR
\providecommand{\rhmcsBarQlikeTrangePXgbFull}{}\renewcommand{\rhmcsBarQlikeTrangePXgbFull}{0.69}  % MCS p-value, every bar, 2014--2024, Qlike, Trange: XGBoost, Full
\providecommand{\rhmcsBarQlikeTrangePRfHar}{}\renewcommand{\rhmcsBarQlikeTrangePRfHar}{\ensuremath{<}0.001}  % MCS p-value, every bar, 2014--2024, Qlike, Trange: Random forest, HAR
\providecommand{\rhmcsBarQlikeTrangePRfFull}{}\renewcommand{\rhmcsBarQlikeTrangePRfFull}{\ensuremath{<}0.001}  % MCS p-value, every bar, 2014--2024, Qlike, Trange: Random forest, Full
\providecommand{\rhmcsBarQlikeTrangeNNinety}{}\renewcommand{\rhmcsBarQlikeTrangeNNinety}{2}  % models in the 90 % MCS (of 10)
\providecommand{\rhmcsBarQlikeTrangeSetNinety}{}\renewcommand{\rhmcsBarQlikeTrangeSetNinety}{LightGBM on the full predictor set and XGBoost on the full predictor set}  % the 90 % MCS, best (last eliminated) first
\providecommand{\rhmcsBarQlikeTrangeStableNinety}{}\renewcommand{\rhmcsBarQlikeTrangeStableNinety}{yes}  % the 90 % MCS is the same at block lengths (5, 20)
\providecommand{\rhmcsBarQlikeTrangePMaxOutNinety}{}\renewcommand{\rhmcsBarQlikeTrangePMaxOutNinety}{\ensuremath{<}0.001}  % largest MCS p-value outside the 90 % set
\providecommand{\rhmcsBarQlikeTrangeNNinetyLo}{}\renewcommand{\rhmcsBarQlikeTrangeNNinetyLo}{2}  % ... at block length 5
\providecommand{\rhmcsBarQlikeTrangeSetNinetyLo}{}\renewcommand{\rhmcsBarQlikeTrangeSetNinetyLo}{LightGBM on the full predictor set and XGBoost on the full predictor set}  % ... at block length 5
\providecommand{\rhmcsBarQlikeTrangeNNinetyHi}{}\renewcommand{\rhmcsBarQlikeTrangeNNinetyHi}{2}  % ... at block length 20
\providecommand{\rhmcsBarQlikeTrangeSetNinetyHi}{}\renewcommand{\rhmcsBarQlikeTrangeSetNinetyHi}{LightGBM on the full predictor set and XGBoost on the full predictor set}  % ... at block length 20
\providecommand{\rhmcsBarQlikeTrangeNSevenFive}{}\renewcommand{\rhmcsBarQlikeTrangeNSevenFive}{2}  % models in the 75 % MCS (of 10)
\providecommand{\rhmcsBarQlikeTrangeSetSevenFive}{}\renewcommand{\rhmcsBarQlikeTrangeSetSevenFive}{LightGBM on the full predictor set and XGBoost on the full predictor set}  % the 75 % MCS, best (last eliminated) first
\providecommand{\rhmcsBarQlikeTrangeStableSevenFive}{}\renewcommand{\rhmcsBarQlikeTrangeStableSevenFive}{yes}  % the 75 % MCS is the same at block lengths (5, 20)
\providecommand{\rhmcsBarQlikeTrangePMaxOutSevenFive}{}\renewcommand{\rhmcsBarQlikeTrangePMaxOutSevenFive}{\ensuremath{<}0.001}  % largest MCS p-value outside the 75 % set
\providecommand{\rhmcsBarQlikeTrangeNSevenFiveLo}{}\renewcommand{\rhmcsBarQlikeTrangeNSevenFiveLo}{2}  % ... at block length 5
\providecommand{\rhmcsBarQlikeTrangeSetSevenFiveLo}{}\renewcommand{\rhmcsBarQlikeTrangeSetSevenFiveLo}{LightGBM on the full predictor set and XGBoost on the full predictor set}  % ... at block length 5
\providecommand{\rhmcsBarQlikeTrangeNSevenFiveHi}{}\renewcommand{\rhmcsBarQlikeTrangeNSevenFiveHi}{2}  % ... at block length 20
\providecommand{\rhmcsBarQlikeTrangeSetSevenFiveHi}{}\renewcommand{\rhmcsBarQlikeTrangeSetSevenFiveHi}{LightGBM on the full predictor set and XGBoost on the full predictor set}  % ... at block length 20
\providecommand{\rhmcsBarQlikeTrangeBest}{}\renewcommand{\rhmcsBarQlikeTrangeBest}{LightGBM on the full predictor set}  % the last model left (MCS p-value 1)
\providecommand{\rhmcsBarMseTmaxPLinHar}{}\renewcommand{\rhmcsBarMseTmaxPLinHar}{1.00}  % MCS p-value, every bar, 2014--2024, Mse, Tmax: Linear, HAR
\providecommand{\rhmcsBarMseTmaxPRidgeFull}{}\renewcommand{\rhmcsBarMseTmaxPRidgeFull}{0.49}  % MCS p-value, every bar, 2014--2024, Mse, Tmax: Ridge, Full
\providecommand{\rhmcsBarMseTmaxPLassoFull}{}\renewcommand{\rhmcsBarMseTmaxPLassoFull}{0.32}  % MCS p-value, every bar, 2014--2024, Mse, Tmax: Lasso, Full
\providecommand{\rhmcsBarMseTmaxPEnetFull}{}\renewcommand{\rhmcsBarMseTmaxPEnetFull}{0.43}  % MCS p-value, every bar, 2014--2024, Mse, Tmax: Elastic net, Full
\providecommand{\rhmcsBarMseTmaxPLgbmHar}{}\renewcommand{\rhmcsBarMseTmaxPLgbmHar}{0.49}  % MCS p-value, every bar, 2014--2024, Mse, Tmax: LightGBM, HAR
\providecommand{\rhmcsBarMseTmaxPLgbmFull}{}\renewcommand{\rhmcsBarMseTmaxPLgbmFull}{0.49}  % MCS p-value, every bar, 2014--2024, Mse, Tmax: LightGBM, Full
\providecommand{\rhmcsBarMseTmaxPXgbHar}{}\renewcommand{\rhmcsBarMseTmaxPXgbHar}{0.49}  % MCS p-value, every bar, 2014--2024, Mse, Tmax: XGBoost, HAR
\providecommand{\rhmcsBarMseTmaxPXgbFull}{}\renewcommand{\rhmcsBarMseTmaxPXgbFull}{0.49}  % MCS p-value, every bar, 2014--2024, Mse, Tmax: XGBoost, Full
\providecommand{\rhmcsBarMseTmaxPRfHar}{}\renewcommand{\rhmcsBarMseTmaxPRfHar}{0.49}  % MCS p-value, every bar, 2014--2024, Mse, Tmax: Random forest, HAR
\providecommand{\rhmcsBarMseTmaxPRfFull}{}\renewcommand{\rhmcsBarMseTmaxPRfFull}{0.49}  % MCS p-value, every bar, 2014--2024, Mse, Tmax: Random forest, Full
\providecommand{\rhmcsBarMseTmaxNNinety}{}\renewcommand{\rhmcsBarMseTmaxNNinety}{10}  % models in the 90 % MCS (of 10)
\providecommand{\rhmcsBarMseTmaxSetNinety}{}\renewcommand{\rhmcsBarMseTmaxSetNinety}{the linear models on the HAR benchmark, XGBoost on the HAR benchmark, LightGBM on the HAR benchmark, the random forest on the HAR benchmark, XGBoost on the full predictor set, LightGBM on the full predictor set, the random forest on the full predictor set, ridge regression on the full predictor set, the elastic net on the full predictor set, and the lasso on the full predictor set}  % the 90 % MCS, best (last eliminated) first
\providecommand{\rhmcsBarMseTmaxStableNinety}{}\renewcommand{\rhmcsBarMseTmaxStableNinety}{yes}  % the 90 % MCS is the same at block lengths (5, 20)
\providecommand{\rhmcsBarMseTmaxPMaxOutNinety}{}\renewcommand{\rhmcsBarMseTmaxPMaxOutNinety}{--}  % largest MCS p-value outside the 90 % set
\providecommand{\rhmcsBarMseTmaxNNinetyLo}{}\renewcommand{\rhmcsBarMseTmaxNNinetyLo}{10}  % ... at block length 5
\providecommand{\rhmcsBarMseTmaxSetNinetyLo}{}\renewcommand{\rhmcsBarMseTmaxSetNinetyLo}{the linear models on the HAR benchmark, XGBoost on the HAR benchmark, LightGBM on the HAR benchmark, the random forest on the HAR benchmark, XGBoost on the full predictor set, LightGBM on the full predictor set, the random forest on the full predictor set, the elastic net on the full predictor set, ridge regression on the full predictor set, and the lasso on the full predictor set}  % ... at block length 5
\providecommand{\rhmcsBarMseTmaxNNinetyHi}{}\renewcommand{\rhmcsBarMseTmaxNNinetyHi}{10}  % ... at block length 20
\providecommand{\rhmcsBarMseTmaxSetNinetyHi}{}\renewcommand{\rhmcsBarMseTmaxSetNinetyHi}{the linear models on the HAR benchmark, XGBoost on the HAR benchmark, LightGBM on the HAR benchmark, the random forest on the HAR benchmark, XGBoost on the full predictor set, LightGBM on the full predictor set, the random forest on the full predictor set, ridge regression on the full predictor set, the elastic net on the full predictor set, and the lasso on the full predictor set}  % ... at block length 20
\providecommand{\rhmcsBarMseTmaxNSevenFive}{}\renewcommand{\rhmcsBarMseTmaxNSevenFive}{10}  % models in the 75 % MCS (of 10)
\providecommand{\rhmcsBarMseTmaxSetSevenFive}{}\renewcommand{\rhmcsBarMseTmaxSetSevenFive}{the linear models on the HAR benchmark, XGBoost on the HAR benchmark, LightGBM on the HAR benchmark, the random forest on the HAR benchmark, XGBoost on the full predictor set, LightGBM on the full predictor set, the random forest on the full predictor set, ridge regression on the full predictor set, the elastic net on the full predictor set, and the lasso on the full predictor set}  % the 75 % MCS, best (last eliminated) first
\providecommand{\rhmcsBarMseTmaxStableSevenFive}{}\renewcommand{\rhmcsBarMseTmaxStableSevenFive}{no}  % the 75 % MCS is the same at block lengths (5, 20)
\providecommand{\rhmcsBarMseTmaxPMaxOutSevenFive}{}\renewcommand{\rhmcsBarMseTmaxPMaxOutSevenFive}{--}  % largest MCS p-value outside the 75 % set
\providecommand{\rhmcsBarMseTmaxNSevenFiveLo}{}\renewcommand{\rhmcsBarMseTmaxNSevenFiveLo}{9}  % ... at block length 5
\providecommand{\rhmcsBarMseTmaxSetSevenFiveLo}{}\renewcommand{\rhmcsBarMseTmaxSetSevenFiveLo}{the linear models on the HAR benchmark, XGBoost on the HAR benchmark, LightGBM on the HAR benchmark, the random forest on the HAR benchmark, XGBoost on the full predictor set, LightGBM on the full predictor set, the random forest on the full predictor set, the elastic net on the full predictor set, and ridge regression on the full predictor set}  % ... at block length 5
\providecommand{\rhmcsBarMseTmaxNSevenFiveHi}{}\renewcommand{\rhmcsBarMseTmaxNSevenFiveHi}{10}  % ... at block length 20
\providecommand{\rhmcsBarMseTmaxSetSevenFiveHi}{}\renewcommand{\rhmcsBarMseTmaxSetSevenFiveHi}{the linear models on the HAR benchmark, XGBoost on the HAR benchmark, LightGBM on the HAR benchmark, the random forest on the HAR benchmark, XGBoost on the full predictor set, LightGBM on the full predictor set, the random forest on the full predictor set, ridge regression on the full predictor set, the elastic net on the full predictor set, and the lasso on the full predictor set}  % ... at block length 20
\providecommand{\rhmcsBarMseTmaxBest}{}\renewcommand{\rhmcsBarMseTmaxBest}{the linear models on the HAR benchmark}  % the last model left (MCS p-value 1)
\providecommand{\rhmcsBarMseTrangePLinHar}{}\renewcommand{\rhmcsBarMseTrangePLinHar}{1.00}  % MCS p-value, every bar, 2014--2024, Mse, Trange: Linear, HAR
\providecommand{\rhmcsBarMseTrangePRidgeFull}{}\renewcommand{\rhmcsBarMseTrangePRidgeFull}{0.42}  % MCS p-value, every bar, 2014--2024, Mse, Trange: Ridge, Full
\providecommand{\rhmcsBarMseTrangePLassoFull}{}\renewcommand{\rhmcsBarMseTrangePLassoFull}{0.42}  % MCS p-value, every bar, 2014--2024, Mse, Trange: Lasso, Full
\providecommand{\rhmcsBarMseTrangePEnetFull}{}\renewcommand{\rhmcsBarMseTrangePEnetFull}{0.42}  % MCS p-value, every bar, 2014--2024, Mse, Trange: Elastic net, Full
\providecommand{\rhmcsBarMseTrangePLgbmHar}{}\renewcommand{\rhmcsBarMseTrangePLgbmHar}{0.31}  % MCS p-value, every bar, 2014--2024, Mse, Trange: LightGBM, HAR
\providecommand{\rhmcsBarMseTrangePLgbmFull}{}\renewcommand{\rhmcsBarMseTrangePLgbmFull}{0.42}  % MCS p-value, every bar, 2014--2024, Mse, Trange: LightGBM, Full
\providecommand{\rhmcsBarMseTrangePXgbHar}{}\renewcommand{\rhmcsBarMseTrangePXgbHar}{0.42}  % MCS p-value, every bar, 2014--2024, Mse, Trange: XGBoost, HAR
\providecommand{\rhmcsBarMseTrangePXgbFull}{}\renewcommand{\rhmcsBarMseTrangePXgbFull}{0.42}  % MCS p-value, every bar, 2014--2024, Mse, Trange: XGBoost, Full
\providecommand{\rhmcsBarMseTrangePRfHar}{}\renewcommand{\rhmcsBarMseTrangePRfHar}{0.31}  % MCS p-value, every bar, 2014--2024, Mse, Trange: Random forest, HAR
\providecommand{\rhmcsBarMseTrangePRfFull}{}\renewcommand{\rhmcsBarMseTrangePRfFull}{0.42}  % MCS p-value, every bar, 2014--2024, Mse, Trange: Random forest, Full
\providecommand{\rhmcsBarMseTrangeNNinety}{}\renewcommand{\rhmcsBarMseTrangeNNinety}{10}  % models in the 90 % MCS (of 10)
\providecommand{\rhmcsBarMseTrangeSetNinety}{}\renewcommand{\rhmcsBarMseTrangeSetNinety}{the linear models on the HAR benchmark, the elastic net on the full predictor set, ridge regression on the full predictor set, XGBoost on the full predictor set, LightGBM on the full predictor set, the random forest on the full predictor set, the lasso on the full predictor set, XGBoost on the HAR benchmark, LightGBM on the HAR benchmark, and the random forest on the HAR benchmark}  % the 90 % MCS, best (last eliminated) first
\providecommand{\rhmcsBarMseTrangeStableNinety}{}\renewcommand{\rhmcsBarMseTrangeStableNinety}{yes}  % the 90 % MCS is the same at block lengths (5, 20)
\providecommand{\rhmcsBarMseTrangePMaxOutNinety}{}\renewcommand{\rhmcsBarMseTrangePMaxOutNinety}{--}  % largest MCS p-value outside the 90 % set
\providecommand{\rhmcsBarMseTrangeNNinetyLo}{}\renewcommand{\rhmcsBarMseTrangeNNinetyLo}{10}  % ... at block length 5
\providecommand{\rhmcsBarMseTrangeSetNinetyLo}{}\renewcommand{\rhmcsBarMseTrangeSetNinetyLo}{the linear models on the HAR benchmark, the elastic net on the full predictor set, ridge regression on the full predictor set, XGBoost on the HAR benchmark, XGBoost on the full predictor set, LightGBM on the full predictor set, the random forest on the full predictor set, LightGBM on the HAR benchmark, the lasso on the full predictor set, and the random forest on the HAR benchmark}  % ... at block length 5
\providecommand{\rhmcsBarMseTrangeNNinetyHi}{}\renewcommand{\rhmcsBarMseTrangeNNinetyHi}{10}  % ... at block length 20
\providecommand{\rhmcsBarMseTrangeSetNinetyHi}{}\renewcommand{\rhmcsBarMseTrangeSetNinetyHi}{the linear models on the HAR benchmark, the elastic net on the full predictor set, ridge regression on the full predictor set, XGBoost on the full predictor set, the lasso on the full predictor set, LightGBM on the full predictor set, the random forest on the full predictor set, XGBoost on the HAR benchmark, LightGBM on the HAR benchmark, and the random forest on the HAR benchmark}  % ... at block length 20
\providecommand{\rhmcsBarMseTrangeNSevenFive}{}\renewcommand{\rhmcsBarMseTrangeNSevenFive}{10}  % models in the 75 % MCS (of 10)
\providecommand{\rhmcsBarMseTrangeSetSevenFive}{}\renewcommand{\rhmcsBarMseTrangeSetSevenFive}{the linear models on the HAR benchmark, the elastic net on the full predictor set, ridge regression on the full predictor set, XGBoost on the full predictor set, LightGBM on the full predictor set, the random forest on the full predictor set, the lasso on the full predictor set, XGBoost on the HAR benchmark, LightGBM on the HAR benchmark, and the random forest on the HAR benchmark}  % the 75 % MCS, best (last eliminated) first
\providecommand{\rhmcsBarMseTrangeStableSevenFive}{}\renewcommand{\rhmcsBarMseTrangeStableSevenFive}{yes}  % the 75 % MCS is the same at block lengths (5, 20)
\providecommand{\rhmcsBarMseTrangePMaxOutSevenFive}{}\renewcommand{\rhmcsBarMseTrangePMaxOutSevenFive}{--}  % largest MCS p-value outside the 75 % set
\providecommand{\rhmcsBarMseTrangeNSevenFiveLo}{}\renewcommand{\rhmcsBarMseTrangeNSevenFiveLo}{10}  % ... at block length 5
\providecommand{\rhmcsBarMseTrangeSetSevenFiveLo}{}\renewcommand{\rhmcsBarMseTrangeSetSevenFiveLo}{the linear models on the HAR benchmark, the elastic net on the full predictor set, ridge regression on the full predictor set, XGBoost on the HAR benchmark, XGBoost on the full predictor set, LightGBM on the full predictor set, the random forest on the full predictor set, LightGBM on the HAR benchmark, the lasso on the full predictor set, and the random forest on the HAR benchmark}  % ... at block length 5
\providecommand{\rhmcsBarMseTrangeNSevenFiveHi}{}\renewcommand{\rhmcsBarMseTrangeNSevenFiveHi}{10}  % ... at block length 20
\providecommand{\rhmcsBarMseTrangeSetSevenFiveHi}{}\renewcommand{\rhmcsBarMseTrangeSetSevenFiveHi}{the linear models on the HAR benchmark, the elastic net on the full predictor set, ridge regression on the full predictor set, XGBoost on the full predictor set, the lasso on the full predictor set, LightGBM on the full predictor set, the random forest on the full predictor set, XGBoost on the HAR benchmark, LightGBM on the HAR benchmark, and the random forest on the HAR benchmark}  % ... at block length 20
\providecommand{\rhmcsBarMseTrangeBest}{}\renewcommand{\rhmcsBarMseTrangeBest}{the linear models on the HAR benchmark}  % the last model left (MCS p-value 1)
\providecommand{\rhmcsAllQlikeTmaxPLinHar}{}\renewcommand{\rhmcsAllQlikeTmaxPLinHar}{0.19}  % MCS p-value, 16:00 bar, 2014--2024, Qlike, Tmax: Linear, HAR
\providecommand{\rhmcsAllQlikeTmaxPRidgeFull}{}\renewcommand{\rhmcsAllQlikeTmaxPRidgeFull}{0.19}  % MCS p-value, 16:00 bar, 2014--2024, Qlike, Tmax: Ridge, Full
\providecommand{\rhmcsAllQlikeTmaxPLassoFull}{}\renewcommand{\rhmcsAllQlikeTmaxPLassoFull}{0.071}  % MCS p-value, 16:00 bar, 2014--2024, Qlike, Tmax: Lasso, Full
\providecommand{\rhmcsAllQlikeTmaxPEnetFull}{}\renewcommand{\rhmcsAllQlikeTmaxPEnetFull}{0.19}  % MCS p-value, 16:00 bar, 2014--2024, Qlike, Tmax: Elastic net, Full
\providecommand{\rhmcsAllQlikeTmaxPLgbmHar}{}\renewcommand{\rhmcsAllQlikeTmaxPLgbmHar}{0.19}  % MCS p-value, 16:00 bar, 2014--2024, Qlike, Tmax: LightGBM, HAR
\providecommand{\rhmcsAllQlikeTmaxPLgbmFull}{}\renewcommand{\rhmcsAllQlikeTmaxPLgbmFull}{0.29}  % MCS p-value, 16:00 bar, 2014--2024, Qlike, Tmax: LightGBM, Full
\providecommand{\rhmcsAllQlikeTmaxPXgbHar}{}\renewcommand{\rhmcsAllQlikeTmaxPXgbHar}{0.19}  % MCS p-value, 16:00 bar, 2014--2024, Qlike, Tmax: XGBoost, HAR
\providecommand{\rhmcsAllQlikeTmaxPXgbFull}{}\renewcommand{\rhmcsAllQlikeTmaxPXgbFull}{1.00}  % MCS p-value, 16:00 bar, 2014--2024, Qlike, Tmax: XGBoost, Full
\providecommand{\rhmcsAllQlikeTmaxPRfHar}{}\renewcommand{\rhmcsAllQlikeTmaxPRfHar}{0.19}  % MCS p-value, 16:00 bar, 2014--2024, Qlike, Tmax: Random forest, HAR
\providecommand{\rhmcsAllQlikeTmaxPRfFull}{}\renewcommand{\rhmcsAllQlikeTmaxPRfFull}{0.19}  % MCS p-value, 16:00 bar, 2014--2024, Qlike, Tmax: Random forest, Full
\providecommand{\rhmcsAllQlikeTmaxNNinety}{}\renewcommand{\rhmcsAllQlikeTmaxNNinety}{9}  % models in the 90 % MCS (of 10)
\providecommand{\rhmcsAllQlikeTmaxSetNinety}{}\renewcommand{\rhmcsAllQlikeTmaxSetNinety}{XGBoost on the full predictor set, LightGBM on the full predictor set, XGBoost on the HAR benchmark, LightGBM on the HAR benchmark, the random forest on the full predictor set, the elastic net on the full predictor set, the random forest on the HAR benchmark, ridge regression on the full predictor set, and the linear models on the HAR benchmark}  % the 90 % MCS, best (last eliminated) first
\providecommand{\rhmcsAllQlikeTmaxStableNinety}{}\renewcommand{\rhmcsAllQlikeTmaxStableNinety}{yes}  % the 90 % MCS is the same at block lengths (5, 20)
\providecommand{\rhmcsAllQlikeTmaxPMaxOutNinety}{}\renewcommand{\rhmcsAllQlikeTmaxPMaxOutNinety}{0.071}  % largest MCS p-value outside the 90 % set
\providecommand{\rhmcsAllQlikeTmaxNNinetyLo}{}\renewcommand{\rhmcsAllQlikeTmaxNNinetyLo}{9}  % ... at block length 5
\providecommand{\rhmcsAllQlikeTmaxSetNinetyLo}{}\renewcommand{\rhmcsAllQlikeTmaxSetNinetyLo}{XGBoost on the full predictor set, LightGBM on the full predictor set, XGBoost on the HAR benchmark, LightGBM on the HAR benchmark, the random forest on the full predictor set, the elastic net on the full predictor set, the random forest on the HAR benchmark, ridge regression on the full predictor set, and the linear models on the HAR benchmark}  % ... at block length 5
\providecommand{\rhmcsAllQlikeTmaxNNinetyHi}{}\renewcommand{\rhmcsAllQlikeTmaxNNinetyHi}{9}  % ... at block length 20
\providecommand{\rhmcsAllQlikeTmaxSetNinetyHi}{}\renewcommand{\rhmcsAllQlikeTmaxSetNinetyHi}{XGBoost on the full predictor set, LightGBM on the full predictor set, XGBoost on the HAR benchmark, LightGBM on the HAR benchmark, the random forest on the full predictor set, the elastic net on the full predictor set, the random forest on the HAR benchmark, ridge regression on the full predictor set, and the linear models on the HAR benchmark}  % ... at block length 20
\providecommand{\rhmcsAllQlikeTmaxNSevenFive}{}\renewcommand{\rhmcsAllQlikeTmaxNSevenFive}{2}  % models in the 75 % MCS (of 10)
\providecommand{\rhmcsAllQlikeTmaxSetSevenFive}{}\renewcommand{\rhmcsAllQlikeTmaxSetSevenFive}{XGBoost on the full predictor set and LightGBM on the full predictor set}  % the 75 % MCS, best (last eliminated) first
\providecommand{\rhmcsAllQlikeTmaxStableSevenFive}{}\renewcommand{\rhmcsAllQlikeTmaxStableSevenFive}{yes}  % the 75 % MCS is the same at block lengths (5, 20)
\providecommand{\rhmcsAllQlikeTmaxPMaxOutSevenFive}{}\renewcommand{\rhmcsAllQlikeTmaxPMaxOutSevenFive}{0.19}  % largest MCS p-value outside the 75 % set
\providecommand{\rhmcsAllQlikeTmaxNSevenFiveLo}{}\renewcommand{\rhmcsAllQlikeTmaxNSevenFiveLo}{2}  % ... at block length 5
\providecommand{\rhmcsAllQlikeTmaxSetSevenFiveLo}{}\renewcommand{\rhmcsAllQlikeTmaxSetSevenFiveLo}{XGBoost on the full predictor set and LightGBM on the full predictor set}  % ... at block length 5
\providecommand{\rhmcsAllQlikeTmaxNSevenFiveHi}{}\renewcommand{\rhmcsAllQlikeTmaxNSevenFiveHi}{2}  % ... at block length 20
\providecommand{\rhmcsAllQlikeTmaxSetSevenFiveHi}{}\renewcommand{\rhmcsAllQlikeTmaxSetSevenFiveHi}{XGBoost on the full predictor set and LightGBM on the full predictor set}  % ... at block length 20
\providecommand{\rhmcsAllQlikeTmaxBest}{}\renewcommand{\rhmcsAllQlikeTmaxBest}{XGBoost on the full predictor set}  % the last model left (MCS p-value 1)
\providecommand{\rhmcsAllQlikeTrangePLinHar}{}\renewcommand{\rhmcsAllQlikeTrangePLinHar}{\ensuremath{<}0.001}  % MCS p-value, 16:00 bar, 2014--2024, Qlike, Trange: Linear, HAR
\providecommand{\rhmcsAllQlikeTrangePRidgeFull}{}\renewcommand{\rhmcsAllQlikeTrangePRidgeFull}{0.008}  % MCS p-value, 16:00 bar, 2014--2024, Qlike, Trange: Ridge, Full
\providecommand{\rhmcsAllQlikeTrangePLassoFull}{}\renewcommand{\rhmcsAllQlikeTrangePLassoFull}{0.006}  % MCS p-value, 16:00 bar, 2014--2024, Qlike, Trange: Lasso, Full
\providecommand{\rhmcsAllQlikeTrangePEnetFull}{}\renewcommand{\rhmcsAllQlikeTrangePEnetFull}{0.010}  % MCS p-value, 16:00 bar, 2014--2024, Qlike, Trange: Elastic net, Full
\providecommand{\rhmcsAllQlikeTrangePLgbmHar}{}\renewcommand{\rhmcsAllQlikeTrangePLgbmHar}{0.14}  % MCS p-value, 16:00 bar, 2014--2024, Qlike, Trange: LightGBM, HAR
\providecommand{\rhmcsAllQlikeTrangePLgbmFull}{}\renewcommand{\rhmcsAllQlikeTrangePLgbmFull}{0.29}  % MCS p-value, 16:00 bar, 2014--2024, Qlike, Trange: LightGBM, Full
\providecommand{\rhmcsAllQlikeTrangePXgbHar}{}\renewcommand{\rhmcsAllQlikeTrangePXgbHar}{0.14}  % MCS p-value, 16:00 bar, 2014--2024, Qlike, Trange: XGBoost, HAR
\providecommand{\rhmcsAllQlikeTrangePXgbFull}{}\renewcommand{\rhmcsAllQlikeTrangePXgbFull}{1.00}  % MCS p-value, 16:00 bar, 2014--2024, Qlike, Trange: XGBoost, Full
\providecommand{\rhmcsAllQlikeTrangePRfHar}{}\renewcommand{\rhmcsAllQlikeTrangePRfHar}{0.001}  % MCS p-value, 16:00 bar, 2014--2024, Qlike, Trange: Random forest, HAR
\providecommand{\rhmcsAllQlikeTrangePRfFull}{}\renewcommand{\rhmcsAllQlikeTrangePRfFull}{\ensuremath{<}0.001}  % MCS p-value, 16:00 bar, 2014--2024, Qlike, Trange: Random forest, Full
\providecommand{\rhmcsAllQlikeTrangeNNinety}{}\renewcommand{\rhmcsAllQlikeTrangeNNinety}{4}  % models in the 90 % MCS (of 10)
\providecommand{\rhmcsAllQlikeTrangeSetNinety}{}\renewcommand{\rhmcsAllQlikeTrangeSetNinety}{XGBoost on the full predictor set, LightGBM on the full predictor set, XGBoost on the HAR benchmark, and LightGBM on the HAR benchmark}  % the 90 % MCS, best (last eliminated) first
\providecommand{\rhmcsAllQlikeTrangeStableNinety}{}\renewcommand{\rhmcsAllQlikeTrangeStableNinety}{yes}  % the 90 % MCS is the same at block lengths (5, 20)
\providecommand{\rhmcsAllQlikeTrangePMaxOutNinety}{}\renewcommand{\rhmcsAllQlikeTrangePMaxOutNinety}{0.010}  % largest MCS p-value outside the 90 % set
\providecommand{\rhmcsAllQlikeTrangeNNinetyLo}{}\renewcommand{\rhmcsAllQlikeTrangeNNinetyLo}{4}  % ... at block length 5
\providecommand{\rhmcsAllQlikeTrangeSetNinetyLo}{}\renewcommand{\rhmcsAllQlikeTrangeSetNinetyLo}{XGBoost on the full predictor set, LightGBM on the full predictor set, LightGBM on the HAR benchmark, and XGBoost on the HAR benchmark}  % ... at block length 5
\providecommand{\rhmcsAllQlikeTrangeNNinetyHi}{}\renewcommand{\rhmcsAllQlikeTrangeNNinetyHi}{4}  % ... at block length 20
\providecommand{\rhmcsAllQlikeTrangeSetNinetyHi}{}\renewcommand{\rhmcsAllQlikeTrangeSetNinetyHi}{XGBoost on the full predictor set, LightGBM on the full predictor set, LightGBM on the HAR benchmark, and XGBoost on the HAR benchmark}  % ... at block length 20
\providecommand{\rhmcsAllQlikeTrangeNSevenFive}{}\renewcommand{\rhmcsAllQlikeTrangeNSevenFive}{2}  % models in the 75 % MCS (of 10)
\providecommand{\rhmcsAllQlikeTrangeSetSevenFive}{}\renewcommand{\rhmcsAllQlikeTrangeSetSevenFive}{XGBoost on the full predictor set and LightGBM on the full predictor set}  % the 75 % MCS, best (last eliminated) first
\providecommand{\rhmcsAllQlikeTrangeStableSevenFive}{}\renewcommand{\rhmcsAllQlikeTrangeStableSevenFive}{yes}  % the 75 % MCS is the same at block lengths (5, 20)
\providecommand{\rhmcsAllQlikeTrangePMaxOutSevenFive}{}\renewcommand{\rhmcsAllQlikeTrangePMaxOutSevenFive}{0.14}  % largest MCS p-value outside the 75 % set
\providecommand{\rhmcsAllQlikeTrangeNSevenFiveLo}{}\renewcommand{\rhmcsAllQlikeTrangeNSevenFiveLo}{2}  % ... at block length 5
\providecommand{\rhmcsAllQlikeTrangeSetSevenFiveLo}{}\renewcommand{\rhmcsAllQlikeTrangeSetSevenFiveLo}{XGBoost on the full predictor set and LightGBM on the full predictor set}  % ... at block length 5
\providecommand{\rhmcsAllQlikeTrangeNSevenFiveHi}{}\renewcommand{\rhmcsAllQlikeTrangeNSevenFiveHi}{2}  % ... at block length 20
\providecommand{\rhmcsAllQlikeTrangeSetSevenFiveHi}{}\renewcommand{\rhmcsAllQlikeTrangeSetSevenFiveHi}{XGBoost on the full predictor set and LightGBM on the full predictor set}  % ... at block length 20
\providecommand{\rhmcsAllQlikeTrangeBest}{}\renewcommand{\rhmcsAllQlikeTrangeBest}{XGBoost on the full predictor set}  % the last model left (MCS p-value 1)
\providecommand{\rhmcsAllMseTmaxPLinHar}{}\renewcommand{\rhmcsAllMseTmaxPLinHar}{0.59}  % MCS p-value, 16:00 bar, 2014--2024, Mse, Tmax: Linear, HAR
\providecommand{\rhmcsAllMseTmaxPRidgeFull}{}\renewcommand{\rhmcsAllMseTmaxPRidgeFull}{0.43}  % MCS p-value, 16:00 bar, 2014--2024, Mse, Tmax: Ridge, Full
\providecommand{\rhmcsAllMseTmaxPLassoFull}{}\renewcommand{\rhmcsAllMseTmaxPLassoFull}{0.38}  % MCS p-value, 16:00 bar, 2014--2024, Mse, Tmax: Lasso, Full
\providecommand{\rhmcsAllMseTmaxPEnetFull}{}\renewcommand{\rhmcsAllMseTmaxPEnetFull}{0.48}  % MCS p-value, 16:00 bar, 2014--2024, Mse, Tmax: Elastic net, Full
\providecommand{\rhmcsAllMseTmaxPLgbmHar}{}\renewcommand{\rhmcsAllMseTmaxPLgbmHar}{0.59}  % MCS p-value, 16:00 bar, 2014--2024, Mse, Tmax: LightGBM, HAR
\providecommand{\rhmcsAllMseTmaxPLgbmFull}{}\renewcommand{\rhmcsAllMseTmaxPLgbmFull}{0.48}  % MCS p-value, 16:00 bar, 2014--2024, Mse, Tmax: LightGBM, Full
\providecommand{\rhmcsAllMseTmaxPXgbHar}{}\renewcommand{\rhmcsAllMseTmaxPXgbHar}{0.59}  % MCS p-value, 16:00 bar, 2014--2024, Mse, Tmax: XGBoost, HAR
\providecommand{\rhmcsAllMseTmaxPXgbFull}{}\renewcommand{\rhmcsAllMseTmaxPXgbFull}{0.48}  % MCS p-value, 16:00 bar, 2014--2024, Mse, Tmax: XGBoost, Full
\providecommand{\rhmcsAllMseTmaxPRfHar}{}\renewcommand{\rhmcsAllMseTmaxPRfHar}{1.00}  % MCS p-value, 16:00 bar, 2014--2024, Mse, Tmax: Random forest, HAR
\providecommand{\rhmcsAllMseTmaxPRfFull}{}\renewcommand{\rhmcsAllMseTmaxPRfFull}{0.59}  % MCS p-value, 16:00 bar, 2014--2024, Mse, Tmax: Random forest, Full
\providecommand{\rhmcsAllMseTmaxNNinety}{}\renewcommand{\rhmcsAllMseTmaxNNinety}{10}  % models in the 90 % MCS (of 10)
\providecommand{\rhmcsAllMseTmaxSetNinety}{}\renewcommand{\rhmcsAllMseTmaxSetNinety}{the random forest on the HAR benchmark, XGBoost on the HAR benchmark, LightGBM on the HAR benchmark, the linear models on the HAR benchmark, the random forest on the full predictor set, XGBoost on the full predictor set, LightGBM on the full predictor set, the elastic net on the full predictor set, ridge regression on the full predictor set, and the lasso on the full predictor set}  % the 90 % MCS, best (last eliminated) first
\providecommand{\rhmcsAllMseTmaxStableNinety}{}\renewcommand{\rhmcsAllMseTmaxStableNinety}{yes}  % the 90 % MCS is the same at block lengths (5, 20)
\providecommand{\rhmcsAllMseTmaxPMaxOutNinety}{}\renewcommand{\rhmcsAllMseTmaxPMaxOutNinety}{--}  % largest MCS p-value outside the 90 % set
\providecommand{\rhmcsAllMseTmaxNNinetyLo}{}\renewcommand{\rhmcsAllMseTmaxNNinetyLo}{10}  % ... at block length 5
\providecommand{\rhmcsAllMseTmaxSetNinetyLo}{}\renewcommand{\rhmcsAllMseTmaxSetNinetyLo}{the random forest on the HAR benchmark, XGBoost on the HAR benchmark, LightGBM on the HAR benchmark, the linear models on the HAR benchmark, the random forest on the full predictor set, XGBoost on the full predictor set, LightGBM on the full predictor set, the elastic net on the full predictor set, ridge regression on the full predictor set, and the lasso on the full predictor set}  % ... at block length 5
\providecommand{\rhmcsAllMseTmaxNNinetyHi}{}\renewcommand{\rhmcsAllMseTmaxNNinetyHi}{10}  % ... at block length 20
\providecommand{\rhmcsAllMseTmaxSetNinetyHi}{}\renewcommand{\rhmcsAllMseTmaxSetNinetyHi}{the random forest on the HAR benchmark, XGBoost on the HAR benchmark, LightGBM on the HAR benchmark, the linear models on the HAR benchmark, the random forest on the full predictor set, XGBoost on the full predictor set, LightGBM on the full predictor set, the elastic net on the full predictor set, ridge regression on the full predictor set, and the lasso on the full predictor set}  % ... at block length 20
\providecommand{\rhmcsAllMseTmaxNSevenFive}{}\renewcommand{\rhmcsAllMseTmaxNSevenFive}{10}  % models in the 75 % MCS (of 10)
\providecommand{\rhmcsAllMseTmaxSetSevenFive}{}\renewcommand{\rhmcsAllMseTmaxSetSevenFive}{the random forest on the HAR benchmark, XGBoost on the HAR benchmark, LightGBM on the HAR benchmark, the linear models on the HAR benchmark, the random forest on the full predictor set, XGBoost on the full predictor set, LightGBM on the full predictor set, the elastic net on the full predictor set, ridge regression on the full predictor set, and the lasso on the full predictor set}  % the 75 % MCS, best (last eliminated) first
\providecommand{\rhmcsAllMseTmaxStableSevenFive}{}\renewcommand{\rhmcsAllMseTmaxStableSevenFive}{no}  % the 75 % MCS is the same at block lengths (5, 20)
\providecommand{\rhmcsAllMseTmaxPMaxOutSevenFive}{}\renewcommand{\rhmcsAllMseTmaxPMaxOutSevenFive}{--}  % largest MCS p-value outside the 75 % set
\providecommand{\rhmcsAllMseTmaxNSevenFiveLo}{}\renewcommand{\rhmcsAllMseTmaxNSevenFiveLo}{9}  % ... at block length 5
\providecommand{\rhmcsAllMseTmaxSetSevenFiveLo}{}\renewcommand{\rhmcsAllMseTmaxSetSevenFiveLo}{the random forest on the HAR benchmark, XGBoost on the HAR benchmark, LightGBM on the HAR benchmark, the linear models on the HAR benchmark, the random forest on the full predictor set, XGBoost on the full predictor set, LightGBM on the full predictor set, the elastic net on the full predictor set, and ridge regression on the full predictor set}  % ... at block length 5
\providecommand{\rhmcsAllMseTmaxNSevenFiveHi}{}\renewcommand{\rhmcsAllMseTmaxNSevenFiveHi}{10}  % ... at block length 20
\providecommand{\rhmcsAllMseTmaxSetSevenFiveHi}{}\renewcommand{\rhmcsAllMseTmaxSetSevenFiveHi}{the random forest on the HAR benchmark, XGBoost on the HAR benchmark, LightGBM on the HAR benchmark, the linear models on the HAR benchmark, the random forest on the full predictor set, XGBoost on the full predictor set, LightGBM on the full predictor set, the elastic net on the full predictor set, ridge regression on the full predictor set, and the lasso on the full predictor set}  % ... at block length 20
\providecommand{\rhmcsAllMseTmaxBest}{}\renewcommand{\rhmcsAllMseTmaxBest}{the random forest on the HAR benchmark}  % the last model left (MCS p-value 1)
\providecommand{\rhmcsAllMseTrangePLinHar}{}\renewcommand{\rhmcsAllMseTrangePLinHar}{0.60}  % MCS p-value, 16:00 bar, 2014--2024, Mse, Trange: Linear, HAR
\providecommand{\rhmcsAllMseTrangePRidgeFull}{}\renewcommand{\rhmcsAllMseTrangePRidgeFull}{0.60}  % MCS p-value, 16:00 bar, 2014--2024, Mse, Trange: Ridge, Full
\providecommand{\rhmcsAllMseTrangePLassoFull}{}\renewcommand{\rhmcsAllMseTrangePLassoFull}{0.60}  % MCS p-value, 16:00 bar, 2014--2024, Mse, Trange: Lasso, Full
\providecommand{\rhmcsAllMseTrangePEnetFull}{}\renewcommand{\rhmcsAllMseTrangePEnetFull}{0.60}  % MCS p-value, 16:00 bar, 2014--2024, Mse, Trange: Elastic net, Full
\providecommand{\rhmcsAllMseTrangePLgbmHar}{}\renewcommand{\rhmcsAllMseTrangePLgbmHar}{0.35}  % MCS p-value, 16:00 bar, 2014--2024, Mse, Trange: LightGBM, HAR
\providecommand{\rhmcsAllMseTrangePLgbmFull}{}\renewcommand{\rhmcsAllMseTrangePLgbmFull}{0.60}  % MCS p-value, 16:00 bar, 2014--2024, Mse, Trange: LightGBM, Full
\providecommand{\rhmcsAllMseTrangePXgbHar}{}\renewcommand{\rhmcsAllMseTrangePXgbHar}{0.60}  % MCS p-value, 16:00 bar, 2014--2024, Mse, Trange: XGBoost, HAR
\providecommand{\rhmcsAllMseTrangePXgbFull}{}\renewcommand{\rhmcsAllMseTrangePXgbFull}{0.60}  % MCS p-value, 16:00 bar, 2014--2024, Mse, Trange: XGBoost, Full
\providecommand{\rhmcsAllMseTrangePRfHar}{}\renewcommand{\rhmcsAllMseTrangePRfHar}{1.00}  % MCS p-value, 16:00 bar, 2014--2024, Mse, Trange: Random forest, HAR
\providecommand{\rhmcsAllMseTrangePRfFull}{}\renewcommand{\rhmcsAllMseTrangePRfFull}{0.60}  % MCS p-value, 16:00 bar, 2014--2024, Mse, Trange: Random forest, Full
\providecommand{\rhmcsAllMseTrangeNNinety}{}\renewcommand{\rhmcsAllMseTrangeNNinety}{10}  % models in the 90 % MCS (of 10)
\providecommand{\rhmcsAllMseTrangeSetNinety}{}\renewcommand{\rhmcsAllMseTrangeSetNinety}{the random forest on the HAR benchmark, the linear models on the HAR benchmark, XGBoost on the HAR benchmark, the random forest on the full predictor set, XGBoost on the full predictor set, the elastic net on the full predictor set, ridge regression on the full predictor set, the lasso on the full predictor set, LightGBM on the full predictor set, and LightGBM on the HAR benchmark}  % the 90 % MCS, best (last eliminated) first
\providecommand{\rhmcsAllMseTrangeStableNinety}{}\renewcommand{\rhmcsAllMseTrangeStableNinety}{yes}  % the 90 % MCS is the same at block lengths (5, 20)
\providecommand{\rhmcsAllMseTrangePMaxOutNinety}{}\renewcommand{\rhmcsAllMseTrangePMaxOutNinety}{--}  % largest MCS p-value outside the 90 % set
\providecommand{\rhmcsAllMseTrangeNNinetyLo}{}\renewcommand{\rhmcsAllMseTrangeNNinetyLo}{10}  % ... at block length 5
\providecommand{\rhmcsAllMseTrangeSetNinetyLo}{}\renewcommand{\rhmcsAllMseTrangeSetNinetyLo}{the random forest on the HAR benchmark, the linear models on the HAR benchmark, XGBoost on the HAR benchmark, XGBoost on the full predictor set, the random forest on the full predictor set, the elastic net on the full predictor set, ridge regression on the full predictor set, LightGBM on the full predictor set, the lasso on the full predictor set, and LightGBM on the HAR benchmark}  % ... at block length 5
\providecommand{\rhmcsAllMseTrangeNNinetyHi}{}\renewcommand{\rhmcsAllMseTrangeNNinetyHi}{10}  % ... at block length 20
\providecommand{\rhmcsAllMseTrangeSetNinetyHi}{}\renewcommand{\rhmcsAllMseTrangeSetNinetyHi}{the random forest on the HAR benchmark, the linear models on the HAR benchmark, the random forest on the full predictor set, XGBoost on the full predictor set, the elastic net on the full predictor set, ridge regression on the full predictor set, the lasso on the full predictor set, XGBoost on the HAR benchmark, LightGBM on the full predictor set, and LightGBM on the HAR benchmark}  % ... at block length 20
\providecommand{\rhmcsAllMseTrangeNSevenFive}{}\renewcommand{\rhmcsAllMseTrangeNSevenFive}{10}  % models in the 75 % MCS (of 10)
\providecommand{\rhmcsAllMseTrangeSetSevenFive}{}\renewcommand{\rhmcsAllMseTrangeSetSevenFive}{the random forest on the HAR benchmark, the linear models on the HAR benchmark, XGBoost on the HAR benchmark, the random forest on the full predictor set, XGBoost on the full predictor set, the elastic net on the full predictor set, ridge regression on the full predictor set, the lasso on the full predictor set, LightGBM on the full predictor set, and LightGBM on the HAR benchmark}  % the 75 % MCS, best (last eliminated) first
\providecommand{\rhmcsAllMseTrangeStableSevenFive}{}\renewcommand{\rhmcsAllMseTrangeStableSevenFive}{yes}  % the 75 % MCS is the same at block lengths (5, 20)
\providecommand{\rhmcsAllMseTrangePMaxOutSevenFive}{}\renewcommand{\rhmcsAllMseTrangePMaxOutSevenFive}{--}  % largest MCS p-value outside the 75 % set
\providecommand{\rhmcsAllMseTrangeNSevenFiveLo}{}\renewcommand{\rhmcsAllMseTrangeNSevenFiveLo}{10}  % ... at block length 5
\providecommand{\rhmcsAllMseTrangeSetSevenFiveLo}{}\renewcommand{\rhmcsAllMseTrangeSetSevenFiveLo}{the random forest on the HAR benchmark, the linear models on the HAR benchmark, XGBoost on the HAR benchmark, XGBoost on the full predictor set, the random forest on the full predictor set, the elastic net on the full predictor set, ridge regression on the full predictor set, LightGBM on the full predictor set, the lasso on the full predictor set, and LightGBM on the HAR benchmark}  % ... at block length 5
\providecommand{\rhmcsAllMseTrangeNSevenFiveHi}{}\renewcommand{\rhmcsAllMseTrangeNSevenFiveHi}{10}  % ... at block length 20
\providecommand{\rhmcsAllMseTrangeSetSevenFiveHi}{}\renewcommand{\rhmcsAllMseTrangeSetSevenFiveHi}{the random forest on the HAR benchmark, the linear models on the HAR benchmark, the random forest on the full predictor set, XGBoost on the full predictor set, the elastic net on the full predictor set, ridge regression on the full predictor set, the lasso on the full predictor set, XGBoost on the HAR benchmark, LightGBM on the full predictor set, and LightGBM on the HAR benchmark}  % ... at block length 20
\providecommand{\rhmcsAllMseTrangeBest}{}\renewcommand{\rhmcsAllMseTrangeBest}{the random forest on the HAR benchmark}  % the last model left (MCS p-value 1)
\providecommand{\rhmcsOosQlikeTmaxPLinHar}{}\renewcommand{\rhmcsOosQlikeTmaxPLinHar}{0.043}  % MCS p-value, 16:00 bar, T_oos, Qlike, Tmax: Linear, HAR
\providecommand{\rhmcsOosQlikeTmaxPRidgeFull}{}\renewcommand{\rhmcsOosQlikeTmaxPRidgeFull}{0.008}  % MCS p-value, 16:00 bar, T_oos, Qlike, Tmax: Ridge, Full
\providecommand{\rhmcsOosQlikeTmaxPLassoFull}{}\renewcommand{\rhmcsOosQlikeTmaxPLassoFull}{0.008}  % MCS p-value, 16:00 bar, T_oos, Qlike, Tmax: Lasso, Full
\providecommand{\rhmcsOosQlikeTmaxPEnetFull}{}\renewcommand{\rhmcsOosQlikeTmaxPEnetFull}{0.024}  % MCS p-value, 16:00 bar, T_oos, Qlike, Tmax: Elastic net, Full
\providecommand{\rhmcsOosQlikeTmaxPLgbmHar}{}\renewcommand{\rhmcsOosQlikeTmaxPLgbmHar}{0.16}  % MCS p-value, 16:00 bar, T_oos, Qlike, Tmax: LightGBM, HAR
\providecommand{\rhmcsOosQlikeTmaxPLgbmFull}{}\renewcommand{\rhmcsOosQlikeTmaxPLgbmFull}{0.18}  % MCS p-value, 16:00 bar, T_oos, Qlike, Tmax: LightGBM, Full
\providecommand{\rhmcsOosQlikeTmaxPXgbHar}{}\renewcommand{\rhmcsOosQlikeTmaxPXgbHar}{0.16}  % MCS p-value, 16:00 bar, T_oos, Qlike, Tmax: XGBoost, HAR
\providecommand{\rhmcsOosQlikeTmaxPXgbFull}{}\renewcommand{\rhmcsOosQlikeTmaxPXgbFull}{0.18}  % MCS p-value, 16:00 bar, T_oos, Qlike, Tmax: XGBoost, Full
\providecommand{\rhmcsOosQlikeTmaxPRfHar}{}\renewcommand{\rhmcsOosQlikeTmaxPRfHar}{0.043}  % MCS p-value, 16:00 bar, T_oos, Qlike, Tmax: Random forest, HAR
\providecommand{\rhmcsOosQlikeTmaxPRfFull}{}\renewcommand{\rhmcsOosQlikeTmaxPRfFull}{0.14}  % MCS p-value, 16:00 bar, T_oos, Qlike, Tmax: Random forest, Full
\providecommand{\rhmcsOosQlikeTmaxPBench}{}\renewcommand{\rhmcsOosQlikeTmaxPBench}{1.00}  % MCS p-value, 16:00 bar, T_oos, Qlike, Tmax: Closing-interval HAR benchmark
\providecommand{\rhmcsOosQlikeTmaxNNinety}{}\renewcommand{\rhmcsOosQlikeTmaxNNinety}{6}  % models in the 90 % MCS (of 11)
\providecommand{\rhmcsOosQlikeTmaxSetNinety}{}\renewcommand{\rhmcsOosQlikeTmaxSetNinety}{the closing-interval HAR benchmark, XGBoost on the full predictor set, LightGBM on the full predictor set, XGBoost on the HAR benchmark, LightGBM on the HAR benchmark, and the random forest on the full predictor set}  % the 90 % MCS, best (last eliminated) first
\providecommand{\rhmcsOosQlikeTmaxStableNinety}{}\renewcommand{\rhmcsOosQlikeTmaxStableNinety}{yes}  % the 90 % MCS is the same at block lengths (5, 20)
\providecommand{\rhmcsOosQlikeTmaxPMaxOutNinety}{}\renewcommand{\rhmcsOosQlikeTmaxPMaxOutNinety}{0.043}  % largest MCS p-value outside the 90 % set
\providecommand{\rhmcsOosQlikeTmaxNNinetyLo}{}\renewcommand{\rhmcsOosQlikeTmaxNNinetyLo}{6}  % ... at block length 5
\providecommand{\rhmcsOosQlikeTmaxSetNinetyLo}{}\renewcommand{\rhmcsOosQlikeTmaxSetNinetyLo}{the closing-interval HAR benchmark, XGBoost on the full predictor set, LightGBM on the full predictor set, XGBoost on the HAR benchmark, LightGBM on the HAR benchmark, and the random forest on the full predictor set}  % ... at block length 5
\providecommand{\rhmcsOosQlikeTmaxNNinetyHi}{}\renewcommand{\rhmcsOosQlikeTmaxNNinetyHi}{6}  % ... at block length 20
\providecommand{\rhmcsOosQlikeTmaxSetNinetyHi}{}\renewcommand{\rhmcsOosQlikeTmaxSetNinetyHi}{the closing-interval HAR benchmark, XGBoost on the full predictor set, LightGBM on the full predictor set, LightGBM on the HAR benchmark, XGBoost on the HAR benchmark, and the random forest on the full predictor set}  % ... at block length 20
\providecommand{\rhmcsOosQlikeTmaxNSevenFive}{}\renewcommand{\rhmcsOosQlikeTmaxNSevenFive}{1}  % models in the 75 % MCS (of 11)
\providecommand{\rhmcsOosQlikeTmaxSetSevenFive}{}\renewcommand{\rhmcsOosQlikeTmaxSetSevenFive}{the closing-interval HAR benchmark}  % the 75 % MCS, best (last eliminated) first
\providecommand{\rhmcsOosQlikeTmaxStableSevenFive}{}\renewcommand{\rhmcsOosQlikeTmaxStableSevenFive}{yes}  % the 75 % MCS is the same at block lengths (5, 20)
\providecommand{\rhmcsOosQlikeTmaxPMaxOutSevenFive}{}\renewcommand{\rhmcsOosQlikeTmaxPMaxOutSevenFive}{0.18}  % largest MCS p-value outside the 75 % set
\providecommand{\rhmcsOosQlikeTmaxNSevenFiveLo}{}\renewcommand{\rhmcsOosQlikeTmaxNSevenFiveLo}{1}  % ... at block length 5
\providecommand{\rhmcsOosQlikeTmaxSetSevenFiveLo}{}\renewcommand{\rhmcsOosQlikeTmaxSetSevenFiveLo}{the closing-interval HAR benchmark}  % ... at block length 5
\providecommand{\rhmcsOosQlikeTmaxNSevenFiveHi}{}\renewcommand{\rhmcsOosQlikeTmaxNSevenFiveHi}{1}  % ... at block length 20
\providecommand{\rhmcsOosQlikeTmaxSetSevenFiveHi}{}\renewcommand{\rhmcsOosQlikeTmaxSetSevenFiveHi}{the closing-interval HAR benchmark}  % ... at block length 20
\providecommand{\rhmcsOosQlikeTmaxBest}{}\renewcommand{\rhmcsOosQlikeTmaxBest}{the closing-interval HAR benchmark}  % the last model left (MCS p-value 1)
\providecommand{\rhmcsOosQlikeTrangePLinHar}{}\renewcommand{\rhmcsOosQlikeTrangePLinHar}{0.034}  % MCS p-value, 16:00 bar, T_oos, Qlike, Trange: Linear, HAR
\providecommand{\rhmcsOosQlikeTrangePRidgeFull}{}\renewcommand{\rhmcsOosQlikeTrangePRidgeFull}{0.007}  % MCS p-value, 16:00 bar, T_oos, Qlike, Trange: Ridge, Full
\providecommand{\rhmcsOosQlikeTrangePLassoFull}{}\renewcommand{\rhmcsOosQlikeTrangePLassoFull}{0.008}  % MCS p-value, 16:00 bar, T_oos, Qlike, Trange: Lasso, Full
\providecommand{\rhmcsOosQlikeTrangePEnetFull}{}\renewcommand{\rhmcsOosQlikeTrangePEnetFull}{0.013}  % MCS p-value, 16:00 bar, T_oos, Qlike, Trange: Elastic net, Full
\providecommand{\rhmcsOosQlikeTrangePLgbmHar}{}\renewcommand{\rhmcsOosQlikeTrangePLgbmHar}{0.034}  % MCS p-value, 16:00 bar, T_oos, Qlike, Trange: LightGBM, HAR
\providecommand{\rhmcsOosQlikeTrangePLgbmFull}{}\renewcommand{\rhmcsOosQlikeTrangePLgbmFull}{0.29}  % MCS p-value, 16:00 bar, T_oos, Qlike, Trange: LightGBM, Full
\providecommand{\rhmcsOosQlikeTrangePXgbHar}{}\renewcommand{\rhmcsOosQlikeTrangePXgbHar}{0.034}  % MCS p-value, 16:00 bar, T_oos, Qlike, Trange: XGBoost, HAR
\providecommand{\rhmcsOosQlikeTrangePXgbFull}{}\renewcommand{\rhmcsOosQlikeTrangePXgbFull}{0.29}  % MCS p-value, 16:00 bar, T_oos, Qlike, Trange: XGBoost, Full
\providecommand{\rhmcsOosQlikeTrangePRfHar}{}\renewcommand{\rhmcsOosQlikeTrangePRfHar}{0.025}  % MCS p-value, 16:00 bar, T_oos, Qlike, Trange: Random forest, HAR
\providecommand{\rhmcsOosQlikeTrangePRfFull}{}\renewcommand{\rhmcsOosQlikeTrangePRfFull}{0.025}  % MCS p-value, 16:00 bar, T_oos, Qlike, Trange: Random forest, Full
\providecommand{\rhmcsOosQlikeTrangePBench}{}\renewcommand{\rhmcsOosQlikeTrangePBench}{1.00}  % MCS p-value, 16:00 bar, T_oos, Qlike, Trange: Closing-interval HAR benchmark
\providecommand{\rhmcsOosQlikeTrangeNNinety}{}\renewcommand{\rhmcsOosQlikeTrangeNNinety}{3}  % models in the 90 % MCS (of 11)
\providecommand{\rhmcsOosQlikeTrangeSetNinety}{}\renewcommand{\rhmcsOosQlikeTrangeSetNinety}{the closing-interval HAR benchmark, XGBoost on the full predictor set, and LightGBM on the full predictor set}  % the 90 % MCS, best (last eliminated) first
\providecommand{\rhmcsOosQlikeTrangeStableNinety}{}\renewcommand{\rhmcsOosQlikeTrangeStableNinety}{yes}  % the 90 % MCS is the same at block lengths (5, 20)
\providecommand{\rhmcsOosQlikeTrangePMaxOutNinety}{}\renewcommand{\rhmcsOosQlikeTrangePMaxOutNinety}{0.034}  % largest MCS p-value outside the 90 % set
\providecommand{\rhmcsOosQlikeTrangeNNinetyLo}{}\renewcommand{\rhmcsOosQlikeTrangeNNinetyLo}{3}  % ... at block length 5
\providecommand{\rhmcsOosQlikeTrangeSetNinetyLo}{}\renewcommand{\rhmcsOosQlikeTrangeSetNinetyLo}{the closing-interval HAR benchmark, XGBoost on the full predictor set, and LightGBM on the full predictor set}  % ... at block length 5
\providecommand{\rhmcsOosQlikeTrangeNNinetyHi}{}\renewcommand{\rhmcsOosQlikeTrangeNNinetyHi}{3}  % ... at block length 20
\providecommand{\rhmcsOosQlikeTrangeSetNinetyHi}{}\renewcommand{\rhmcsOosQlikeTrangeSetNinetyHi}{the closing-interval HAR benchmark, XGBoost on the full predictor set, and LightGBM on the full predictor set}  % ... at block length 20
\providecommand{\rhmcsOosQlikeTrangeNSevenFive}{}\renewcommand{\rhmcsOosQlikeTrangeNSevenFive}{3}  % models in the 75 % MCS (of 11)
\providecommand{\rhmcsOosQlikeTrangeSetSevenFive}{}\renewcommand{\rhmcsOosQlikeTrangeSetSevenFive}{the closing-interval HAR benchmark, XGBoost on the full predictor set, and LightGBM on the full predictor set}  % the 75 % MCS, best (last eliminated) first
\providecommand{\rhmcsOosQlikeTrangeStableSevenFive}{}\renewcommand{\rhmcsOosQlikeTrangeStableSevenFive}{yes}  % the 75 % MCS is the same at block lengths (5, 20)
\providecommand{\rhmcsOosQlikeTrangePMaxOutSevenFive}{}\renewcommand{\rhmcsOosQlikeTrangePMaxOutSevenFive}{0.034}  % largest MCS p-value outside the 75 % set
\providecommand{\rhmcsOosQlikeTrangeNSevenFiveLo}{}\renewcommand{\rhmcsOosQlikeTrangeNSevenFiveLo}{3}  % ... at block length 5
\providecommand{\rhmcsOosQlikeTrangeSetSevenFiveLo}{}\renewcommand{\rhmcsOosQlikeTrangeSetSevenFiveLo}{the closing-interval HAR benchmark, XGBoost on the full predictor set, and LightGBM on the full predictor set}  % ... at block length 5
\providecommand{\rhmcsOosQlikeTrangeNSevenFiveHi}{}\renewcommand{\rhmcsOosQlikeTrangeNSevenFiveHi}{3}  % ... at block length 20
\providecommand{\rhmcsOosQlikeTrangeSetSevenFiveHi}{}\renewcommand{\rhmcsOosQlikeTrangeSetSevenFiveHi}{the closing-interval HAR benchmark, XGBoost on the full predictor set, and LightGBM on the full predictor set}  % ... at block length 20
\providecommand{\rhmcsOosQlikeTrangeBest}{}\renewcommand{\rhmcsOosQlikeTrangeBest}{the closing-interval HAR benchmark}  % the last model left (MCS p-value 1)
\providecommand{\rhmcsOosMseTmaxPLinHar}{}\renewcommand{\rhmcsOosMseTmaxPLinHar}{0.56}  % MCS p-value, 16:00 bar, T_oos, Mse, Tmax: Linear, HAR
\providecommand{\rhmcsOosMseTmaxPRidgeFull}{}\renewcommand{\rhmcsOosMseTmaxPRidgeFull}{0.44}  % MCS p-value, 16:00 bar, T_oos, Mse, Tmax: Ridge, Full
\providecommand{\rhmcsOosMseTmaxPLassoFull}{}\renewcommand{\rhmcsOosMseTmaxPLassoFull}{0.39}  % MCS p-value, 16:00 bar, T_oos, Mse, Tmax: Lasso, Full
\providecommand{\rhmcsOosMseTmaxPEnetFull}{}\renewcommand{\rhmcsOosMseTmaxPEnetFull}{0.49}  % MCS p-value, 16:00 bar, T_oos, Mse, Tmax: Elastic net, Full
\providecommand{\rhmcsOosMseTmaxPLgbmHar}{}\renewcommand{\rhmcsOosMseTmaxPLgbmHar}{0.56}  % MCS p-value, 16:00 bar, T_oos, Mse, Tmax: LightGBM, HAR
\providecommand{\rhmcsOosMseTmaxPLgbmFull}{}\renewcommand{\rhmcsOosMseTmaxPLgbmFull}{0.54}  % MCS p-value, 16:00 bar, T_oos, Mse, Tmax: LightGBM, Full
\providecommand{\rhmcsOosMseTmaxPXgbHar}{}\renewcommand{\rhmcsOosMseTmaxPXgbHar}{0.56}  % MCS p-value, 16:00 bar, T_oos, Mse, Tmax: XGBoost, HAR
\providecommand{\rhmcsOosMseTmaxPXgbFull}{}\renewcommand{\rhmcsOosMseTmaxPXgbFull}{0.56}  % MCS p-value, 16:00 bar, T_oos, Mse, Tmax: XGBoost, Full
\providecommand{\rhmcsOosMseTmaxPRfHar}{}\renewcommand{\rhmcsOosMseTmaxPRfHar}{1.00}  % MCS p-value, 16:00 bar, T_oos, Mse, Tmax: Random forest, HAR
\providecommand{\rhmcsOosMseTmaxPRfFull}{}\renewcommand{\rhmcsOosMseTmaxPRfFull}{0.56}  % MCS p-value, 16:00 bar, T_oos, Mse, Tmax: Random forest, Full
\providecommand{\rhmcsOosMseTmaxPBench}{}\renewcommand{\rhmcsOosMseTmaxPBench}{0.56}  % MCS p-value, 16:00 bar, T_oos, Mse, Tmax: Closing-interval HAR benchmark
\providecommand{\rhmcsOosMseTmaxNNinety}{}\renewcommand{\rhmcsOosMseTmaxNNinety}{11}  % models in the 90 % MCS (of 11)
\providecommand{\rhmcsOosMseTmaxSetNinety}{}\renewcommand{\rhmcsOosMseTmaxSetNinety}{the random forest on the HAR benchmark, XGBoost on the HAR benchmark, LightGBM on the HAR benchmark, the random forest on the full predictor set, the linear models on the HAR benchmark, the closing-interval HAR benchmark, XGBoost on the full predictor set, LightGBM on the full predictor set, the elastic net on the full predictor set, ridge regression on the full predictor set, and the lasso on the full predictor set}  % the 90 % MCS, best (last eliminated) first
\providecommand{\rhmcsOosMseTmaxStableNinety}{}\renewcommand{\rhmcsOosMseTmaxStableNinety}{yes}  % the 90 % MCS is the same at block lengths (5, 20)
\providecommand{\rhmcsOosMseTmaxPMaxOutNinety}{}\renewcommand{\rhmcsOosMseTmaxPMaxOutNinety}{--}  % largest MCS p-value outside the 90 % set
\providecommand{\rhmcsOosMseTmaxNNinetyLo}{}\renewcommand{\rhmcsOosMseTmaxNNinetyLo}{11}  % ... at block length 5
\providecommand{\rhmcsOosMseTmaxSetNinetyLo}{}\renewcommand{\rhmcsOosMseTmaxSetNinetyLo}{the random forest on the HAR benchmark, XGBoost on the HAR benchmark, LightGBM on the HAR benchmark, the linear models on the HAR benchmark, the random forest on the full predictor set, the closing-interval HAR benchmark, XGBoost on the full predictor set, LightGBM on the full predictor set, the elastic net on the full predictor set, ridge regression on the full predictor set, and the lasso on the full predictor set}  % ... at block length 5
\providecommand{\rhmcsOosMseTmaxNNinetyHi}{}\renewcommand{\rhmcsOosMseTmaxNNinetyHi}{11}  % ... at block length 20
\providecommand{\rhmcsOosMseTmaxSetNinetyHi}{}\renewcommand{\rhmcsOosMseTmaxSetNinetyHi}{the random forest on the HAR benchmark, XGBoost on the HAR benchmark, LightGBM on the HAR benchmark, the linear models on the HAR benchmark, the random forest on the full predictor set, the closing-interval HAR benchmark, XGBoost on the full predictor set, LightGBM on the full predictor set, the elastic net on the full predictor set, ridge regression on the full predictor set, and the lasso on the full predictor set}  % ... at block length 20
\providecommand{\rhmcsOosMseTmaxNSevenFive}{}\renewcommand{\rhmcsOosMseTmaxNSevenFive}{11}  % models in the 75 % MCS (of 11)
\providecommand{\rhmcsOosMseTmaxSetSevenFive}{}\renewcommand{\rhmcsOosMseTmaxSetSevenFive}{the random forest on the HAR benchmark, XGBoost on the HAR benchmark, LightGBM on the HAR benchmark, the random forest on the full predictor set, the linear models on the HAR benchmark, the closing-interval HAR benchmark, XGBoost on the full predictor set, LightGBM on the full predictor set, the elastic net on the full predictor set, ridge regression on the full predictor set, and the lasso on the full predictor set}  % the 75 % MCS, best (last eliminated) first
\providecommand{\rhmcsOosMseTmaxStableSevenFive}{}\renewcommand{\rhmcsOosMseTmaxStableSevenFive}{no}  % the 75 % MCS is the same at block lengths (5, 20)
\providecommand{\rhmcsOosMseTmaxPMaxOutSevenFive}{}\renewcommand{\rhmcsOosMseTmaxPMaxOutSevenFive}{--}  % largest MCS p-value outside the 75 % set
\providecommand{\rhmcsOosMseTmaxNSevenFiveLo}{}\renewcommand{\rhmcsOosMseTmaxNSevenFiveLo}{10}  % ... at block length 5
\providecommand{\rhmcsOosMseTmaxSetSevenFiveLo}{}\renewcommand{\rhmcsOosMseTmaxSetSevenFiveLo}{the random forest on the HAR benchmark, XGBoost on the HAR benchmark, LightGBM on the HAR benchmark, the linear models on the HAR benchmark, the random forest on the full predictor set, the closing-interval HAR benchmark, XGBoost on the full predictor set, LightGBM on the full predictor set, the elastic net on the full predictor set, and ridge regression on the full predictor set}  % ... at block length 5
\providecommand{\rhmcsOosMseTmaxNSevenFiveHi}{}\renewcommand{\rhmcsOosMseTmaxNSevenFiveHi}{11}  % ... at block length 20
\providecommand{\rhmcsOosMseTmaxSetSevenFiveHi}{}\renewcommand{\rhmcsOosMseTmaxSetSevenFiveHi}{the random forest on the HAR benchmark, XGBoost on the HAR benchmark, LightGBM on the HAR benchmark, the linear models on the HAR benchmark, the random forest on the full predictor set, the closing-interval HAR benchmark, XGBoost on the full predictor set, LightGBM on the full predictor set, the elastic net on the full predictor set, ridge regression on the full predictor set, and the lasso on the full predictor set}  % ... at block length 20
\providecommand{\rhmcsOosMseTmaxBest}{}\renewcommand{\rhmcsOosMseTmaxBest}{the random forest on the HAR benchmark}  % the last model left (MCS p-value 1)
\providecommand{\rhmcsOosMseTrangePLinHar}{}\renewcommand{\rhmcsOosMseTrangePLinHar}{0.56}  % MCS p-value, 16:00 bar, T_oos, Mse, Trange: Linear, HAR
\providecommand{\rhmcsOosMseTrangePRidgeFull}{}\renewcommand{\rhmcsOosMseTrangePRidgeFull}{0.56}  % MCS p-value, 16:00 bar, T_oos, Mse, Trange: Ridge, Full
\providecommand{\rhmcsOosMseTrangePLassoFull}{}\renewcommand{\rhmcsOosMseTrangePLassoFull}{0.56}  % MCS p-value, 16:00 bar, T_oos, Mse, Trange: Lasso, Full
\providecommand{\rhmcsOosMseTrangePEnetFull}{}\renewcommand{\rhmcsOosMseTrangePEnetFull}{0.56}  % MCS p-value, 16:00 bar, T_oos, Mse, Trange: Elastic net, Full
\providecommand{\rhmcsOosMseTrangePLgbmHar}{}\renewcommand{\rhmcsOosMseTrangePLgbmHar}{0.56}  % MCS p-value, 16:00 bar, T_oos, Mse, Trange: LightGBM, HAR
\providecommand{\rhmcsOosMseTrangePLgbmFull}{}\renewcommand{\rhmcsOosMseTrangePLgbmFull}{0.56}  % MCS p-value, 16:00 bar, T_oos, Mse, Trange: LightGBM, Full
\providecommand{\rhmcsOosMseTrangePXgbHar}{}\renewcommand{\rhmcsOosMseTrangePXgbHar}{0.56}  % MCS p-value, 16:00 bar, T_oos, Mse, Trange: XGBoost, HAR
\providecommand{\rhmcsOosMseTrangePXgbFull}{}\renewcommand{\rhmcsOosMseTrangePXgbFull}{0.56}  % MCS p-value, 16:00 bar, T_oos, Mse, Trange: XGBoost, Full
\providecommand{\rhmcsOosMseTrangePRfHar}{}\renewcommand{\rhmcsOosMseTrangePRfHar}{1.00}  % MCS p-value, 16:00 bar, T_oos, Mse, Trange: Random forest, HAR
\providecommand{\rhmcsOosMseTrangePRfFull}{}\renewcommand{\rhmcsOosMseTrangePRfFull}{0.56}  % MCS p-value, 16:00 bar, T_oos, Mse, Trange: Random forest, Full
\providecommand{\rhmcsOosMseTrangePBench}{}\renewcommand{\rhmcsOosMseTrangePBench}{0.56}  % MCS p-value, 16:00 bar, T_oos, Mse, Trange: Closing-interval HAR benchmark
\providecommand{\rhmcsOosMseTrangeNNinety}{}\renewcommand{\rhmcsOosMseTrangeNNinety}{11}  % models in the 90 % MCS (of 11)
\providecommand{\rhmcsOosMseTrangeSetNinety}{}\renewcommand{\rhmcsOosMseTrangeSetNinety}{the random forest on the HAR benchmark, XGBoost on the HAR benchmark, the linear models on the HAR benchmark, the random forest on the full predictor set, XGBoost on the full predictor set, LightGBM on the full predictor set, the elastic net on the full predictor set, ridge regression on the full predictor set, the lasso on the full predictor set, the closing-interval HAR benchmark, and LightGBM on the HAR benchmark}  % the 90 % MCS, best (last eliminated) first
\providecommand{\rhmcsOosMseTrangeStableNinety}{}\renewcommand{\rhmcsOosMseTrangeStableNinety}{yes}  % the 90 % MCS is the same at block lengths (5, 20)
\providecommand{\rhmcsOosMseTrangePMaxOutNinety}{}\renewcommand{\rhmcsOosMseTrangePMaxOutNinety}{--}  % largest MCS p-value outside the 90 % set
\providecommand{\rhmcsOosMseTrangeNNinetyLo}{}\renewcommand{\rhmcsOosMseTrangeNNinetyLo}{11}  % ... at block length 5
\providecommand{\rhmcsOosMseTrangeSetNinetyLo}{}\renewcommand{\rhmcsOosMseTrangeSetNinetyLo}{the random forest on the HAR benchmark, XGBoost on the HAR benchmark, the linear models on the HAR benchmark, XGBoost on the full predictor set, the random forest on the full predictor set, the elastic net on the full predictor set, LightGBM on the full predictor set, LightGBM on the HAR benchmark, ridge regression on the full predictor set, the lasso on the full predictor set, and the closing-interval HAR benchmark}  % ... at block length 5
\providecommand{\rhmcsOosMseTrangeNNinetyHi}{}\renewcommand{\rhmcsOosMseTrangeNNinetyHi}{11}  % ... at block length 20
\providecommand{\rhmcsOosMseTrangeSetNinetyHi}{}\renewcommand{\rhmcsOosMseTrangeSetNinetyHi}{the random forest on the HAR benchmark, the linear models on the HAR benchmark, the random forest on the full predictor set, XGBoost on the HAR benchmark, XGBoost on the full predictor set, LightGBM on the full predictor set, the elastic net on the full predictor set, ridge regression on the full predictor set, the lasso on the full predictor set, the closing-interval HAR benchmark, and LightGBM on the HAR benchmark}  % ... at block length 20
\providecommand{\rhmcsOosMseTrangeNSevenFive}{}\renewcommand{\rhmcsOosMseTrangeNSevenFive}{11}  % models in the 75 % MCS (of 11)
\providecommand{\rhmcsOosMseTrangeSetSevenFive}{}\renewcommand{\rhmcsOosMseTrangeSetSevenFive}{the random forest on the HAR benchmark, XGBoost on the HAR benchmark, the linear models on the HAR benchmark, the random forest on the full predictor set, XGBoost on the full predictor set, LightGBM on the full predictor set, the elastic net on the full predictor set, ridge regression on the full predictor set, the lasso on the full predictor set, the closing-interval HAR benchmark, and LightGBM on the HAR benchmark}  % the 75 % MCS, best (last eliminated) first
\providecommand{\rhmcsOosMseTrangeStableSevenFive}{}\renewcommand{\rhmcsOosMseTrangeStableSevenFive}{yes}  % the 75 % MCS is the same at block lengths (5, 20)
\providecommand{\rhmcsOosMseTrangePMaxOutSevenFive}{}\renewcommand{\rhmcsOosMseTrangePMaxOutSevenFive}{--}  % largest MCS p-value outside the 75 % set
\providecommand{\rhmcsOosMseTrangeNSevenFiveLo}{}\renewcommand{\rhmcsOosMseTrangeNSevenFiveLo}{11}  % ... at block length 5
\providecommand{\rhmcsOosMseTrangeSetSevenFiveLo}{}\renewcommand{\rhmcsOosMseTrangeSetSevenFiveLo}{the random forest on the HAR benchmark, XGBoost on the HAR benchmark, the linear models on the HAR benchmark, XGBoost on the full predictor set, the random forest on the full predictor set, the elastic net on the full predictor set, LightGBM on the full predictor set, LightGBM on the HAR benchmark, ridge regression on the full predictor set, the lasso on the full predictor set, and the closing-interval HAR benchmark}  % ... at block length 5
\providecommand{\rhmcsOosMseTrangeNSevenFiveHi}{}\renewcommand{\rhmcsOosMseTrangeNSevenFiveHi}{11}  % ... at block length 20
\providecommand{\rhmcsOosMseTrangeSetSevenFiveHi}{}\renewcommand{\rhmcsOosMseTrangeSetSevenFiveHi}{the random forest on the HAR benchmark, the linear models on the HAR benchmark, the random forest on the full predictor set, XGBoost on the HAR benchmark, XGBoost on the full predictor set, LightGBM on the full predictor set, the elastic net on the full predictor set, ridge regression on the full predictor set, the lasso on the full predictor set, the closing-interval HAR benchmark, and LightGBM on the HAR benchmark}  % ... at block length 20
\providecommand{\rhmcsOosMseTrangeBest}{}\renewcommand{\rhmcsOosMseTrangeBest}{the random forest on the HAR benchmark}  % the last model left (MCS p-value 1)
\providecommand{\rhmcsChQlikeTmaxPLinHar}{}\renewcommand{\rhmcsChQlikeTmaxPLinHar}{0.050}  % MCS p-value, 16:00 bar, the benchmark's span, Qlike, Tmax: Linear, HAR
\providecommand{\rhmcsChQlikeTmaxPRidgeFull}{}\renewcommand{\rhmcsChQlikeTmaxPRidgeFull}{0.050}  % MCS p-value, 16:00 bar, the benchmark's span, Qlike, Tmax: Ridge, Full
\providecommand{\rhmcsChQlikeTmaxPLassoFull}{}\renewcommand{\rhmcsChQlikeTmaxPLassoFull}{0.050}  % MCS p-value, 16:00 bar, the benchmark's span, Qlike, Tmax: Lasso, Full
\providecommand{\rhmcsChQlikeTmaxPEnetFull}{}\renewcommand{\rhmcsChQlikeTmaxPEnetFull}{0.102}  % MCS p-value, 16:00 bar, the benchmark's span, Qlike, Tmax: Elastic net, Full
\providecommand{\rhmcsChQlikeTmaxPLgbmHar}{}\renewcommand{\rhmcsChQlikeTmaxPLgbmHar}{0.102}  % MCS p-value, 16:00 bar, the benchmark's span, Qlike, Tmax: LightGBM, HAR
\providecommand{\rhmcsChQlikeTmaxPLgbmFull}{}\renewcommand{\rhmcsChQlikeTmaxPLgbmFull}{0.16}  % MCS p-value, 16:00 bar, the benchmark's span, Qlike, Tmax: LightGBM, Full
\providecommand{\rhmcsChQlikeTmaxPXgbHar}{}\renewcommand{\rhmcsChQlikeTmaxPXgbHar}{0.102}  % MCS p-value, 16:00 bar, the benchmark's span, Qlike, Tmax: XGBoost, HAR
\providecommand{\rhmcsChQlikeTmaxPXgbFull}{}\renewcommand{\rhmcsChQlikeTmaxPXgbFull}{0.22}  % MCS p-value, 16:00 bar, the benchmark's span, Qlike, Tmax: XGBoost, Full
\providecommand{\rhmcsChQlikeTmaxPRfHar}{}\renewcommand{\rhmcsChQlikeTmaxPRfHar}{0.050}  % MCS p-value, 16:00 bar, the benchmark's span, Qlike, Tmax: Random forest, HAR
\providecommand{\rhmcsChQlikeTmaxPRfFull}{}\renewcommand{\rhmcsChQlikeTmaxPRfFull}{0.102}  % MCS p-value, 16:00 bar, the benchmark's span, Qlike, Tmax: Random forest, Full
\providecommand{\rhmcsChQlikeTmaxPBench}{}\renewcommand{\rhmcsChQlikeTmaxPBench}{1.00}  % MCS p-value, 16:00 bar, the benchmark's span, Qlike, Tmax: Closing-interval HAR benchmark
\providecommand{\rhmcsChQlikeTmaxNNinety}{}\renewcommand{\rhmcsChQlikeTmaxNNinety}{7}  % models in the 90 % MCS (of 11)
\providecommand{\rhmcsChQlikeTmaxSetNinety}{}\renewcommand{\rhmcsChQlikeTmaxSetNinety}{the closing-interval HAR benchmark, XGBoost on the full predictor set, LightGBM on the full predictor set, LightGBM on the HAR benchmark, XGBoost on the HAR benchmark, the random forest on the full predictor set, and the elastic net on the full predictor set}  % the 90 % MCS, best (last eliminated) first
\providecommand{\rhmcsChQlikeTmaxStableNinety}{}\renewcommand{\rhmcsChQlikeTmaxStableNinety}{no}  % the 90 % MCS is the same at block lengths (5, 20)
\providecommand{\rhmcsChQlikeTmaxPMaxOutNinety}{}\renewcommand{\rhmcsChQlikeTmaxPMaxOutNinety}{0.050}  % largest MCS p-value outside the 90 % set
\providecommand{\rhmcsChQlikeTmaxNNinetyLo}{}\renewcommand{\rhmcsChQlikeTmaxNNinetyLo}{3}  % ... at block length 5
\providecommand{\rhmcsChQlikeTmaxSetNinetyLo}{}\renewcommand{\rhmcsChQlikeTmaxSetNinetyLo}{the closing-interval HAR benchmark, XGBoost on the full predictor set, and LightGBM on the full predictor set}  % ... at block length 5
\providecommand{\rhmcsChQlikeTmaxNNinetyHi}{}\renewcommand{\rhmcsChQlikeTmaxNNinetyHi}{7}  % ... at block length 20
\providecommand{\rhmcsChQlikeTmaxSetNinetyHi}{}\renewcommand{\rhmcsChQlikeTmaxSetNinetyHi}{the closing-interval HAR benchmark, XGBoost on the full predictor set, LightGBM on the full predictor set, LightGBM on the HAR benchmark, XGBoost on the HAR benchmark, the random forest on the full predictor set, and the elastic net on the full predictor set}  % ... at block length 20
\providecommand{\rhmcsChQlikeTmaxNSevenFive}{}\renewcommand{\rhmcsChQlikeTmaxNSevenFive}{1}  % models in the 75 % MCS (of 11)
\providecommand{\rhmcsChQlikeTmaxSetSevenFive}{}\renewcommand{\rhmcsChQlikeTmaxSetSevenFive}{the closing-interval HAR benchmark}  % the 75 % MCS, best (last eliminated) first
\providecommand{\rhmcsChQlikeTmaxStableSevenFive}{}\renewcommand{\rhmcsChQlikeTmaxStableSevenFive}{yes}  % the 75 % MCS is the same at block lengths (5, 20)
\providecommand{\rhmcsChQlikeTmaxPMaxOutSevenFive}{}\renewcommand{\rhmcsChQlikeTmaxPMaxOutSevenFive}{0.22}  % largest MCS p-value outside the 75 % set
\providecommand{\rhmcsChQlikeTmaxNSevenFiveLo}{}\renewcommand{\rhmcsChQlikeTmaxNSevenFiveLo}{1}  % ... at block length 5
\providecommand{\rhmcsChQlikeTmaxSetSevenFiveLo}{}\renewcommand{\rhmcsChQlikeTmaxSetSevenFiveLo}{the closing-interval HAR benchmark}  % ... at block length 5
\providecommand{\rhmcsChQlikeTmaxNSevenFiveHi}{}\renewcommand{\rhmcsChQlikeTmaxNSevenFiveHi}{1}  % ... at block length 20
\providecommand{\rhmcsChQlikeTmaxSetSevenFiveHi}{}\renewcommand{\rhmcsChQlikeTmaxSetSevenFiveHi}{the closing-interval HAR benchmark}  % ... at block length 20
\providecommand{\rhmcsChQlikeTmaxBest}{}\renewcommand{\rhmcsChQlikeTmaxBest}{the closing-interval HAR benchmark}  % the last model left (MCS p-value 1)
\providecommand{\rhmcsChQlikeTrangePLinHar}{}\renewcommand{\rhmcsChQlikeTrangePLinHar}{\ensuremath{<}0.001}  % MCS p-value, 16:00 bar, the benchmark's span, Qlike, Trange: Linear, HAR
\providecommand{\rhmcsChQlikeTrangePRidgeFull}{}\renewcommand{\rhmcsChQlikeTrangePRidgeFull}{0.008}  % MCS p-value, 16:00 bar, the benchmark's span, Qlike, Trange: Ridge, Full
\providecommand{\rhmcsChQlikeTrangePLassoFull}{}\renewcommand{\rhmcsChQlikeTrangePLassoFull}{0.007}  % MCS p-value, 16:00 bar, the benchmark's span, Qlike, Trange: Lasso, Full
\providecommand{\rhmcsChQlikeTrangePEnetFull}{}\renewcommand{\rhmcsChQlikeTrangePEnetFull}{0.008}  % MCS p-value, 16:00 bar, the benchmark's span, Qlike, Trange: Elastic net, Full
\providecommand{\rhmcsChQlikeTrangePLgbmHar}{}\renewcommand{\rhmcsChQlikeTrangePLgbmHar}{0.008}  % MCS p-value, 16:00 bar, the benchmark's span, Qlike, Trange: LightGBM, HAR
\providecommand{\rhmcsChQlikeTrangePLgbmFull}{}\renewcommand{\rhmcsChQlikeTrangePLgbmFull}{0.31}  % MCS p-value, 16:00 bar, the benchmark's span, Qlike, Trange: LightGBM, Full
\providecommand{\rhmcsChQlikeTrangePXgbHar}{}\renewcommand{\rhmcsChQlikeTrangePXgbHar}{0.007}  % MCS p-value, 16:00 bar, the benchmark's span, Qlike, Trange: XGBoost, HAR
\providecommand{\rhmcsChQlikeTrangePXgbFull}{}\renewcommand{\rhmcsChQlikeTrangePXgbFull}{0.31}  % MCS p-value, 16:00 bar, the benchmark's span, Qlike, Trange: XGBoost, Full
\providecommand{\rhmcsChQlikeTrangePRfHar}{}\renewcommand{\rhmcsChQlikeTrangePRfHar}{\ensuremath{<}0.001}  % MCS p-value, 16:00 bar, the benchmark's span, Qlike, Trange: Random forest, HAR
\providecommand{\rhmcsChQlikeTrangePRfFull}{}\renewcommand{\rhmcsChQlikeTrangePRfFull}{\ensuremath{<}0.001}  % MCS p-value, 16:00 bar, the benchmark's span, Qlike, Trange: Random forest, Full
\providecommand{\rhmcsChQlikeTrangePBench}{}\renewcommand{\rhmcsChQlikeTrangePBench}{1.00}  % MCS p-value, 16:00 bar, the benchmark's span, Qlike, Trange: Closing-interval HAR benchmark
\providecommand{\rhmcsChQlikeTrangeNNinety}{}\renewcommand{\rhmcsChQlikeTrangeNNinety}{3}  % models in the 90 % MCS (of 11)
\providecommand{\rhmcsChQlikeTrangeSetNinety}{}\renewcommand{\rhmcsChQlikeTrangeSetNinety}{the closing-interval HAR benchmark, XGBoost on the full predictor set, and LightGBM on the full predictor set}  % the 90 % MCS, best (last eliminated) first
\providecommand{\rhmcsChQlikeTrangeStableNinety}{}\renewcommand{\rhmcsChQlikeTrangeStableNinety}{yes}  % the 90 % MCS is the same at block lengths (5, 20)
\providecommand{\rhmcsChQlikeTrangePMaxOutNinety}{}\renewcommand{\rhmcsChQlikeTrangePMaxOutNinety}{0.008}  % largest MCS p-value outside the 90 % set
\providecommand{\rhmcsChQlikeTrangeNNinetyLo}{}\renewcommand{\rhmcsChQlikeTrangeNNinetyLo}{3}  % ... at block length 5
\providecommand{\rhmcsChQlikeTrangeSetNinetyLo}{}\renewcommand{\rhmcsChQlikeTrangeSetNinetyLo}{the closing-interval HAR benchmark, XGBoost on the full predictor set, and LightGBM on the full predictor set}  % ... at block length 5
\providecommand{\rhmcsChQlikeTrangeNNinetyHi}{}\renewcommand{\rhmcsChQlikeTrangeNNinetyHi}{3}  % ... at block length 20
\providecommand{\rhmcsChQlikeTrangeSetNinetyHi}{}\renewcommand{\rhmcsChQlikeTrangeSetNinetyHi}{the closing-interval HAR benchmark, XGBoost on the full predictor set, and LightGBM on the full predictor set}  % ... at block length 20
\providecommand{\rhmcsChQlikeTrangeNSevenFive}{}\renewcommand{\rhmcsChQlikeTrangeNSevenFive}{3}  % models in the 75 % MCS (of 11)
\providecommand{\rhmcsChQlikeTrangeSetSevenFive}{}\renewcommand{\rhmcsChQlikeTrangeSetSevenFive}{the closing-interval HAR benchmark, XGBoost on the full predictor set, and LightGBM on the full predictor set}  % the 75 % MCS, best (last eliminated) first
\providecommand{\rhmcsChQlikeTrangeStableSevenFive}{}\renewcommand{\rhmcsChQlikeTrangeStableSevenFive}{yes}  % the 75 % MCS is the same at block lengths (5, 20)
\providecommand{\rhmcsChQlikeTrangePMaxOutSevenFive}{}\renewcommand{\rhmcsChQlikeTrangePMaxOutSevenFive}{0.008}  % largest MCS p-value outside the 75 % set
\providecommand{\rhmcsChQlikeTrangeNSevenFiveLo}{}\renewcommand{\rhmcsChQlikeTrangeNSevenFiveLo}{3}  % ... at block length 5
\providecommand{\rhmcsChQlikeTrangeSetSevenFiveLo}{}\renewcommand{\rhmcsChQlikeTrangeSetSevenFiveLo}{the closing-interval HAR benchmark, XGBoost on the full predictor set, and LightGBM on the full predictor set}  % ... at block length 5
\providecommand{\rhmcsChQlikeTrangeNSevenFiveHi}{}\renewcommand{\rhmcsChQlikeTrangeNSevenFiveHi}{3}  % ... at block length 20
\providecommand{\rhmcsChQlikeTrangeSetSevenFiveHi}{}\renewcommand{\rhmcsChQlikeTrangeSetSevenFiveHi}{the closing-interval HAR benchmark, XGBoost on the full predictor set, and LightGBM on the full predictor set}  % ... at block length 20
\providecommand{\rhmcsChQlikeTrangeBest}{}\renewcommand{\rhmcsChQlikeTrangeBest}{the closing-interval HAR benchmark}  % the last model left (MCS p-value 1)
\providecommand{\rhmcsChMseTmaxPLinHar}{}\renewcommand{\rhmcsChMseTmaxPLinHar}{0.59}  % MCS p-value, 16:00 bar, the benchmark's span, Mse, Tmax: Linear, HAR
\providecommand{\rhmcsChMseTmaxPRidgeFull}{}\renewcommand{\rhmcsChMseTmaxPRidgeFull}{0.47}  % MCS p-value, 16:00 bar, the benchmark's span, Mse, Tmax: Ridge, Full
\providecommand{\rhmcsChMseTmaxPLassoFull}{}\renewcommand{\rhmcsChMseTmaxPLassoFull}{0.42}  % MCS p-value, 16:00 bar, the benchmark's span, Mse, Tmax: Lasso, Full
\providecommand{\rhmcsChMseTmaxPEnetFull}{}\renewcommand{\rhmcsChMseTmaxPEnetFull}{0.53}  % MCS p-value, 16:00 bar, the benchmark's span, Mse, Tmax: Elastic net, Full
\providecommand{\rhmcsChMseTmaxPLgbmHar}{}\renewcommand{\rhmcsChMseTmaxPLgbmHar}{0.59}  % MCS p-value, 16:00 bar, the benchmark's span, Mse, Tmax: LightGBM, HAR
\providecommand{\rhmcsChMseTmaxPLgbmFull}{}\renewcommand{\rhmcsChMseTmaxPLgbmFull}{0.54}  % MCS p-value, 16:00 bar, the benchmark's span, Mse, Tmax: LightGBM, Full
\providecommand{\rhmcsChMseTmaxPXgbHar}{}\renewcommand{\rhmcsChMseTmaxPXgbHar}{0.59}  % MCS p-value, 16:00 bar, the benchmark's span, Mse, Tmax: XGBoost, HAR
\providecommand{\rhmcsChMseTmaxPXgbFull}{}\renewcommand{\rhmcsChMseTmaxPXgbFull}{0.58}  % MCS p-value, 16:00 bar, the benchmark's span, Mse, Tmax: XGBoost, Full
\providecommand{\rhmcsChMseTmaxPRfHar}{}\renewcommand{\rhmcsChMseTmaxPRfHar}{1.00}  % MCS p-value, 16:00 bar, the benchmark's span, Mse, Tmax: Random forest, HAR
\providecommand{\rhmcsChMseTmaxPRfFull}{}\renewcommand{\rhmcsChMseTmaxPRfFull}{0.59}  % MCS p-value, 16:00 bar, the benchmark's span, Mse, Tmax: Random forest, Full
\providecommand{\rhmcsChMseTmaxPBench}{}\renewcommand{\rhmcsChMseTmaxPBench}{0.58}  % MCS p-value, 16:00 bar, the benchmark's span, Mse, Tmax: Closing-interval HAR benchmark
\providecommand{\rhmcsChMseTmaxNNinety}{}\renewcommand{\rhmcsChMseTmaxNNinety}{11}  % models in the 90 % MCS (of 11)
\providecommand{\rhmcsChMseTmaxSetNinety}{}\renewcommand{\rhmcsChMseTmaxSetNinety}{the random forest on the HAR benchmark, XGBoost on the HAR benchmark, LightGBM on the HAR benchmark, the linear models on the HAR benchmark, the random forest on the full predictor set, the closing-interval HAR benchmark, XGBoost on the full predictor set, LightGBM on the full predictor set, the elastic net on the full predictor set, ridge regression on the full predictor set, and the lasso on the full predictor set}  % the 90 % MCS, best (last eliminated) first
\providecommand{\rhmcsChMseTmaxStableNinety}{}\renewcommand{\rhmcsChMseTmaxStableNinety}{yes}  % the 90 % MCS is the same at block lengths (5, 20)
\providecommand{\rhmcsChMseTmaxPMaxOutNinety}{}\renewcommand{\rhmcsChMseTmaxPMaxOutNinety}{--}  % largest MCS p-value outside the 90 % set
\providecommand{\rhmcsChMseTmaxNNinetyLo}{}\renewcommand{\rhmcsChMseTmaxNNinetyLo}{11}  % ... at block length 5
\providecommand{\rhmcsChMseTmaxSetNinetyLo}{}\renewcommand{\rhmcsChMseTmaxSetNinetyLo}{the random forest on the HAR benchmark, XGBoost on the HAR benchmark, LightGBM on the HAR benchmark, the linear models on the HAR benchmark, the random forest on the full predictor set, the closing-interval HAR benchmark, XGBoost on the full predictor set, LightGBM on the full predictor set, the elastic net on the full predictor set, ridge regression on the full predictor set, and the lasso on the full predictor set}  % ... at block length 5
\providecommand{\rhmcsChMseTmaxNNinetyHi}{}\renewcommand{\rhmcsChMseTmaxNNinetyHi}{11}  % ... at block length 20
\providecommand{\rhmcsChMseTmaxSetNinetyHi}{}\renewcommand{\rhmcsChMseTmaxSetNinetyHi}{the random forest on the HAR benchmark, XGBoost on the HAR benchmark, LightGBM on the HAR benchmark, the linear models on the HAR benchmark, the random forest on the full predictor set, the closing-interval HAR benchmark, XGBoost on the full predictor set, LightGBM on the full predictor set, the elastic net on the full predictor set, ridge regression on the full predictor set, and the lasso on the full predictor set}  % ... at block length 20
\providecommand{\rhmcsChMseTmaxNSevenFive}{}\renewcommand{\rhmcsChMseTmaxNSevenFive}{11}  % models in the 75 % MCS (of 11)
\providecommand{\rhmcsChMseTmaxSetSevenFive}{}\renewcommand{\rhmcsChMseTmaxSetSevenFive}{the random forest on the HAR benchmark, XGBoost on the HAR benchmark, LightGBM on the HAR benchmark, the linear models on the HAR benchmark, the random forest on the full predictor set, the closing-interval HAR benchmark, XGBoost on the full predictor set, LightGBM on the full predictor set, the elastic net on the full predictor set, ridge regression on the full predictor set, and the lasso on the full predictor set}  % the 75 % MCS, best (last eliminated) first
\providecommand{\rhmcsChMseTmaxStableSevenFive}{}\renewcommand{\rhmcsChMseTmaxStableSevenFive}{yes}  % the 75 % MCS is the same at block lengths (5, 20)
\providecommand{\rhmcsChMseTmaxPMaxOutSevenFive}{}\renewcommand{\rhmcsChMseTmaxPMaxOutSevenFive}{--}  % largest MCS p-value outside the 75 % set
\providecommand{\rhmcsChMseTmaxNSevenFiveLo}{}\renewcommand{\rhmcsChMseTmaxNSevenFiveLo}{11}  % ... at block length 5
\providecommand{\rhmcsChMseTmaxSetSevenFiveLo}{}\renewcommand{\rhmcsChMseTmaxSetSevenFiveLo}{the random forest on the HAR benchmark, XGBoost on the HAR benchmark, LightGBM on the HAR benchmark, the linear models on the HAR benchmark, the random forest on the full predictor set, the closing-interval HAR benchmark, XGBoost on the full predictor set, LightGBM on the full predictor set, the elastic net on the full predictor set, ridge regression on the full predictor set, and the lasso on the full predictor set}  % ... at block length 5
\providecommand{\rhmcsChMseTmaxNSevenFiveHi}{}\renewcommand{\rhmcsChMseTmaxNSevenFiveHi}{11}  % ... at block length 20
\providecommand{\rhmcsChMseTmaxSetSevenFiveHi}{}\renewcommand{\rhmcsChMseTmaxSetSevenFiveHi}{the random forest on the HAR benchmark, XGBoost on the HAR benchmark, LightGBM on the HAR benchmark, the linear models on the HAR benchmark, the random forest on the full predictor set, the closing-interval HAR benchmark, XGBoost on the full predictor set, LightGBM on the full predictor set, the elastic net on the full predictor set, ridge regression on the full predictor set, and the lasso on the full predictor set}  % ... at block length 20
\providecommand{\rhmcsChMseTmaxBest}{}\renewcommand{\rhmcsChMseTmaxBest}{the random forest on the HAR benchmark}  % the last model left (MCS p-value 1)
\providecommand{\rhmcsChMseTrangePLinHar}{}\renewcommand{\rhmcsChMseTrangePLinHar}{0.60}  % MCS p-value, 16:00 bar, the benchmark's span, Mse, Trange: Linear, HAR
\providecommand{\rhmcsChMseTrangePRidgeFull}{}\renewcommand{\rhmcsChMseTrangePRidgeFull}{0.60}  % MCS p-value, 16:00 bar, the benchmark's span, Mse, Trange: Ridge, Full
\providecommand{\rhmcsChMseTrangePLassoFull}{}\renewcommand{\rhmcsChMseTrangePLassoFull}{0.60}  % MCS p-value, 16:00 bar, the benchmark's span, Mse, Trange: Lasso, Full
\providecommand{\rhmcsChMseTrangePEnetFull}{}\renewcommand{\rhmcsChMseTrangePEnetFull}{0.60}  % MCS p-value, 16:00 bar, the benchmark's span, Mse, Trange: Elastic net, Full
\providecommand{\rhmcsChMseTrangePLgbmHar}{}\renewcommand{\rhmcsChMseTrangePLgbmHar}{0.41}  % MCS p-value, 16:00 bar, the benchmark's span, Mse, Trange: LightGBM, HAR
\providecommand{\rhmcsChMseTrangePLgbmFull}{}\renewcommand{\rhmcsChMseTrangePLgbmFull}{0.60}  % MCS p-value, 16:00 bar, the benchmark's span, Mse, Trange: LightGBM, Full
\providecommand{\rhmcsChMseTrangePXgbHar}{}\renewcommand{\rhmcsChMseTrangePXgbHar}{0.60}  % MCS p-value, 16:00 bar, the benchmark's span, Mse, Trange: XGBoost, HAR
\providecommand{\rhmcsChMseTrangePXgbFull}{}\renewcommand{\rhmcsChMseTrangePXgbFull}{0.60}  % MCS p-value, 16:00 bar, the benchmark's span, Mse, Trange: XGBoost, Full
\providecommand{\rhmcsChMseTrangePRfHar}{}\renewcommand{\rhmcsChMseTrangePRfHar}{1.00}  % MCS p-value, 16:00 bar, the benchmark's span, Mse, Trange: Random forest, HAR
\providecommand{\rhmcsChMseTrangePRfFull}{}\renewcommand{\rhmcsChMseTrangePRfFull}{0.60}  % MCS p-value, 16:00 bar, the benchmark's span, Mse, Trange: Random forest, Full
\providecommand{\rhmcsChMseTrangePBench}{}\renewcommand{\rhmcsChMseTrangePBench}{0.60}  % MCS p-value, 16:00 bar, the benchmark's span, Mse, Trange: Closing-interval HAR benchmark
\providecommand{\rhmcsChMseTrangeNNinety}{}\renewcommand{\rhmcsChMseTrangeNNinety}{11}  % models in the 90 % MCS (of 11)
\providecommand{\rhmcsChMseTrangeSetNinety}{}\renewcommand{\rhmcsChMseTrangeSetNinety}{the random forest on the HAR benchmark, the linear models on the HAR benchmark, XGBoost on the HAR benchmark, the random forest on the full predictor set, XGBoost on the full predictor set, the elastic net on the full predictor set, ridge regression on the full predictor set, the lasso on the full predictor set, LightGBM on the full predictor set, the closing-interval HAR benchmark, and LightGBM on the HAR benchmark}  % the 90 % MCS, best (last eliminated) first
\providecommand{\rhmcsChMseTrangeStableNinety}{}\renewcommand{\rhmcsChMseTrangeStableNinety}{yes}  % the 90 % MCS is the same at block lengths (5, 20)
\providecommand{\rhmcsChMseTrangePMaxOutNinety}{}\renewcommand{\rhmcsChMseTrangePMaxOutNinety}{--}  % largest MCS p-value outside the 90 % set
\providecommand{\rhmcsChMseTrangeNNinetyLo}{}\renewcommand{\rhmcsChMseTrangeNNinetyLo}{11}  % ... at block length 5
\providecommand{\rhmcsChMseTrangeSetNinetyLo}{}\renewcommand{\rhmcsChMseTrangeSetNinetyLo}{the random forest on the HAR benchmark, the linear models on the HAR benchmark, XGBoost on the HAR benchmark, XGBoost on the full predictor set, the random forest on the full predictor set, the elastic net on the full predictor set, ridge regression on the full predictor set, LightGBM on the full predictor set, the lasso on the full predictor set, the closing-interval HAR benchmark, and LightGBM on the HAR benchmark}  % ... at block length 5
\providecommand{\rhmcsChMseTrangeNNinetyHi}{}\renewcommand{\rhmcsChMseTrangeNNinetyHi}{11}  % ... at block length 20
\providecommand{\rhmcsChMseTrangeSetNinetyHi}{}\renewcommand{\rhmcsChMseTrangeSetNinetyHi}{the random forest on the HAR benchmark, the linear models on the HAR benchmark, the random forest on the full predictor set, XGBoost on the full predictor set, the elastic net on the full predictor set, ridge regression on the full predictor set, the lasso on the full predictor set, XGBoost on the HAR benchmark, the closing-interval HAR benchmark, LightGBM on the full predictor set, and LightGBM on the HAR benchmark}  % ... at block length 20
\providecommand{\rhmcsChMseTrangeNSevenFive}{}\renewcommand{\rhmcsChMseTrangeNSevenFive}{11}  % models in the 75 % MCS (of 11)
\providecommand{\rhmcsChMseTrangeSetSevenFive}{}\renewcommand{\rhmcsChMseTrangeSetSevenFive}{the random forest on the HAR benchmark, the linear models on the HAR benchmark, XGBoost on the HAR benchmark, the random forest on the full predictor set, XGBoost on the full predictor set, the elastic net on the full predictor set, ridge regression on the full predictor set, the lasso on the full predictor set, LightGBM on the full predictor set, the closing-interval HAR benchmark, and LightGBM on the HAR benchmark}  % the 75 % MCS, best (last eliminated) first
\providecommand{\rhmcsChMseTrangeStableSevenFive}{}\renewcommand{\rhmcsChMseTrangeStableSevenFive}{yes}  % the 75 % MCS is the same at block lengths (5, 20)
\providecommand{\rhmcsChMseTrangePMaxOutSevenFive}{}\renewcommand{\rhmcsChMseTrangePMaxOutSevenFive}{--}  % largest MCS p-value outside the 75 % set
\providecommand{\rhmcsChMseTrangeNSevenFiveLo}{}\renewcommand{\rhmcsChMseTrangeNSevenFiveLo}{11}  % ... at block length 5
\providecommand{\rhmcsChMseTrangeSetSevenFiveLo}{}\renewcommand{\rhmcsChMseTrangeSetSevenFiveLo}{the random forest on the HAR benchmark, the linear models on the HAR benchmark, XGBoost on the HAR benchmark, XGBoost on the full predictor set, the random forest on the full predictor set, the elastic net on the full predictor set, ridge regression on the full predictor set, LightGBM on the full predictor set, the lasso on the full predictor set, the closing-interval HAR benchmark, and LightGBM on the HAR benchmark}  % ... at block length 5
\providecommand{\rhmcsChMseTrangeNSevenFiveHi}{}\renewcommand{\rhmcsChMseTrangeNSevenFiveHi}{11}  % ... at block length 20
\providecommand{\rhmcsChMseTrangeSetSevenFiveHi}{}\renewcommand{\rhmcsChMseTrangeSetSevenFiveHi}{the random forest on the HAR benchmark, the linear models on the HAR benchmark, the random forest on the full predictor set, XGBoost on the full predictor set, the elastic net on the full predictor set, ridge regression on the full predictor set, the lasso on the full predictor set, XGBoost on the HAR benchmark, the closing-interval HAR benchmark, LightGBM on the full predictor set, and LightGBM on the HAR benchmark}  % ... at block length 20
\providecommand{\rhmcsChMseTrangeBest}{}\renewcommand{\rhmcsChMseTrangeBest}{the random forest on the HAR benchmark}  % the last model left (MCS p-value 1)
\providecommand{\rhmcsBarQAfSpaC}{}\renewcommand{\rhmcsBarQAfSpaC}{\ensuremath{<}0.001}  % SPA p-value (consistent), family BarQAf
\providecommand{\rhmcsBarQAfSpaL}{}\renewcommand{\rhmcsBarQAfSpaL}{\ensuremath{<}0.001}  % ... lower
\providecommand{\rhmcsBarQAfSpaU}{}\renewcommand{\rhmcsBarQAfSpaU}{\ensuremath{<}0.001}  % ... upper
\providecommand{\rhmcsBarQAfNRej}{}\renewcommand{\rhmcsBarQAfNRej}{6}  % StepM rejections at FWER 5 % (of 6)
\providecommand{\rhmcsBarQAfNRejWord}{}\renewcommand{\rhmcsBarQAfNRejWord}{six}  % ... as a word
\providecommand{\rhmcsBarQAfNRaw}{}\renewcommand{\rhmcsBarQAfNRaw}{6}  % single-test bootstrap p < 0.05 (no adjustment)
\providecommand{\rhmcsBarQAfNRawWord}{}\renewcommand{\rhmcsBarQAfNRawWord}{six}  % ... as a word
\providecommand{\rhmcsBarQAfNPos}{}\renewcommand{\rhmcsBarQAfNPos}{6}  % alternatives with a positive estimate (better than the benchmark)
\providecommand{\rhmcsBarQAfPMin}{}\renewcommand{\rhmcsBarQAfPMin}{\ensuremath{<}0.001}  % smallest StepM-adjusted p-value
\providecommand{\rhmcsBarQAfPMax}{}\renewcommand{\rhmcsBarQAfPMax}{\ensuremath{<}0.001}  % largest StepM-adjusted p-value
\providecommand{\rhmcsBarQAfPTwoMin}{}\renewcommand{\rhmcsBarQAfPTwoMin}{\ensuremath{<}0.001}  % smallest two-sided StepM-adjusted p-value
\providecommand{\rhmcsBarQAfPTwoMax}{}\renewcommand{\rhmcsBarQAfPTwoMax}{\ensuremath{<}0.001}  % largest two-sided StepM-adjusted p-value
\providecommand{\rhmcsBarQAfRawMin}{}\renewcommand{\rhmcsBarQAfRawMin}{\ensuremath{<}0.001}  % smallest single-test bootstrap p-value
\providecommand{\rhmcsBarQAfRej}{}\renewcommand{\rhmcsBarQAfRej}{ridge regression on the full predictor set, the lasso on the full predictor set, the elastic net on the full predictor set, LightGBM on the full predictor set, XGBoost on the full predictor set, and the random forest on the full predictor set}  % the alternatives StepM rejects at 5 % (none: none)
\providecommand{\rhmcsBarQAfNRejTwo}{}\renewcommand{\rhmcsBarQAfNRejTwo}{6}  % StepM rejections at FWER 5 %, two-sided
\providecommand{\rhmcsBarQAfNRejTwoWord}{}\renewcommand{\rhmcsBarQAfNRejTwoWord}{six}  % ... as a word
\providecommand{\rhmcsBarQAfNRawTwo}{}\renewcommand{\rhmcsBarQAfNRawTwo}{6}  % single-test bootstrap p < 0.05, two-sided
\providecommand{\rhmcsBarQAfRejTwo}{}\renewcommand{\rhmcsBarQAfRejTwo}{ridge regression on the full predictor set, the lasso on the full predictor set, the elastic net on the full predictor set, LightGBM on the full predictor set, XGBoost on the full predictor set, and the random forest on the full predictor set}  % the alternatives the two-sided StepM rejects at 5 % (none: none)
\providecommand{\rhmcsBarQAfPRidge}{}\renewcommand{\rhmcsBarQAfPRidge}{\ensuremath{<}0.001}  % StepM-adjusted p-value (one-sided), Ridge
\providecommand{\rhmcsBarQAfRawRidge}{}\renewcommand{\rhmcsBarQAfRawRidge}{\ensuremath{<}0.001}  % single-test bootstrap p-value (one-sided), Ridge
\providecommand{\rhmcsBarQAfPTwoRidge}{}\renewcommand{\rhmcsBarQAfPTwoRidge}{\ensuremath{<}0.001}  % StepM-adjusted p-value (two-sided), Ridge
\providecommand{\rhmcsBarQAfRawTwoRidge}{}\renewcommand{\rhmcsBarQAfRawTwoRidge}{\ensuremath{<}0.001}  % single-test bootstrap p-value (two-sided), Ridge
\providecommand{\rhmcsBarQAfPLasso}{}\renewcommand{\rhmcsBarQAfPLasso}{\ensuremath{<}0.001}  % StepM-adjusted p-value (one-sided), Lasso
\providecommand{\rhmcsBarQAfRawLasso}{}\renewcommand{\rhmcsBarQAfRawLasso}{\ensuremath{<}0.001}  % single-test bootstrap p-value (one-sided), Lasso
\providecommand{\rhmcsBarQAfPTwoLasso}{}\renewcommand{\rhmcsBarQAfPTwoLasso}{\ensuremath{<}0.001}  % StepM-adjusted p-value (two-sided), Lasso
\providecommand{\rhmcsBarQAfRawTwoLasso}{}\renewcommand{\rhmcsBarQAfRawTwoLasso}{\ensuremath{<}0.001}  % single-test bootstrap p-value (two-sided), Lasso
\providecommand{\rhmcsBarQAfPEnet}{}\renewcommand{\rhmcsBarQAfPEnet}{\ensuremath{<}0.001}  % StepM-adjusted p-value (one-sided), Elastic net
\providecommand{\rhmcsBarQAfRawEnet}{}\renewcommand{\rhmcsBarQAfRawEnet}{\ensuremath{<}0.001}  % single-test bootstrap p-value (one-sided), Elastic net
\providecommand{\rhmcsBarQAfPTwoEnet}{}\renewcommand{\rhmcsBarQAfPTwoEnet}{\ensuremath{<}0.001}  % StepM-adjusted p-value (two-sided), Elastic net
\providecommand{\rhmcsBarQAfRawTwoEnet}{}\renewcommand{\rhmcsBarQAfRawTwoEnet}{\ensuremath{<}0.001}  % single-test bootstrap p-value (two-sided), Elastic net
\providecommand{\rhmcsBarQAfPLgbm}{}\renewcommand{\rhmcsBarQAfPLgbm}{\ensuremath{<}0.001}  % StepM-adjusted p-value (one-sided), LightGBM
\providecommand{\rhmcsBarQAfRawLgbm}{}\renewcommand{\rhmcsBarQAfRawLgbm}{\ensuremath{<}0.001}  % single-test bootstrap p-value (one-sided), LightGBM
\providecommand{\rhmcsBarQAfPTwoLgbm}{}\renewcommand{\rhmcsBarQAfPTwoLgbm}{\ensuremath{<}0.001}  % StepM-adjusted p-value (two-sided), LightGBM
\providecommand{\rhmcsBarQAfRawTwoLgbm}{}\renewcommand{\rhmcsBarQAfRawTwoLgbm}{\ensuremath{<}0.001}  % single-test bootstrap p-value (two-sided), LightGBM
\providecommand{\rhmcsBarQAfPXgb}{}\renewcommand{\rhmcsBarQAfPXgb}{\ensuremath{<}0.001}  % StepM-adjusted p-value (one-sided), XGBoost
\providecommand{\rhmcsBarQAfRawXgb}{}\renewcommand{\rhmcsBarQAfRawXgb}{\ensuremath{<}0.001}  % single-test bootstrap p-value (one-sided), XGBoost
\providecommand{\rhmcsBarQAfPTwoXgb}{}\renewcommand{\rhmcsBarQAfPTwoXgb}{\ensuremath{<}0.001}  % StepM-adjusted p-value (two-sided), XGBoost
\providecommand{\rhmcsBarQAfRawTwoXgb}{}\renewcommand{\rhmcsBarQAfRawTwoXgb}{\ensuremath{<}0.001}  % single-test bootstrap p-value (two-sided), XGBoost
\providecommand{\rhmcsBarQAfPRf}{}\renewcommand{\rhmcsBarQAfPRf}{\ensuremath{<}0.001}  % StepM-adjusted p-value (one-sided), Random forest
\providecommand{\rhmcsBarQAfRawRf}{}\renewcommand{\rhmcsBarQAfRawRf}{\ensuremath{<}0.001}  % single-test bootstrap p-value (one-sided), Random forest
\providecommand{\rhmcsBarQAfPTwoRf}{}\renewcommand{\rhmcsBarQAfPTwoRf}{\ensuremath{<}0.001}  % StepM-adjusted p-value (two-sided), Random forest
\providecommand{\rhmcsBarQAfRawTwoRf}{}\renewcommand{\rhmcsBarQAfRawTwoRf}{\ensuremath{<}0.001}  % single-test bootstrap p-value (two-sided), Random forest
\providecommand{\rhmcsBarQAfNRejLo}{}\renewcommand{\rhmcsBarQAfNRejLo}{6}  % StepM rejections at 5 %, block length 5
\providecommand{\rhmcsBarQAfSpaCLo}{}\renewcommand{\rhmcsBarQAfSpaCLo}{\ensuremath{<}0.001}  % SPA p-value (consistent), block length 5
\providecommand{\rhmcsBarQAfNRejTwoLo}{}\renewcommand{\rhmcsBarQAfNRejTwoLo}{6}  % ... two-sided, block length 5
\providecommand{\rhmcsBarQAfPMinLo}{}\renewcommand{\rhmcsBarQAfPMinLo}{\ensuremath{<}0.001}  % smallest adjusted p, block length 5
\providecommand{\rhmcsBarQAfPTwoMinLo}{}\renewcommand{\rhmcsBarQAfPTwoMinLo}{\ensuremath{<}0.001}  % ... two-sided, block length 5
\providecommand{\rhmcsBarQAfNRejHi}{}\renewcommand{\rhmcsBarQAfNRejHi}{6}  % StepM rejections at 5 %, block length 20
\providecommand{\rhmcsBarQAfSpaCHi}{}\renewcommand{\rhmcsBarQAfSpaCHi}{\ensuremath{<}0.001}  % SPA p-value (consistent), block length 20
\providecommand{\rhmcsBarQAfNRejTwoHi}{}\renewcommand{\rhmcsBarQAfNRejTwoHi}{6}  % ... two-sided, block length 20
\providecommand{\rhmcsBarQAfPMinHi}{}\renewcommand{\rhmcsBarQAfPMinHi}{\ensuremath{<}0.001}  % smallest adjusted p, block length 20
\providecommand{\rhmcsBarQAfPTwoMinHi}{}\renewcommand{\rhmcsBarQAfPTwoMinHi}{\ensuremath{<}0.001}  % ... two-sided, block length 20
\providecommand{\rhmcsBarMAfSpaC}{}\renewcommand{\rhmcsBarMAfSpaC}{0.64}  % SPA p-value (consistent), family BarMAf
\providecommand{\rhmcsBarMAfSpaL}{}\renewcommand{\rhmcsBarMAfSpaL}{0.22}  % ... lower
\providecommand{\rhmcsBarMAfSpaU}{}\renewcommand{\rhmcsBarMAfSpaU}{0.64}  % ... upper
\providecommand{\rhmcsBarMAfNRej}{}\renewcommand{\rhmcsBarMAfNRej}{0}  % StepM rejections at FWER 5 % (of 6)
\providecommand{\rhmcsBarMAfNRejWord}{}\renewcommand{\rhmcsBarMAfNRejWord}{none}  % ... as a word
\providecommand{\rhmcsBarMAfNRaw}{}\renewcommand{\rhmcsBarMAfNRaw}{0}  % single-test bootstrap p < 0.05 (no adjustment)
\providecommand{\rhmcsBarMAfNRawWord}{}\renewcommand{\rhmcsBarMAfNRawWord}{none}  % ... as a word
\providecommand{\rhmcsBarMAfNPos}{}\renewcommand{\rhmcsBarMAfNPos}{0}  % alternatives with a positive estimate (better than the benchmark)
\providecommand{\rhmcsBarMAfPMin}{}\renewcommand{\rhmcsBarMAfPMin}{0.88}  % smallest StepM-adjusted p-value
\providecommand{\rhmcsBarMAfPMax}{}\renewcommand{\rhmcsBarMAfPMax}{0.89}  % largest StepM-adjusted p-value
\providecommand{\rhmcsBarMAfPTwoMin}{}\renewcommand{\rhmcsBarMAfPTwoMin}{0.28}  % smallest two-sided StepM-adjusted p-value
\providecommand{\rhmcsBarMAfPTwoMax}{}\renewcommand{\rhmcsBarMAfPTwoMax}{0.38}  % largest two-sided StepM-adjusted p-value
\providecommand{\rhmcsBarMAfRawMin}{}\renewcommand{\rhmcsBarMAfRawMin}{0.84}  % smallest single-test bootstrap p-value
\providecommand{\rhmcsBarMAfRej}{}\renewcommand{\rhmcsBarMAfRej}{none}  % the alternatives StepM rejects at 5 % (none: none)
\providecommand{\rhmcsBarMAfNRejTwo}{}\renewcommand{\rhmcsBarMAfNRejTwo}{0}  % StepM rejections at FWER 5 %, two-sided
\providecommand{\rhmcsBarMAfNRejTwoWord}{}\renewcommand{\rhmcsBarMAfNRejTwoWord}{none}  % ... as a word
\providecommand{\rhmcsBarMAfNRawTwo}{}\renewcommand{\rhmcsBarMAfNRawTwo}{0}  % single-test bootstrap p < 0.05, two-sided
\providecommand{\rhmcsBarMAfRejTwo}{}\renewcommand{\rhmcsBarMAfRejTwo}{none}  % the alternatives the two-sided StepM rejects at 5 % (none: none)
\providecommand{\rhmcsBarMAfPRidge}{}\renewcommand{\rhmcsBarMAfPRidge}{0.89}  % StepM-adjusted p-value (one-sided), Ridge
\providecommand{\rhmcsBarMAfRawRidge}{}\renewcommand{\rhmcsBarMAfRawRidge}{0.87}  % single-test bootstrap p-value (one-sided), Ridge
\providecommand{\rhmcsBarMAfPTwoRidge}{}\renewcommand{\rhmcsBarMAfPTwoRidge}{0.32}  % StepM-adjusted p-value (two-sided), Ridge
\providecommand{\rhmcsBarMAfRawTwoRidge}{}\renewcommand{\rhmcsBarMAfRawTwoRidge}{0.23}  % single-test bootstrap p-value (two-sided), Ridge
\providecommand{\rhmcsBarMAfPLasso}{}\renewcommand{\rhmcsBarMAfPLasso}{0.89}  % StepM-adjusted p-value (one-sided), Lasso
\providecommand{\rhmcsBarMAfRawLasso}{}\renewcommand{\rhmcsBarMAfRawLasso}{0.87}  % single-test bootstrap p-value (one-sided), Lasso
\providecommand{\rhmcsBarMAfPTwoLasso}{}\renewcommand{\rhmcsBarMAfPTwoLasso}{0.28}  % StepM-adjusted p-value (two-sided), Lasso
\providecommand{\rhmcsBarMAfRawTwoLasso}{}\renewcommand{\rhmcsBarMAfRawTwoLasso}{0.20}  % single-test bootstrap p-value (two-sided), Lasso
\providecommand{\rhmcsBarMAfPEnet}{}\renewcommand{\rhmcsBarMAfPEnet}{0.89}  % StepM-adjusted p-value (one-sided), Elastic net
\providecommand{\rhmcsBarMAfRawEnet}{}\renewcommand{\rhmcsBarMAfRawEnet}{0.87}  % single-test bootstrap p-value (one-sided), Elastic net
\providecommand{\rhmcsBarMAfPTwoEnet}{}\renewcommand{\rhmcsBarMAfPTwoEnet}{0.32}  % StepM-adjusted p-value (two-sided), Elastic net
\providecommand{\rhmcsBarMAfRawTwoEnet}{}\renewcommand{\rhmcsBarMAfRawTwoEnet}{0.23}  % single-test bootstrap p-value (two-sided), Elastic net
\providecommand{\rhmcsBarMAfPLgbm}{}\renewcommand{\rhmcsBarMAfPLgbm}{0.88}  % StepM-adjusted p-value (one-sided), LightGBM
\providecommand{\rhmcsBarMAfRawLgbm}{}\renewcommand{\rhmcsBarMAfRawLgbm}{0.84}  % single-test bootstrap p-value (one-sided), LightGBM
\providecommand{\rhmcsBarMAfPTwoLgbm}{}\renewcommand{\rhmcsBarMAfPTwoLgbm}{0.38}  % StepM-adjusted p-value (two-sided), LightGBM
\providecommand{\rhmcsBarMAfRawTwoLgbm}{}\renewcommand{\rhmcsBarMAfRawTwoLgbm}{0.36}  % single-test bootstrap p-value (two-sided), LightGBM
\providecommand{\rhmcsBarMAfPXgb}{}\renewcommand{\rhmcsBarMAfPXgb}{0.88}  % StepM-adjusted p-value (one-sided), XGBoost
\providecommand{\rhmcsBarMAfRawXgb}{}\renewcommand{\rhmcsBarMAfRawXgb}{0.84}  % single-test bootstrap p-value (one-sided), XGBoost
\providecommand{\rhmcsBarMAfPTwoXgb}{}\renewcommand{\rhmcsBarMAfPTwoXgb}{0.38}  % StepM-adjusted p-value (two-sided), XGBoost
\providecommand{\rhmcsBarMAfRawTwoXgb}{}\renewcommand{\rhmcsBarMAfRawTwoXgb}{0.35}  % single-test bootstrap p-value (two-sided), XGBoost
\providecommand{\rhmcsBarMAfPRf}{}\renewcommand{\rhmcsBarMAfPRf}{0.89}  % StepM-adjusted p-value (one-sided), Random forest
\providecommand{\rhmcsBarMAfRawRf}{}\renewcommand{\rhmcsBarMAfRawRf}{0.86}  % single-test bootstrap p-value (one-sided), Random forest
\providecommand{\rhmcsBarMAfPTwoRf}{}\renewcommand{\rhmcsBarMAfPTwoRf}{0.35}  % StepM-adjusted p-value (two-sided), Random forest
\providecommand{\rhmcsBarMAfRawTwoRf}{}\renewcommand{\rhmcsBarMAfRawTwoRf}{0.29}  % single-test bootstrap p-value (two-sided), Random forest
\providecommand{\rhmcsBarMAfNRejLo}{}\renewcommand{\rhmcsBarMAfNRejLo}{0}  % StepM rejections at 5 %, block length 5
\providecommand{\rhmcsBarMAfSpaCLo}{}\renewcommand{\rhmcsBarMAfSpaCLo}{0.64}  % SPA p-value (consistent), block length 5
\providecommand{\rhmcsBarMAfNRejTwoLo}{}\renewcommand{\rhmcsBarMAfNRejTwoLo}{0}  % ... two-sided, block length 5
\providecommand{\rhmcsBarMAfPMinLo}{}\renewcommand{\rhmcsBarMAfPMinLo}{0.91}  % smallest adjusted p, block length 5
\providecommand{\rhmcsBarMAfPTwoMinLo}{}\renewcommand{\rhmcsBarMAfPTwoMinLo}{0.21}  % ... two-sided, block length 5
\providecommand{\rhmcsBarMAfNRejHi}{}\renewcommand{\rhmcsBarMAfNRejHi}{0}  % StepM rejections at 5 %, block length 20
\providecommand{\rhmcsBarMAfSpaCHi}{}\renewcommand{\rhmcsBarMAfSpaCHi}{0.63}  % SPA p-value (consistent), block length 20
\providecommand{\rhmcsBarMAfNRejTwoHi}{}\renewcommand{\rhmcsBarMAfNRejTwoHi}{0}  % ... two-sided, block length 20
\providecommand{\rhmcsBarMAfPMinHi}{}\renewcommand{\rhmcsBarMAfPMinHi}{0.84}  % smallest adjusted p, block length 20
\providecommand{\rhmcsBarMAfPTwoMinHi}{}\renewcommand{\rhmcsBarMAfPTwoMinHi}{0.39}  % ... two-sided, block length 20
\providecommand{\rhmcsAllQAfSpaC}{}\renewcommand{\rhmcsAllQAfSpaC}{0.045}  % SPA p-value (consistent), family AllQAf
\providecommand{\rhmcsAllQAfSpaL}{}\renewcommand{\rhmcsAllQAfSpaL}{0.037}  % ... lower
\providecommand{\rhmcsAllQAfSpaU}{}\renewcommand{\rhmcsAllQAfSpaU}{0.045}  % ... upper
\providecommand{\rhmcsAllQAfNRej}{}\renewcommand{\rhmcsAllQAfNRej}{1}  % StepM rejections at FWER 5 % (of 6)
\providecommand{\rhmcsAllQAfNRejWord}{}\renewcommand{\rhmcsAllQAfNRejWord}{one}  % ... as a word
\providecommand{\rhmcsAllQAfNRaw}{}\renewcommand{\rhmcsAllQAfNRaw}{2}  % single-test bootstrap p < 0.05 (no adjustment)
\providecommand{\rhmcsAllQAfNRawWord}{}\renewcommand{\rhmcsAllQAfNRawWord}{two}  % ... as a word
\providecommand{\rhmcsAllQAfNPos}{}\renewcommand{\rhmcsAllQAfNPos}{3}  % alternatives with a positive estimate (better than the benchmark)
\providecommand{\rhmcsAllQAfPMin}{}\renewcommand{\rhmcsAllQAfPMin}{0.045}  % smallest StepM-adjusted p-value
\providecommand{\rhmcsAllQAfPMax}{}\renewcommand{\rhmcsAllQAfPMax}{0.90}  % largest StepM-adjusted p-value
\providecommand{\rhmcsAllQAfPTwoMin}{}\renewcommand{\rhmcsAllQAfPTwoMin}{0.11}  % smallest two-sided StepM-adjusted p-value
\providecommand{\rhmcsAllQAfPTwoMax}{}\renewcommand{\rhmcsAllQAfPTwoMax}{0.81}  % largest two-sided StepM-adjusted p-value
\providecommand{\rhmcsAllQAfRawMin}{}\renewcommand{\rhmcsAllQAfRawMin}{0.020}  % smallest single-test bootstrap p-value
\providecommand{\rhmcsAllQAfRej}{}\renewcommand{\rhmcsAllQAfRej}{XGBoost on the full predictor set}  % the alternatives StepM rejects at 5 % (none: none)
\providecommand{\rhmcsAllQAfNRejTwo}{}\renewcommand{\rhmcsAllQAfNRejTwo}{0}  % StepM rejections at FWER 5 %, two-sided
\providecommand{\rhmcsAllQAfNRejTwoWord}{}\renewcommand{\rhmcsAllQAfNRejTwoWord}{none}  % ... as a word
\providecommand{\rhmcsAllQAfNRawTwo}{}\renewcommand{\rhmcsAllQAfNRawTwo}{0}  % single-test bootstrap p < 0.05, two-sided
\providecommand{\rhmcsAllQAfRejTwo}{}\renewcommand{\rhmcsAllQAfRejTwo}{none}  % the alternatives the two-sided StepM rejects at 5 % (none: none)
\providecommand{\rhmcsAllQAfPRidge}{}\renewcommand{\rhmcsAllQAfPRidge}{0.78}  % StepM-adjusted p-value (one-sided), Ridge
\providecommand{\rhmcsAllQAfRawRidge}{}\renewcommand{\rhmcsAllQAfRawRidge}{0.76}  % single-test bootstrap p-value (one-sided), Ridge
\providecommand{\rhmcsAllQAfPTwoRidge}{}\renewcommand{\rhmcsAllQAfPTwoRidge}{0.55}  % StepM-adjusted p-value (two-sided), Ridge
\providecommand{\rhmcsAllQAfRawTwoRidge}{}\renewcommand{\rhmcsAllQAfRawTwoRidge}{0.50}  % single-test bootstrap p-value (two-sided), Ridge
\providecommand{\rhmcsAllQAfPLasso}{}\renewcommand{\rhmcsAllQAfPLasso}{0.90}  % StepM-adjusted p-value (one-sided), Lasso
\providecommand{\rhmcsAllQAfRawLasso}{}\renewcommand{\rhmcsAllQAfRawLasso}{0.90}  % single-test bootstrap p-value (one-sided), Lasso
\providecommand{\rhmcsAllQAfPTwoLasso}{}\renewcommand{\rhmcsAllQAfPTwoLasso}{0.32}  % StepM-adjusted p-value (two-sided), Lasso
\providecommand{\rhmcsAllQAfRawTwoLasso}{}\renewcommand{\rhmcsAllQAfRawTwoLasso}{0.18}  % single-test bootstrap p-value (two-sided), Lasso
\providecommand{\rhmcsAllQAfPEnet}{}\renewcommand{\rhmcsAllQAfPEnet}{0.67}  % StepM-adjusted p-value (one-sided), Elastic net
\providecommand{\rhmcsAllQAfRawEnet}{}\renewcommand{\rhmcsAllQAfRawEnet}{0.62}  % single-test bootstrap p-value (one-sided), Elastic net
\providecommand{\rhmcsAllQAfPTwoEnet}{}\renewcommand{\rhmcsAllQAfPTwoEnet}{0.81}  % StepM-adjusted p-value (two-sided), Elastic net
\providecommand{\rhmcsAllQAfRawTwoEnet}{}\renewcommand{\rhmcsAllQAfRawTwoEnet}{0.81}  % single-test bootstrap p-value (two-sided), Elastic net
\providecommand{\rhmcsAllQAfPLgbm}{}\renewcommand{\rhmcsAllQAfPLgbm}{0.081}  % StepM-adjusted p-value (one-sided), LightGBM
\providecommand{\rhmcsAllQAfRawLgbm}{}\renewcommand{\rhmcsAllQAfRawLgbm}{0.037}  % single-test bootstrap p-value (one-sided), LightGBM
\providecommand{\rhmcsAllQAfPTwoLgbm}{}\renewcommand{\rhmcsAllQAfPTwoLgbm}{0.18}  % StepM-adjusted p-value (two-sided), LightGBM
\providecommand{\rhmcsAllQAfRawTwoLgbm}{}\renewcommand{\rhmcsAllQAfRawTwoLgbm}{0.081}  % single-test bootstrap p-value (two-sided), LightGBM
\providecommand{\rhmcsAllQAfPXgb}{}\renewcommand{\rhmcsAllQAfPXgb}{0.045}  % StepM-adjusted p-value (one-sided), XGBoost
\providecommand{\rhmcsAllQAfRawXgb}{}\renewcommand{\rhmcsAllQAfRawXgb}{0.020}  % single-test bootstrap p-value (one-sided), XGBoost
\providecommand{\rhmcsAllQAfPTwoXgb}{}\renewcommand{\rhmcsAllQAfPTwoXgb}{0.11}  % StepM-adjusted p-value (two-sided), XGBoost
\providecommand{\rhmcsAllQAfRawTwoXgb}{}\renewcommand{\rhmcsAllQAfRawTwoXgb}{0.050}  % single-test bootstrap p-value (two-sided), XGBoost
\providecommand{\rhmcsAllQAfPRf}{}\renewcommand{\rhmcsAllQAfPRf}{0.26}  % StepM-adjusted p-value (one-sided), Random forest
\providecommand{\rhmcsAllQAfRawRf}{}\renewcommand{\rhmcsAllQAfRawRf}{0.15}  % single-test bootstrap p-value (one-sided), Random forest
\providecommand{\rhmcsAllQAfPTwoRf}{}\renewcommand{\rhmcsAllQAfPTwoRf}{0.48}  % StepM-adjusted p-value (two-sided), Random forest
\providecommand{\rhmcsAllQAfRawTwoRf}{}\renewcommand{\rhmcsAllQAfRawTwoRf}{0.29}  % single-test bootstrap p-value (two-sided), Random forest
\providecommand{\rhmcsAllQAfNRejLo}{}\renewcommand{\rhmcsAllQAfNRejLo}{1}  % StepM rejections at 5 %, block length 5
\providecommand{\rhmcsAllQAfSpaCLo}{}\renewcommand{\rhmcsAllQAfSpaCLo}{0.036}  % SPA p-value (consistent), block length 5
\providecommand{\rhmcsAllQAfNRejTwoLo}{}\renewcommand{\rhmcsAllQAfNRejTwoLo}{0}  % ... two-sided, block length 5
\providecommand{\rhmcsAllQAfPMinLo}{}\renewcommand{\rhmcsAllQAfPMinLo}{0.036}  % smallest adjusted p, block length 5
\providecommand{\rhmcsAllQAfPTwoMinLo}{}\renewcommand{\rhmcsAllQAfPTwoMinLo}{0.091}  % ... two-sided, block length 5
\providecommand{\rhmcsAllQAfNRejHi}{}\renewcommand{\rhmcsAllQAfNRejHi}{0}  % StepM rejections at 5 %, block length 20
\providecommand{\rhmcsAllQAfSpaCHi}{}\renewcommand{\rhmcsAllQAfSpaCHi}{0.051}  % SPA p-value (consistent), block length 20
\providecommand{\rhmcsAllQAfNRejTwoHi}{}\renewcommand{\rhmcsAllQAfNRejTwoHi}{0}  % ... two-sided, block length 20
\providecommand{\rhmcsAllQAfPMinHi}{}\renewcommand{\rhmcsAllQAfPMinHi}{0.051}  % smallest adjusted p, block length 20
\providecommand{\rhmcsAllQAfPTwoMinHi}{}\renewcommand{\rhmcsAllQAfPTwoMinHi}{0.12}  % ... two-sided, block length 20
\providecommand{\rhmcsAllMAfSpaC}{}\renewcommand{\rhmcsAllMAfSpaC}{0.66}  % SPA p-value (consistent), family AllMAf
\providecommand{\rhmcsAllMAfSpaL}{}\renewcommand{\rhmcsAllMAfSpaL}{0.20}  % ... lower
\providecommand{\rhmcsAllMAfSpaU}{}\renewcommand{\rhmcsAllMAfSpaU}{0.66}  % ... upper
\providecommand{\rhmcsAllMAfNRej}{}\renewcommand{\rhmcsAllMAfNRej}{0}  % StepM rejections at FWER 5 % (of 6)
\providecommand{\rhmcsAllMAfNRejWord}{}\renewcommand{\rhmcsAllMAfNRejWord}{none}  % ... as a word
\providecommand{\rhmcsAllMAfNRaw}{}\renewcommand{\rhmcsAllMAfNRaw}{0}  % single-test bootstrap p < 0.05 (no adjustment)
\providecommand{\rhmcsAllMAfNRawWord}{}\renewcommand{\rhmcsAllMAfNRawWord}{none}  % ... as a word
\providecommand{\rhmcsAllMAfNPos}{}\renewcommand{\rhmcsAllMAfNPos}{0}  % alternatives with a positive estimate (better than the benchmark)
\providecommand{\rhmcsAllMAfPMin}{}\renewcommand{\rhmcsAllMAfPMin}{0.91}  % smallest StepM-adjusted p-value
\providecommand{\rhmcsAllMAfPMax}{}\renewcommand{\rhmcsAllMAfPMax}{0.91}  % largest StepM-adjusted p-value
\providecommand{\rhmcsAllMAfPTwoMin}{}\renewcommand{\rhmcsAllMAfPTwoMin}{0.36}  % smallest two-sided StepM-adjusted p-value
\providecommand{\rhmcsAllMAfPTwoMax}{}\renewcommand{\rhmcsAllMAfPTwoMax}{0.37}  % largest two-sided StepM-adjusted p-value
\providecommand{\rhmcsAllMAfRawMin}{}\renewcommand{\rhmcsAllMAfRawMin}{0.86}  % smallest single-test bootstrap p-value
\providecommand{\rhmcsAllMAfRej}{}\renewcommand{\rhmcsAllMAfRej}{none}  % the alternatives StepM rejects at 5 % (none: none)
\providecommand{\rhmcsAllMAfNRejTwo}{}\renewcommand{\rhmcsAllMAfNRejTwo}{0}  % StepM rejections at FWER 5 %, two-sided
\providecommand{\rhmcsAllMAfNRejTwoWord}{}\renewcommand{\rhmcsAllMAfNRejTwoWord}{none}  % ... as a word
\providecommand{\rhmcsAllMAfNRawTwo}{}\renewcommand{\rhmcsAllMAfNRawTwo}{0}  % single-test bootstrap p < 0.05, two-sided
\providecommand{\rhmcsAllMAfRejTwo}{}\renewcommand{\rhmcsAllMAfRejTwo}{none}  % the alternatives the two-sided StepM rejects at 5 % (none: none)
\providecommand{\rhmcsAllMAfPRidge}{}\renewcommand{\rhmcsAllMAfPRidge}{0.91}  % StepM-adjusted p-value (one-sided), Ridge
\providecommand{\rhmcsAllMAfRawRidge}{}\renewcommand{\rhmcsAllMAfRawRidge}{0.87}  % single-test bootstrap p-value (one-sided), Ridge
\providecommand{\rhmcsAllMAfPTwoRidge}{}\renewcommand{\rhmcsAllMAfPTwoRidge}{0.36}  % StepM-adjusted p-value (two-sided), Ridge
\providecommand{\rhmcsAllMAfRawTwoRidge}{}\renewcommand{\rhmcsAllMAfRawTwoRidge}{0.26}  % single-test bootstrap p-value (two-sided), Ridge
\providecommand{\rhmcsAllMAfPLasso}{}\renewcommand{\rhmcsAllMAfPLasso}{0.91}  % StepM-adjusted p-value (one-sided), Lasso
\providecommand{\rhmcsAllMAfRawLasso}{}\renewcommand{\rhmcsAllMAfRawLasso}{0.87}  % single-test bootstrap p-value (one-sided), Lasso
\providecommand{\rhmcsAllMAfPTwoLasso}{}\renewcommand{\rhmcsAllMAfPTwoLasso}{0.36}  % StepM-adjusted p-value (two-sided), Lasso
\providecommand{\rhmcsAllMAfRawTwoLasso}{}\renewcommand{\rhmcsAllMAfRawTwoLasso}{0.26}  % single-test bootstrap p-value (two-sided), Lasso
\providecommand{\rhmcsAllMAfPEnet}{}\renewcommand{\rhmcsAllMAfPEnet}{0.91}  % StepM-adjusted p-value (one-sided), Elastic net
\providecommand{\rhmcsAllMAfRawEnet}{}\renewcommand{\rhmcsAllMAfRawEnet}{0.87}  % single-test bootstrap p-value (one-sided), Elastic net
\providecommand{\rhmcsAllMAfPTwoEnet}{}\renewcommand{\rhmcsAllMAfPTwoEnet}{0.37}  % StepM-adjusted p-value (two-sided), Elastic net
\providecommand{\rhmcsAllMAfRawTwoEnet}{}\renewcommand{\rhmcsAllMAfRawTwoEnet}{0.26}  % single-test bootstrap p-value (two-sided), Elastic net
\providecommand{\rhmcsAllMAfPLgbm}{}\renewcommand{\rhmcsAllMAfPLgbm}{0.91}  % StepM-adjusted p-value (one-sided), LightGBM
\providecommand{\rhmcsAllMAfRawLgbm}{}\renewcommand{\rhmcsAllMAfRawLgbm}{0.86}  % single-test bootstrap p-value (one-sided), LightGBM
\providecommand{\rhmcsAllMAfPTwoLgbm}{}\renewcommand{\rhmcsAllMAfPTwoLgbm}{0.37}  % StepM-adjusted p-value (two-sided), LightGBM
\providecommand{\rhmcsAllMAfRawTwoLgbm}{}\renewcommand{\rhmcsAllMAfRawTwoLgbm}{0.33}  % single-test bootstrap p-value (two-sided), LightGBM
\providecommand{\rhmcsAllMAfPXgb}{}\renewcommand{\rhmcsAllMAfPXgb}{0.91}  % StepM-adjusted p-value (one-sided), XGBoost
\providecommand{\rhmcsAllMAfRawXgb}{}\renewcommand{\rhmcsAllMAfRawXgb}{0.86}  % single-test bootstrap p-value (one-sided), XGBoost
\providecommand{\rhmcsAllMAfPTwoXgb}{}\renewcommand{\rhmcsAllMAfPTwoXgb}{0.37}  % StepM-adjusted p-value (two-sided), XGBoost
\providecommand{\rhmcsAllMAfRawTwoXgb}{}\renewcommand{\rhmcsAllMAfRawTwoXgb}{0.32}  % single-test bootstrap p-value (two-sided), XGBoost
\providecommand{\rhmcsAllMAfPRf}{}\renewcommand{\rhmcsAllMAfPRf}{0.91}  % StepM-adjusted p-value (one-sided), Random forest
\providecommand{\rhmcsAllMAfRawRf}{}\renewcommand{\rhmcsAllMAfRawRf}{0.86}  % single-test bootstrap p-value (one-sided), Random forest
\providecommand{\rhmcsAllMAfPTwoRf}{}\renewcommand{\rhmcsAllMAfPTwoRf}{0.37}  % StepM-adjusted p-value (two-sided), Random forest
\providecommand{\rhmcsAllMAfRawTwoRf}{}\renewcommand{\rhmcsAllMAfRawTwoRf}{0.30}  % single-test bootstrap p-value (two-sided), Random forest
\providecommand{\rhmcsAllMAfNRejLo}{}\renewcommand{\rhmcsAllMAfNRejLo}{0}  % StepM rejections at 5 %, block length 5
\providecommand{\rhmcsAllMAfSpaCLo}{}\renewcommand{\rhmcsAllMAfSpaCLo}{0.73}  % SPA p-value (consistent), block length 5
\providecommand{\rhmcsAllMAfNRejTwoLo}{}\renewcommand{\rhmcsAllMAfNRejTwoLo}{0}  % ... two-sided, block length 5
\providecommand{\rhmcsAllMAfPMinLo}{}\renewcommand{\rhmcsAllMAfPMinLo}{0.97}  % smallest adjusted p, block length 5
\providecommand{\rhmcsAllMAfPTwoMinLo}{}\renewcommand{\rhmcsAllMAfPTwoMinLo}{0.24}  % ... two-sided, block length 5
\providecommand{\rhmcsAllMAfNRejHi}{}\renewcommand{\rhmcsAllMAfNRejHi}{0}  % StepM rejections at 5 %, block length 20
\providecommand{\rhmcsAllMAfSpaCHi}{}\renewcommand{\rhmcsAllMAfSpaCHi}{0.62}  % SPA p-value (consistent), block length 20
\providecommand{\rhmcsAllMAfNRejTwoHi}{}\renewcommand{\rhmcsAllMAfNRejTwoHi}{0}  % ... two-sided, block length 20
\providecommand{\rhmcsAllMAfPMinHi}{}\renewcommand{\rhmcsAllMAfPMinHi}{0.87}  % smallest adjusted p, block length 20
\providecommand{\rhmcsAllMAfPTwoMinHi}{}\renewcommand{\rhmcsAllMAfPTwoMinHi}{0.45}  % ... two-sided, block length 20
\providecommand{\rhmcsOosQAfSpaC}{}\renewcommand{\rhmcsOosQAfSpaC}{0.34}  % SPA p-value (consistent), family OosQAf
\providecommand{\rhmcsOosQAfSpaL}{}\renewcommand{\rhmcsOosQAfSpaL}{0.34}  % ... lower
\providecommand{\rhmcsOosQAfSpaU}{}\renewcommand{\rhmcsOosQAfSpaU}{0.48}  % ... upper
\providecommand{\rhmcsOosQAfNRej}{}\renewcommand{\rhmcsOosQAfNRej}{0}  % StepM rejections at FWER 5 % (of 6)
\providecommand{\rhmcsOosQAfNRejWord}{}\renewcommand{\rhmcsOosQAfNRejWord}{none}  % ... as a word
\providecommand{\rhmcsOosQAfNRaw}{}\renewcommand{\rhmcsOosQAfNRaw}{0}  % single-test bootstrap p < 0.05 (no adjustment)
\providecommand{\rhmcsOosQAfNRawWord}{}\renewcommand{\rhmcsOosQAfNRawWord}{none}  % ... as a word
\providecommand{\rhmcsOosQAfNPos}{}\renewcommand{\rhmcsOosQAfNPos}{3}  % alternatives with a positive estimate (better than the benchmark)
\providecommand{\rhmcsOosQAfPMin}{}\renewcommand{\rhmcsOosQAfPMin}{0.48}  % smallest StepM-adjusted p-value
\providecommand{\rhmcsOosQAfPMax}{}\renewcommand{\rhmcsOosQAfPMax}{0.99}  % largest StepM-adjusted p-value
\providecommand{\rhmcsOosQAfPTwoMin}{}\renewcommand{\rhmcsOosQAfPTwoMin}{0.032}  % smallest two-sided StepM-adjusted p-value
\providecommand{\rhmcsOosQAfPTwoMax}{}\renewcommand{\rhmcsOosQAfPTwoMax}{0.78}  % largest two-sided StepM-adjusted p-value
\providecommand{\rhmcsOosQAfRawMin}{}\renewcommand{\rhmcsOosQAfRawMin}{0.26}  % smallest single-test bootstrap p-value
\providecommand{\rhmcsOosQAfRej}{}\renewcommand{\rhmcsOosQAfRej}{none}  % the alternatives StepM rejects at 5 % (none: none)
\providecommand{\rhmcsOosQAfNRejTwo}{}\renewcommand{\rhmcsOosQAfNRejTwo}{1}  % StepM rejections at FWER 5 %, two-sided
\providecommand{\rhmcsOosQAfNRejTwoWord}{}\renewcommand{\rhmcsOosQAfNRejTwoWord}{one}  % ... as a word
\providecommand{\rhmcsOosQAfNRawTwo}{}\renewcommand{\rhmcsOosQAfNRawTwo}{3}  % single-test bootstrap p < 0.05, two-sided
\providecommand{\rhmcsOosQAfRejTwo}{}\renewcommand{\rhmcsOosQAfRejTwo}{the lasso on the full predictor set}  % the alternatives the two-sided StepM rejects at 5 % (none: none)
\providecommand{\rhmcsOosQAfPRidge}{}\renewcommand{\rhmcsOosQAfPRidge}{0.99}  % StepM-adjusted p-value (one-sided), Ridge
\providecommand{\rhmcsOosQAfRawRidge}{}\renewcommand{\rhmcsOosQAfRawRidge}{0.99}  % single-test bootstrap p-value (one-sided), Ridge
\providecommand{\rhmcsOosQAfPTwoRidge}{}\renewcommand{\rhmcsOosQAfPTwoRidge}{0.052}  % StepM-adjusted p-value (two-sided), Ridge
\providecommand{\rhmcsOosQAfRawTwoRidge}{}\renewcommand{\rhmcsOosQAfRawTwoRidge}{0.019}  % single-test bootstrap p-value (two-sided), Ridge
\providecommand{\rhmcsOosQAfPLasso}{}\renewcommand{\rhmcsOosQAfPLasso}{0.99}  % StepM-adjusted p-value (one-sided), Lasso
\providecommand{\rhmcsOosQAfRawLasso}{}\renewcommand{\rhmcsOosQAfRawLasso}{0.99}  % single-test bootstrap p-value (one-sided), Lasso
\providecommand{\rhmcsOosQAfPTwoLasso}{}\renewcommand{\rhmcsOosQAfPTwoLasso}{0.032}  % StepM-adjusted p-value (two-sided), Lasso
\providecommand{\rhmcsOosQAfRawTwoLasso}{}\renewcommand{\rhmcsOosQAfRawTwoLasso}{0.012}  % single-test bootstrap p-value (two-sided), Lasso
\providecommand{\rhmcsOosQAfPEnet}{}\renewcommand{\rhmcsOosQAfPEnet}{0.98}  % StepM-adjusted p-value (one-sided), Elastic net
\providecommand{\rhmcsOosQAfRawEnet}{}\renewcommand{\rhmcsOosQAfRawEnet}{0.97}  % single-test bootstrap p-value (one-sided), Elastic net
\providecommand{\rhmcsOosQAfPTwoEnet}{}\renewcommand{\rhmcsOosQAfPTwoEnet}{0.11}  % StepM-adjusted p-value (two-sided), Elastic net
\providecommand{\rhmcsOosQAfRawTwoEnet}{}\renewcommand{\rhmcsOosQAfRawTwoEnet}{0.049}  % single-test bootstrap p-value (two-sided), Elastic net
\providecommand{\rhmcsOosQAfPLgbm}{}\renewcommand{\rhmcsOosQAfPLgbm}{0.53}  % StepM-adjusted p-value (one-sided), LightGBM
\providecommand{\rhmcsOosQAfRawLgbm}{}\renewcommand{\rhmcsOosQAfRawLgbm}{0.31}  % single-test bootstrap p-value (one-sided), LightGBM
\providecommand{\rhmcsOosQAfPTwoLgbm}{}\renewcommand{\rhmcsOosQAfPTwoLgbm}{0.76}  % StepM-adjusted p-value (two-sided), LightGBM
\providecommand{\rhmcsOosQAfRawTwoLgbm}{}\renewcommand{\rhmcsOosQAfRawTwoLgbm}{0.62}  % single-test bootstrap p-value (two-sided), LightGBM
\providecommand{\rhmcsOosQAfPXgb}{}\renewcommand{\rhmcsOosQAfPXgb}{0.48}  % StepM-adjusted p-value (one-sided), XGBoost
\providecommand{\rhmcsOosQAfRawXgb}{}\renewcommand{\rhmcsOosQAfRawXgb}{0.26}  % single-test bootstrap p-value (one-sided), XGBoost
\providecommand{\rhmcsOosQAfPTwoXgb}{}\renewcommand{\rhmcsOosQAfPTwoXgb}{0.68}  % StepM-adjusted p-value (two-sided), XGBoost
\providecommand{\rhmcsOosQAfRawTwoXgb}{}\renewcommand{\rhmcsOosQAfRawTwoXgb}{0.53}  % single-test bootstrap p-value (two-sided), XGBoost
\providecommand{\rhmcsOosQAfPRf}{}\renewcommand{\rhmcsOosQAfPRf}{0.59}  % StepM-adjusted p-value (one-sided), Random forest
\providecommand{\rhmcsOosQAfRawRf}{}\renewcommand{\rhmcsOosQAfRawRf}{0.40}  % single-test bootstrap p-value (one-sided), Random forest
\providecommand{\rhmcsOosQAfPTwoRf}{}\renewcommand{\rhmcsOosQAfPTwoRf}{0.78}  % StepM-adjusted p-value (two-sided), Random forest
\providecommand{\rhmcsOosQAfRawTwoRf}{}\renewcommand{\rhmcsOosQAfRawTwoRf}{0.78}  % single-test bootstrap p-value (two-sided), Random forest
\providecommand{\rhmcsOosQAfNRejLo}{}\renewcommand{\rhmcsOosQAfNRejLo}{0}  % StepM rejections at 5 %, block length 5
\providecommand{\rhmcsOosQAfSpaCLo}{}\renewcommand{\rhmcsOosQAfSpaCLo}{0.34}  % SPA p-value (consistent), block length 5
\providecommand{\rhmcsOosQAfNRejTwoLo}{}\renewcommand{\rhmcsOosQAfNRejTwoLo}{2}  % ... two-sided, block length 5
\providecommand{\rhmcsOosQAfPMinLo}{}\renewcommand{\rhmcsOosQAfPMinLo}{0.47}  % smallest adjusted p, block length 5
\providecommand{\rhmcsOosQAfPTwoMinLo}{}\renewcommand{\rhmcsOosQAfPTwoMinLo}{0.017}  % ... two-sided, block length 5
\providecommand{\rhmcsOosQAfNRejHi}{}\renewcommand{\rhmcsOosQAfNRejHi}{0}  % StepM rejections at 5 %, block length 20
\providecommand{\rhmcsOosQAfSpaCHi}{}\renewcommand{\rhmcsOosQAfSpaCHi}{0.43}  % SPA p-value (consistent), block length 20
\providecommand{\rhmcsOosQAfNRejTwoHi}{}\renewcommand{\rhmcsOosQAfNRejTwoHi}{0}  % ... two-sided, block length 20
\providecommand{\rhmcsOosQAfPMinHi}{}\renewcommand{\rhmcsOosQAfPMinHi}{0.46}  % smallest adjusted p, block length 20
\providecommand{\rhmcsOosQAfPTwoMinHi}{}\renewcommand{\rhmcsOosQAfPTwoMinHi}{0.051}  % ... two-sided, block length 20
\providecommand{\rhmcsOosQBmSpaC}{}\renewcommand{\rhmcsOosQBmSpaC}{0.53}  % SPA p-value (consistent), family OosQBm
\providecommand{\rhmcsOosQBmSpaL}{}\renewcommand{\rhmcsOosQBmSpaL}{0.083}  % ... lower
\providecommand{\rhmcsOosQBmSpaU}{}\renewcommand{\rhmcsOosQBmSpaU}{0.69}  % ... upper
\providecommand{\rhmcsOosQBmNRej}{}\renewcommand{\rhmcsOosQBmNRej}{0}  % StepM rejections at FWER 5 % (of 6)
\providecommand{\rhmcsOosQBmNRejWord}{}\renewcommand{\rhmcsOosQBmNRejWord}{none}  % ... as a word
\providecommand{\rhmcsOosQBmNRaw}{}\renewcommand{\rhmcsOosQBmNRaw}{0}  % single-test bootstrap p < 0.05 (no adjustment)
\providecommand{\rhmcsOosQBmNRawWord}{}\renewcommand{\rhmcsOosQBmNRawWord}{none}  % ... as a word
\providecommand{\rhmcsOosQBmNPos}{}\renewcommand{\rhmcsOosQBmNPos}{0}  % alternatives with a positive estimate (better than the benchmark)
\providecommand{\rhmcsOosQBmPMin}{}\renewcommand{\rhmcsOosQBmPMin}{0.98}  % smallest StepM-adjusted p-value
\providecommand{\rhmcsOosQBmPMax}{}\renewcommand{\rhmcsOosQBmPMax}{1.00}  % largest StepM-adjusted p-value
\providecommand{\rhmcsOosQBmPTwoMin}{}\renewcommand{\rhmcsOosQBmPTwoMin}{0.001}  % smallest two-sided StepM-adjusted p-value
\providecommand{\rhmcsOosQBmPTwoMax}{}\renewcommand{\rhmcsOosQBmPTwoMax}{0.17}  % largest two-sided StepM-adjusted p-value
\providecommand{\rhmcsOosQBmRawMin}{}\renewcommand{\rhmcsOosQBmRawMin}{0.92}  % smallest single-test bootstrap p-value
\providecommand{\rhmcsOosQBmRej}{}\renewcommand{\rhmcsOosQBmRej}{none}  % the alternatives StepM rejects at 5 % (none: none)
\providecommand{\rhmcsOosQBmNRejTwo}{}\renewcommand{\rhmcsOosQBmNRejTwo}{4}  % StepM rejections at FWER 5 %, two-sided
\providecommand{\rhmcsOosQBmNRejTwoWord}{}\renewcommand{\rhmcsOosQBmNRejTwoWord}{four}  % ... as a word
\providecommand{\rhmcsOosQBmNRawTwo}{}\renewcommand{\rhmcsOosQBmNRawTwo}{4}  % single-test bootstrap p < 0.05, two-sided
\providecommand{\rhmcsOosQBmRejTwo}{}\renewcommand{\rhmcsOosQBmRejTwo}{ridge regression on the full predictor set, the lasso on the full predictor set, the elastic net on the full predictor set, and the random forest on the full predictor set}  % the alternatives the two-sided StepM rejects at 5 % (none: none)
\providecommand{\rhmcsOosQBmPRidge}{}\renewcommand{\rhmcsOosQBmPRidge}{1.00}  % StepM-adjusted p-value (one-sided), Ridge
\providecommand{\rhmcsOosQBmRawRidge}{}\renewcommand{\rhmcsOosQBmRawRidge}{1.00}  % single-test bootstrap p-value (one-sided), Ridge
\providecommand{\rhmcsOosQBmPTwoRidge}{}\renewcommand{\rhmcsOosQBmPTwoRidge}{0.001}  % StepM-adjusted p-value (two-sided), Ridge
\providecommand{\rhmcsOosQBmRawTwoRidge}{}\renewcommand{\rhmcsOosQBmRawTwoRidge}{\ensuremath{<}0.001}  % single-test bootstrap p-value (two-sided), Ridge
\providecommand{\rhmcsOosQBmPLasso}{}\renewcommand{\rhmcsOosQBmPLasso}{1.00}  % StepM-adjusted p-value (one-sided), Lasso
\providecommand{\rhmcsOosQBmRawLasso}{}\renewcommand{\rhmcsOosQBmRawLasso}{1.00}  % single-test bootstrap p-value (one-sided), Lasso
\providecommand{\rhmcsOosQBmPTwoLasso}{}\renewcommand{\rhmcsOosQBmPTwoLasso}{0.001}  % StepM-adjusted p-value (two-sided), Lasso
\providecommand{\rhmcsOosQBmRawTwoLasso}{}\renewcommand{\rhmcsOosQBmRawTwoLasso}{\ensuremath{<}0.001}  % single-test bootstrap p-value (two-sided), Lasso
\providecommand{\rhmcsOosQBmPEnet}{}\renewcommand{\rhmcsOosQBmPEnet}{1.00}  % StepM-adjusted p-value (one-sided), Elastic net
\providecommand{\rhmcsOosQBmRawEnet}{}\renewcommand{\rhmcsOosQBmRawEnet}{1.00}  % single-test bootstrap p-value (one-sided), Elastic net
\providecommand{\rhmcsOosQBmPTwoEnet}{}\renewcommand{\rhmcsOosQBmPTwoEnet}{0.002}  % StepM-adjusted p-value (two-sided), Elastic net
\providecommand{\rhmcsOosQBmRawTwoEnet}{}\renewcommand{\rhmcsOosQBmRawTwoEnet}{\ensuremath{<}0.001}  % single-test bootstrap p-value (two-sided), Elastic net
\providecommand{\rhmcsOosQBmPLgbm}{}\renewcommand{\rhmcsOosQBmPLgbm}{0.98}  % StepM-adjusted p-value (one-sided), LightGBM
\providecommand{\rhmcsOosQBmRawLgbm}{}\renewcommand{\rhmcsOosQBmRawLgbm}{0.92}  % single-test bootstrap p-value (one-sided), LightGBM
\providecommand{\rhmcsOosQBmPTwoLgbm}{}\renewcommand{\rhmcsOosQBmPTwoLgbm}{0.17}  % StepM-adjusted p-value (two-sided), LightGBM
\providecommand{\rhmcsOosQBmRawTwoLgbm}{}\renewcommand{\rhmcsOosQBmRawTwoLgbm}{0.15}  % single-test bootstrap p-value (two-sided), LightGBM
\providecommand{\rhmcsOosQBmPXgb}{}\renewcommand{\rhmcsOosQBmPXgb}{0.98}  % StepM-adjusted p-value (one-sided), XGBoost
\providecommand{\rhmcsOosQBmRawXgb}{}\renewcommand{\rhmcsOosQBmRawXgb}{0.92}  % single-test bootstrap p-value (one-sided), XGBoost
\providecommand{\rhmcsOosQBmPTwoXgb}{}\renewcommand{\rhmcsOosQBmPTwoXgb}{0.17}  % StepM-adjusted p-value (two-sided), XGBoost
\providecommand{\rhmcsOosQBmRawTwoXgb}{}\renewcommand{\rhmcsOosQBmRawTwoXgb}{0.16}  % single-test bootstrap p-value (two-sided), XGBoost
\providecommand{\rhmcsOosQBmPRf}{}\renewcommand{\rhmcsOosQBmPRf}{1.00}  % StepM-adjusted p-value (one-sided), Random forest
\providecommand{\rhmcsOosQBmRawRf}{}\renewcommand{\rhmcsOosQBmRawRf}{0.98}  % single-test bootstrap p-value (one-sided), Random forest
\providecommand{\rhmcsOosQBmPTwoRf}{}\renewcommand{\rhmcsOosQBmPTwoRf}{0.039}  % StepM-adjusted p-value (two-sided), Random forest
\providecommand{\rhmcsOosQBmRawTwoRf}{}\renewcommand{\rhmcsOosQBmRawTwoRf}{0.030}  % single-test bootstrap p-value (two-sided), Random forest
\providecommand{\rhmcsOosQBmNRejLo}{}\renewcommand{\rhmcsOosQBmNRejLo}{0}  % StepM rejections at 5 %, block length 5
\providecommand{\rhmcsOosQBmSpaCLo}{}\renewcommand{\rhmcsOosQBmSpaCLo}{0.53}  % SPA p-value (consistent), block length 5
\providecommand{\rhmcsOosQBmNRejTwoLo}{}\renewcommand{\rhmcsOosQBmNRejTwoLo}{4}  % ... two-sided, block length 5
\providecommand{\rhmcsOosQBmPMinLo}{}\renewcommand{\rhmcsOosQBmPMinLo}{0.97}  % smallest adjusted p, block length 5
\providecommand{\rhmcsOosQBmPTwoMinLo}{}\renewcommand{\rhmcsOosQBmPTwoMinLo}{\ensuremath{<}0.001}  % ... two-sided, block length 5
\providecommand{\rhmcsOosQBmNRejHi}{}\renewcommand{\rhmcsOosQBmNRejHi}{0}  % StepM rejections at 5 %, block length 20
\providecommand{\rhmcsOosQBmSpaCHi}{}\renewcommand{\rhmcsOosQBmSpaCHi}{0.53}  % SPA p-value (consistent), block length 20
\providecommand{\rhmcsOosQBmNRejTwoHi}{}\renewcommand{\rhmcsOosQBmNRejTwoHi}{4}  % ... two-sided, block length 20
\providecommand{\rhmcsOosQBmPMinHi}{}\renewcommand{\rhmcsOosQBmPMinHi}{0.98}  % smallest adjusted p, block length 20
\providecommand{\rhmcsOosQBmPTwoMinHi}{}\renewcommand{\rhmcsOosQBmPTwoMinHi}{0.003}  % ... two-sided, block length 20
\providecommand{\rhmcsOosQBmAllSpaC}{}\renewcommand{\rhmcsOosQBmAllSpaC}{0.53}  % SPA p-value (consistent), family OosQBmAll
\providecommand{\rhmcsOosQBmAllSpaL}{}\renewcommand{\rhmcsOosQBmAllSpaL}{0.084}  % ... lower
\providecommand{\rhmcsOosQBmAllSpaU}{}\renewcommand{\rhmcsOosQBmAllSpaU}{0.79}  % ... upper
\providecommand{\rhmcsOosQBmAllNRej}{}\renewcommand{\rhmcsOosQBmAllNRej}{0}  % StepM rejections at FWER 5 % (of 10)
\providecommand{\rhmcsOosQBmAllNRejWord}{}\renewcommand{\rhmcsOosQBmAllNRejWord}{none}  % ... as a word
\providecommand{\rhmcsOosQBmAllNRaw}{}\renewcommand{\rhmcsOosQBmAllNRaw}{0}  % single-test bootstrap p < 0.05 (no adjustment)
\providecommand{\rhmcsOosQBmAllNRawWord}{}\renewcommand{\rhmcsOosQBmAllNRawWord}{none}  % ... as a word
\providecommand{\rhmcsOosQBmAllNPos}{}\renewcommand{\rhmcsOosQBmAllNPos}{0}  % alternatives with a positive estimate (better than the benchmark)
\providecommand{\rhmcsOosQBmAllPMin}{}\renewcommand{\rhmcsOosQBmAllPMin}{0.99}  % smallest StepM-adjusted p-value
\providecommand{\rhmcsOosQBmAllPMax}{}\renewcommand{\rhmcsOosQBmAllPMax}{1.00}  % largest StepM-adjusted p-value
\providecommand{\rhmcsOosQBmAllPTwoMin}{}\renewcommand{\rhmcsOosQBmAllPTwoMin}{0.003}  % smallest two-sided StepM-adjusted p-value
\providecommand{\rhmcsOosQBmAllPTwoMax}{}\renewcommand{\rhmcsOosQBmAllPTwoMax}{0.17}  % largest two-sided StepM-adjusted p-value
\providecommand{\rhmcsOosQBmAllRawMin}{}\renewcommand{\rhmcsOosQBmAllRawMin}{0.92}  % smallest single-test bootstrap p-value
\providecommand{\rhmcsOosQBmAllRej}{}\renewcommand{\rhmcsOosQBmAllRej}{none}  % the alternatives StepM rejects at 5 % (none: none)
\providecommand{\rhmcsOosQBmAllNRejTwo}{}\renewcommand{\rhmcsOosQBmAllNRejTwo}{8}  % StepM rejections at FWER 5 %, two-sided
\providecommand{\rhmcsOosQBmAllNRejTwoWord}{}\renewcommand{\rhmcsOosQBmAllNRejTwoWord}{eight}  % ... as a word
\providecommand{\rhmcsOosQBmAllNRawTwo}{}\renewcommand{\rhmcsOosQBmAllNRawTwo}{8}  % single-test bootstrap p < 0.05, two-sided
\providecommand{\rhmcsOosQBmAllRejTwo}{}\renewcommand{\rhmcsOosQBmAllRejTwo}{the linear models on the HAR benchmark, ridge regression on the full predictor set, the lasso on the full predictor set, the elastic net on the full predictor set, LightGBM on the HAR benchmark, XGBoost on the HAR benchmark, the random forest on the HAR benchmark, and the random forest on the full predictor set}  % the alternatives the two-sided StepM rejects at 5 % (none: none)
\providecommand{\rhmcsOosQBmAllPLinHar}{}\renewcommand{\rhmcsOosQBmAllPLinHar}{1.00}  % StepM-adjusted p-value (one-sided), Linear, HAR
\providecommand{\rhmcsOosQBmAllRawLinHar}{}\renewcommand{\rhmcsOosQBmAllRawLinHar}{0.99}  % single-test bootstrap p-value (one-sided), Linear, HAR
\providecommand{\rhmcsOosQBmAllPTwoLinHar}{}\renewcommand{\rhmcsOosQBmAllPTwoLinHar}{0.014}  % StepM-adjusted p-value (two-sided), Linear, HAR
\providecommand{\rhmcsOosQBmAllRawTwoLinHar}{}\renewcommand{\rhmcsOosQBmAllRawTwoLinHar}{0.007}  % single-test bootstrap p-value (two-sided), Linear, HAR
\providecommand{\rhmcsOosQBmAllPRidgeFull}{}\renewcommand{\rhmcsOosQBmAllPRidgeFull}{1.00}  % StepM-adjusted p-value (one-sided), Ridge, Full
\providecommand{\rhmcsOosQBmAllRawRidgeFull}{}\renewcommand{\rhmcsOosQBmAllRawRidgeFull}{1.00}  % single-test bootstrap p-value (one-sided), Ridge, Full
\providecommand{\rhmcsOosQBmAllPTwoRidgeFull}{}\renewcommand{\rhmcsOosQBmAllPTwoRidgeFull}{0.003}  % StepM-adjusted p-value (two-sided), Ridge, Full
\providecommand{\rhmcsOosQBmAllRawTwoRidgeFull}{}\renewcommand{\rhmcsOosQBmAllRawTwoRidgeFull}{\ensuremath{<}0.001}  % single-test bootstrap p-value (two-sided), Ridge, Full
\providecommand{\rhmcsOosQBmAllPLassoFull}{}\renewcommand{\rhmcsOosQBmAllPLassoFull}{1.00}  % StepM-adjusted p-value (one-sided), Lasso, Full
\providecommand{\rhmcsOosQBmAllRawLassoFull}{}\renewcommand{\rhmcsOosQBmAllRawLassoFull}{1.00}  % single-test bootstrap p-value (one-sided), Lasso, Full
\providecommand{\rhmcsOosQBmAllPTwoLassoFull}{}\renewcommand{\rhmcsOosQBmAllPTwoLassoFull}{0.003}  % StepM-adjusted p-value (two-sided), Lasso, Full
\providecommand{\rhmcsOosQBmAllRawTwoLassoFull}{}\renewcommand{\rhmcsOosQBmAllRawTwoLassoFull}{\ensuremath{<}0.001}  % single-test bootstrap p-value (two-sided), Lasso, Full
\providecommand{\rhmcsOosQBmAllPEnetFull}{}\renewcommand{\rhmcsOosQBmAllPEnetFull}{1.00}  % StepM-adjusted p-value (one-sided), Elastic net, Full
\providecommand{\rhmcsOosQBmAllRawEnetFull}{}\renewcommand{\rhmcsOosQBmAllRawEnetFull}{1.00}  % single-test bootstrap p-value (one-sided), Elastic net, Full
\providecommand{\rhmcsOosQBmAllPTwoEnetFull}{}\renewcommand{\rhmcsOosQBmAllPTwoEnetFull}{0.004}  % StepM-adjusted p-value (two-sided), Elastic net, Full
\providecommand{\rhmcsOosQBmAllRawTwoEnetFull}{}\renewcommand{\rhmcsOosQBmAllRawTwoEnetFull}{\ensuremath{<}0.001}  % single-test bootstrap p-value (two-sided), Elastic net, Full
\providecommand{\rhmcsOosQBmAllPLgbmHar}{}\renewcommand{\rhmcsOosQBmAllPLgbmHar}{1.00}  % StepM-adjusted p-value (one-sided), LightGBM, HAR
\providecommand{\rhmcsOosQBmAllRawLgbmHar}{}\renewcommand{\rhmcsOosQBmAllRawLgbmHar}{0.99}  % single-test bootstrap p-value (one-sided), LightGBM, HAR
\providecommand{\rhmcsOosQBmAllPTwoLgbmHar}{}\renewcommand{\rhmcsOosQBmAllPTwoLgbmHar}{0.021}  % StepM-adjusted p-value (two-sided), LightGBM, HAR
\providecommand{\rhmcsOosQBmAllRawTwoLgbmHar}{}\renewcommand{\rhmcsOosQBmAllRawTwoLgbmHar}{0.009}  % single-test bootstrap p-value (two-sided), LightGBM, HAR
\providecommand{\rhmcsOosQBmAllPLgbmFull}{}\renewcommand{\rhmcsOosQBmAllPLgbmFull}{0.99}  % StepM-adjusted p-value (one-sided), LightGBM, Full
\providecommand{\rhmcsOosQBmAllRawLgbmFull}{}\renewcommand{\rhmcsOosQBmAllRawLgbmFull}{0.92}  % single-test bootstrap p-value (one-sided), LightGBM, Full
\providecommand{\rhmcsOosQBmAllPTwoLgbmFull}{}\renewcommand{\rhmcsOosQBmAllPTwoLgbmFull}{0.17}  % StepM-adjusted p-value (two-sided), LightGBM, Full
\providecommand{\rhmcsOosQBmAllRawTwoLgbmFull}{}\renewcommand{\rhmcsOosQBmAllRawTwoLgbmFull}{0.15}  % single-test bootstrap p-value (two-sided), LightGBM, Full
\providecommand{\rhmcsOosQBmAllPXgbHar}{}\renewcommand{\rhmcsOosQBmAllPXgbHar}{1.00}  % StepM-adjusted p-value (one-sided), XGBoost, HAR
\providecommand{\rhmcsOosQBmAllRawXgbHar}{}\renewcommand{\rhmcsOosQBmAllRawXgbHar}{0.99}  % single-test bootstrap p-value (one-sided), XGBoost, HAR
\providecommand{\rhmcsOosQBmAllPTwoXgbHar}{}\renewcommand{\rhmcsOosQBmAllPTwoXgbHar}{0.016}  % StepM-adjusted p-value (two-sided), XGBoost, HAR
\providecommand{\rhmcsOosQBmAllRawTwoXgbHar}{}\renewcommand{\rhmcsOosQBmAllRawTwoXgbHar}{0.006}  % single-test bootstrap p-value (two-sided), XGBoost, HAR
\providecommand{\rhmcsOosQBmAllPXgbFull}{}\renewcommand{\rhmcsOosQBmAllPXgbFull}{0.99}  % StepM-adjusted p-value (one-sided), XGBoost, Full
\providecommand{\rhmcsOosQBmAllRawXgbFull}{}\renewcommand{\rhmcsOosQBmAllRawXgbFull}{0.92}  % single-test bootstrap p-value (one-sided), XGBoost, Full
\providecommand{\rhmcsOosQBmAllPTwoXgbFull}{}\renewcommand{\rhmcsOosQBmAllPTwoXgbFull}{0.17}  % StepM-adjusted p-value (two-sided), XGBoost, Full
\providecommand{\rhmcsOosQBmAllRawTwoXgbFull}{}\renewcommand{\rhmcsOosQBmAllRawTwoXgbFull}{0.16}  % single-test bootstrap p-value (two-sided), XGBoost, Full
\providecommand{\rhmcsOosQBmAllPRfHar}{}\renewcommand{\rhmcsOosQBmAllPRfHar}{1.00}  % StepM-adjusted p-value (one-sided), Random forest, HAR
\providecommand{\rhmcsOosQBmAllRawRfHar}{}\renewcommand{\rhmcsOosQBmAllRawRfHar}{1.00}  % single-test bootstrap p-value (one-sided), Random forest, HAR
\providecommand{\rhmcsOosQBmAllPTwoRfHar}{}\renewcommand{\rhmcsOosQBmAllPTwoRfHar}{0.006}  % StepM-adjusted p-value (two-sided), Random forest, HAR
\providecommand{\rhmcsOosQBmAllRawTwoRfHar}{}\renewcommand{\rhmcsOosQBmAllRawTwoRfHar}{0.002}  % single-test bootstrap p-value (two-sided), Random forest, HAR
\providecommand{\rhmcsOosQBmAllPRfFull}{}\renewcommand{\rhmcsOosQBmAllPRfFull}{1.00}  % StepM-adjusted p-value (one-sided), Random forest, Full
\providecommand{\rhmcsOosQBmAllRawRfFull}{}\renewcommand{\rhmcsOosQBmAllRawRfFull}{0.98}  % single-test bootstrap p-value (one-sided), Random forest, Full
\providecommand{\rhmcsOosQBmAllPTwoRfFull}{}\renewcommand{\rhmcsOosQBmAllPTwoRfFull}{0.039}  % StepM-adjusted p-value (two-sided), Random forest, Full
\providecommand{\rhmcsOosQBmAllRawTwoRfFull}{}\renewcommand{\rhmcsOosQBmAllRawTwoRfFull}{0.030}  % single-test bootstrap p-value (two-sided), Random forest, Full
\providecommand{\rhmcsOosQBmAllNRejLo}{}\renewcommand{\rhmcsOosQBmAllNRejLo}{0}  % StepM rejections at 5 %, block length 5
\providecommand{\rhmcsOosQBmAllSpaCLo}{}\renewcommand{\rhmcsOosQBmAllSpaCLo}{0.53}  % SPA p-value (consistent), block length 5
\providecommand{\rhmcsOosQBmAllNRejTwoLo}{}\renewcommand{\rhmcsOosQBmAllNRejTwoLo}{8}  % ... two-sided, block length 5
\providecommand{\rhmcsOosQBmAllPMinLo}{}\renewcommand{\rhmcsOosQBmAllPMinLo}{0.99}  % smallest adjusted p, block length 5
\providecommand{\rhmcsOosQBmAllPTwoMinLo}{}\renewcommand{\rhmcsOosQBmAllPTwoMinLo}{0.002}  % ... two-sided, block length 5
\providecommand{\rhmcsOosQBmAllNRejHi}{}\renewcommand{\rhmcsOosQBmAllNRejHi}{0}  % StepM rejections at 5 %, block length 20
\providecommand{\rhmcsOosQBmAllSpaCHi}{}\renewcommand{\rhmcsOosQBmAllSpaCHi}{0.53}  % SPA p-value (consistent), block length 20
\providecommand{\rhmcsOosQBmAllNRejTwoHi}{}\renewcommand{\rhmcsOosQBmAllNRejTwoHi}{8}  % ... two-sided, block length 20
\providecommand{\rhmcsOosQBmAllPMinHi}{}\renewcommand{\rhmcsOosQBmAllPMinHi}{0.99}  % smallest adjusted p, block length 20
\providecommand{\rhmcsOosQBmAllPTwoMinHi}{}\renewcommand{\rhmcsOosQBmAllPTwoMinHi}{0.005}  % ... two-sided, block length 20
\providecommand{\rhmcsOosMAfSpaC}{}\renewcommand{\rhmcsOosMAfSpaC}{0.67}  % SPA p-value (consistent), family OosMAf
\providecommand{\rhmcsOosMAfSpaL}{}\renewcommand{\rhmcsOosMAfSpaL}{0.19}  % ... lower
\providecommand{\rhmcsOosMAfSpaU}{}\renewcommand{\rhmcsOosMAfSpaU}{0.67}  % ... upper
\providecommand{\rhmcsOosMAfNRej}{}\renewcommand{\rhmcsOosMAfNRej}{0}  % StepM rejections at FWER 5 % (of 6)
\providecommand{\rhmcsOosMAfNRejWord}{}\renewcommand{\rhmcsOosMAfNRejWord}{none}  % ... as a word
\providecommand{\rhmcsOosMAfNRaw}{}\renewcommand{\rhmcsOosMAfNRaw}{0}  % single-test bootstrap p < 0.05 (no adjustment)
\providecommand{\rhmcsOosMAfNRawWord}{}\renewcommand{\rhmcsOosMAfNRawWord}{none}  % ... as a word
\providecommand{\rhmcsOosMAfNPos}{}\renewcommand{\rhmcsOosMAfNPos}{0}  % alternatives with a positive estimate (better than the benchmark)
\providecommand{\rhmcsOosMAfPMin}{}\renewcommand{\rhmcsOosMAfPMin}{0.92}  % smallest StepM-adjusted p-value
\providecommand{\rhmcsOosMAfPMax}{}\renewcommand{\rhmcsOosMAfPMax}{0.92}  % largest StepM-adjusted p-value
\providecommand{\rhmcsOosMAfPTwoMin}{}\renewcommand{\rhmcsOosMAfPTwoMin}{0.32}  % smallest two-sided StepM-adjusted p-value
\providecommand{\rhmcsOosMAfPTwoMax}{}\renewcommand{\rhmcsOosMAfPTwoMax}{0.34}  % largest two-sided StepM-adjusted p-value
\providecommand{\rhmcsOosMAfRawMin}{}\renewcommand{\rhmcsOosMAfRawMin}{0.86}  % smallest single-test bootstrap p-value
\providecommand{\rhmcsOosMAfRej}{}\renewcommand{\rhmcsOosMAfRej}{none}  % the alternatives StepM rejects at 5 % (none: none)
\providecommand{\rhmcsOosMAfNRejTwo}{}\renewcommand{\rhmcsOosMAfNRejTwo}{0}  % StepM rejections at FWER 5 %, two-sided
\providecommand{\rhmcsOosMAfNRejTwoWord}{}\renewcommand{\rhmcsOosMAfNRejTwoWord}{none}  % ... as a word
\providecommand{\rhmcsOosMAfNRawTwo}{}\renewcommand{\rhmcsOosMAfNRawTwo}{0}  % single-test bootstrap p < 0.05, two-sided
\providecommand{\rhmcsOosMAfRejTwo}{}\renewcommand{\rhmcsOosMAfRejTwo}{none}  % the alternatives the two-sided StepM rejects at 5 % (none: none)
\providecommand{\rhmcsOosMAfPRidge}{}\renewcommand{\rhmcsOosMAfPRidge}{0.92}  % StepM-adjusted p-value (one-sided), Ridge
\providecommand{\rhmcsOosMAfRawRidge}{}\renewcommand{\rhmcsOosMAfRawRidge}{0.88}  % single-test bootstrap p-value (one-sided), Ridge
\providecommand{\rhmcsOosMAfPTwoRidge}{}\renewcommand{\rhmcsOosMAfPTwoRidge}{0.33}  % StepM-adjusted p-value (two-sided), Ridge
\providecommand{\rhmcsOosMAfRawTwoRidge}{}\renewcommand{\rhmcsOosMAfRawTwoRidge}{0.23}  % single-test bootstrap p-value (two-sided), Ridge
\providecommand{\rhmcsOosMAfPLasso}{}\renewcommand{\rhmcsOosMAfPLasso}{0.92}  % StepM-adjusted p-value (one-sided), Lasso
\providecommand{\rhmcsOosMAfRawLasso}{}\renewcommand{\rhmcsOosMAfRawLasso}{0.88}  % single-test bootstrap p-value (one-sided), Lasso
\providecommand{\rhmcsOosMAfPTwoLasso}{}\renewcommand{\rhmcsOosMAfPTwoLasso}{0.32}  % StepM-adjusted p-value (two-sided), Lasso
\providecommand{\rhmcsOosMAfRawTwoLasso}{}\renewcommand{\rhmcsOosMAfRawTwoLasso}{0.22}  % single-test bootstrap p-value (two-sided), Lasso
\providecommand{\rhmcsOosMAfPEnet}{}\renewcommand{\rhmcsOosMAfPEnet}{0.92}  % StepM-adjusted p-value (one-sided), Elastic net
\providecommand{\rhmcsOosMAfRawEnet}{}\renewcommand{\rhmcsOosMAfRawEnet}{0.88}  % single-test bootstrap p-value (one-sided), Elastic net
\providecommand{\rhmcsOosMAfPTwoEnet}{}\renewcommand{\rhmcsOosMAfPTwoEnet}{0.34}  % StepM-adjusted p-value (two-sided), Elastic net
\providecommand{\rhmcsOosMAfRawTwoEnet}{}\renewcommand{\rhmcsOosMAfRawTwoEnet}{0.23}  % single-test bootstrap p-value (two-sided), Elastic net
\providecommand{\rhmcsOosMAfPLgbm}{}\renewcommand{\rhmcsOosMAfPLgbm}{0.92}  % StepM-adjusted p-value (one-sided), LightGBM
\providecommand{\rhmcsOosMAfRawLgbm}{}\renewcommand{\rhmcsOosMAfRawLgbm}{0.86}  % single-test bootstrap p-value (one-sided), LightGBM
\providecommand{\rhmcsOosMAfPTwoLgbm}{}\renewcommand{\rhmcsOosMAfPTwoLgbm}{0.34}  % StepM-adjusted p-value (two-sided), LightGBM
\providecommand{\rhmcsOosMAfRawTwoLgbm}{}\renewcommand{\rhmcsOosMAfRawTwoLgbm}{0.31}  % single-test bootstrap p-value (two-sided), LightGBM
\providecommand{\rhmcsOosMAfPXgb}{}\renewcommand{\rhmcsOosMAfPXgb}{0.92}  % StepM-adjusted p-value (one-sided), XGBoost
\providecommand{\rhmcsOosMAfRawXgb}{}\renewcommand{\rhmcsOosMAfRawXgb}{0.86}  % single-test bootstrap p-value (one-sided), XGBoost
\providecommand{\rhmcsOosMAfPTwoXgb}{}\renewcommand{\rhmcsOosMAfPTwoXgb}{0.34}  % StepM-adjusted p-value (two-sided), XGBoost
\providecommand{\rhmcsOosMAfRawTwoXgb}{}\renewcommand{\rhmcsOosMAfRawTwoXgb}{0.29}  % single-test bootstrap p-value (two-sided), XGBoost
\providecommand{\rhmcsOosMAfPRf}{}\renewcommand{\rhmcsOosMAfPRf}{0.92}  % StepM-adjusted p-value (one-sided), Random forest
\providecommand{\rhmcsOosMAfRawRf}{}\renewcommand{\rhmcsOosMAfRawRf}{0.86}  % single-test bootstrap p-value (one-sided), Random forest
\providecommand{\rhmcsOosMAfPTwoRf}{}\renewcommand{\rhmcsOosMAfPTwoRf}{0.34}  % StepM-adjusted p-value (two-sided), Random forest
\providecommand{\rhmcsOosMAfRawTwoRf}{}\renewcommand{\rhmcsOosMAfRawTwoRf}{0.28}  % single-test bootstrap p-value (two-sided), Random forest
\providecommand{\rhmcsOosMAfNRejLo}{}\renewcommand{\rhmcsOosMAfNRejLo}{0}  % StepM rejections at 5 %, block length 5
\providecommand{\rhmcsOosMAfSpaCLo}{}\renewcommand{\rhmcsOosMAfSpaCLo}{0.73}  % SPA p-value (consistent), block length 5
\providecommand{\rhmcsOosMAfNRejTwoLo}{}\renewcommand{\rhmcsOosMAfNRejTwoLo}{0}  % ... two-sided, block length 5
\providecommand{\rhmcsOosMAfPMinLo}{}\renewcommand{\rhmcsOosMAfPMinLo}{0.97}  % smallest adjusted p, block length 5
\providecommand{\rhmcsOosMAfPTwoMinLo}{}\renewcommand{\rhmcsOosMAfPTwoMinLo}{0.21}  % ... two-sided, block length 5
\providecommand{\rhmcsOosMAfNRejHi}{}\renewcommand{\rhmcsOosMAfNRejHi}{0}  % StepM rejections at 5 %, block length 20
\providecommand{\rhmcsOosMAfSpaCHi}{}\renewcommand{\rhmcsOosMAfSpaCHi}{0.62}  % SPA p-value (consistent), block length 20
\providecommand{\rhmcsOosMAfNRejTwoHi}{}\renewcommand{\rhmcsOosMAfNRejTwoHi}{0}  % ... two-sided, block length 20
\providecommand{\rhmcsOosMAfPMinHi}{}\renewcommand{\rhmcsOosMAfPMinHi}{0.88}  % smallest adjusted p, block length 20
\providecommand{\rhmcsOosMAfPTwoMinHi}{}\renewcommand{\rhmcsOosMAfPTwoMinHi}{0.38}  % ... two-sided, block length 20
\providecommand{\rhmcsOosMBmSpaC}{}\renewcommand{\rhmcsOosMBmSpaC}{0.63}  % SPA p-value (consistent), family OosMBm
\providecommand{\rhmcsOosMBmSpaL}{}\renewcommand{\rhmcsOosMBmSpaL}{0.33}  % ... lower
\providecommand{\rhmcsOosMBmSpaU}{}\renewcommand{\rhmcsOosMBmSpaU}{0.63}  % ... upper
\providecommand{\rhmcsOosMBmNRej}{}\renewcommand{\rhmcsOosMBmNRej}{0}  % StepM rejections at FWER 5 % (of 6)
\providecommand{\rhmcsOosMBmNRejWord}{}\renewcommand{\rhmcsOosMBmNRejWord}{none}  % ... as a word
\providecommand{\rhmcsOosMBmNRaw}{}\renewcommand{\rhmcsOosMBmNRaw}{0}  % single-test bootstrap p < 0.05 (no adjustment)
\providecommand{\rhmcsOosMBmNRawWord}{}\renewcommand{\rhmcsOosMBmNRawWord}{none}  % ... as a word
\providecommand{\rhmcsOosMBmNPos}{}\renewcommand{\rhmcsOosMBmNPos}{1}  % alternatives with a positive estimate (better than the benchmark)
\providecommand{\rhmcsOosMBmPMin}{}\renewcommand{\rhmcsOosMBmPMin}{0.63}  % smallest StepM-adjusted p-value
\providecommand{\rhmcsOosMBmPMax}{}\renewcommand{\rhmcsOosMBmPMax}{0.88}  % largest StepM-adjusted p-value
\providecommand{\rhmcsOosMBmPTwoMin}{}\renewcommand{\rhmcsOosMBmPTwoMin}{0.47}  % smallest two-sided StepM-adjusted p-value
\providecommand{\rhmcsOosMBmPTwoMax}{}\renewcommand{\rhmcsOosMBmPTwoMax}{0.74}  % largest two-sided StepM-adjusted p-value
\providecommand{\rhmcsOosMBmRawMin}{}\renewcommand{\rhmcsOosMBmRawMin}{0.33}  % smallest single-test bootstrap p-value
\providecommand{\rhmcsOosMBmRej}{}\renewcommand{\rhmcsOosMBmRej}{none}  % the alternatives StepM rejects at 5 % (none: none)
\providecommand{\rhmcsOosMBmNRejTwo}{}\renewcommand{\rhmcsOosMBmNRejTwo}{0}  % StepM rejections at FWER 5 %, two-sided
\providecommand{\rhmcsOosMBmNRejTwoWord}{}\renewcommand{\rhmcsOosMBmNRejTwoWord}{none}  % ... as a word
\providecommand{\rhmcsOosMBmNRawTwo}{}\renewcommand{\rhmcsOosMBmNRawTwo}{0}  % single-test bootstrap p < 0.05, two-sided
\providecommand{\rhmcsOosMBmRejTwo}{}\renewcommand{\rhmcsOosMBmRejTwo}{none}  % the alternatives the two-sided StepM rejects at 5 % (none: none)
\providecommand{\rhmcsOosMBmPRidge}{}\renewcommand{\rhmcsOosMBmPRidge}{0.88}  % StepM-adjusted p-value (one-sided), Ridge
\providecommand{\rhmcsOosMBmRawRidge}{}\renewcommand{\rhmcsOosMBmRawRidge}{0.88}  % single-test bootstrap p-value (one-sided), Ridge
\providecommand{\rhmcsOosMBmPTwoRidge}{}\renewcommand{\rhmcsOosMBmPTwoRidge}{0.48}  % StepM-adjusted p-value (two-sided), Ridge
\providecommand{\rhmcsOosMBmRawTwoRidge}{}\renewcommand{\rhmcsOosMBmRawTwoRidge}{0.26}  % single-test bootstrap p-value (two-sided), Ridge
\providecommand{\rhmcsOosMBmPLasso}{}\renewcommand{\rhmcsOosMBmPLasso}{0.88}  % StepM-adjusted p-value (one-sided), Lasso
\providecommand{\rhmcsOosMBmRawLasso}{}\renewcommand{\rhmcsOosMBmRawLasso}{0.88}  % single-test bootstrap p-value (one-sided), Lasso
\providecommand{\rhmcsOosMBmPTwoLasso}{}\renewcommand{\rhmcsOosMBmPTwoLasso}{0.47}  % StepM-adjusted p-value (two-sided), Lasso
\providecommand{\rhmcsOosMBmRawTwoLasso}{}\renewcommand{\rhmcsOosMBmRawTwoLasso}{0.253}  % single-test bootstrap p-value (two-sided), Lasso
\providecommand{\rhmcsOosMBmPEnet}{}\renewcommand{\rhmcsOosMBmPEnet}{0.88}  % StepM-adjusted p-value (one-sided), Elastic net
\providecommand{\rhmcsOosMBmRawEnet}{}\renewcommand{\rhmcsOosMBmRawEnet}{0.88}  % single-test bootstrap p-value (one-sided), Elastic net
\providecommand{\rhmcsOosMBmPTwoEnet}{}\renewcommand{\rhmcsOosMBmPTwoEnet}{0.49}  % StepM-adjusted p-value (two-sided), Elastic net
\providecommand{\rhmcsOosMBmRawTwoEnet}{}\renewcommand{\rhmcsOosMBmRawTwoEnet}{0.26}  % single-test bootstrap p-value (two-sided), Elastic net
\providecommand{\rhmcsOosMBmPLgbm}{}\renewcommand{\rhmcsOosMBmPLgbm}{0.87}  % StepM-adjusted p-value (one-sided), LightGBM
\providecommand{\rhmcsOosMBmRawLgbm}{}\renewcommand{\rhmcsOosMBmRawLgbm}{0.81}  % single-test bootstrap p-value (one-sided), LightGBM
\providecommand{\rhmcsOosMBmPTwoLgbm}{}\renewcommand{\rhmcsOosMBmPTwoLgbm}{0.62}  % StepM-adjusted p-value (two-sided), LightGBM
\providecommand{\rhmcsOosMBmRawTwoLgbm}{}\renewcommand{\rhmcsOosMBmRawTwoLgbm}{0.40}  % single-test bootstrap p-value (two-sided), LightGBM
\providecommand{\rhmcsOosMBmPXgb}{}\renewcommand{\rhmcsOosMBmPXgb}{0.86}  % StepM-adjusted p-value (one-sided), XGBoost
\providecommand{\rhmcsOosMBmRawXgb}{}\renewcommand{\rhmcsOosMBmRawXgb}{0.76}  % single-test bootstrap p-value (one-sided), XGBoost
\providecommand{\rhmcsOosMBmPTwoXgb}{}\renewcommand{\rhmcsOosMBmPTwoXgb}{0.74}  % StepM-adjusted p-value (two-sided), XGBoost
\providecommand{\rhmcsOosMBmRawTwoXgb}{}\renewcommand{\rhmcsOosMBmRawTwoXgb}{0.55}  % single-test bootstrap p-value (two-sided), XGBoost
\providecommand{\rhmcsOosMBmPRf}{}\renewcommand{\rhmcsOosMBmPRf}{0.63}  % StepM-adjusted p-value (one-sided), Random forest
\providecommand{\rhmcsOosMBmRawRf}{}\renewcommand{\rhmcsOosMBmRawRf}{0.33}  % single-test bootstrap p-value (one-sided), Random forest
\providecommand{\rhmcsOosMBmPTwoRf}{}\renewcommand{\rhmcsOosMBmPTwoRf}{0.74}  % StepM-adjusted p-value (two-sided), Random forest
\providecommand{\rhmcsOosMBmRawTwoRf}{}\renewcommand{\rhmcsOosMBmRawTwoRf}{0.72}  % single-test bootstrap p-value (two-sided), Random forest
\providecommand{\rhmcsOosMBmNRejLo}{}\renewcommand{\rhmcsOosMBmNRejLo}{0}  % StepM rejections at 5 %, block length 5
\providecommand{\rhmcsOosMBmSpaCLo}{}\renewcommand{\rhmcsOosMBmSpaCLo}{0.72}  % SPA p-value (consistent), block length 5
\providecommand{\rhmcsOosMBmNRejTwoLo}{}\renewcommand{\rhmcsOosMBmNRejTwoLo}{0}  % ... two-sided, block length 5
\providecommand{\rhmcsOosMBmPMinLo}{}\renewcommand{\rhmcsOosMBmPMinLo}{0.72}  % smallest adjusted p, block length 5
\providecommand{\rhmcsOosMBmPTwoMinLo}{}\renewcommand{\rhmcsOosMBmPTwoMinLo}{0.32}  % ... two-sided, block length 5
\providecommand{\rhmcsOosMBmNRejHi}{}\renewcommand{\rhmcsOosMBmNRejHi}{0}  % StepM rejections at 5 %, block length 20
\providecommand{\rhmcsOosMBmSpaCHi}{}\renewcommand{\rhmcsOosMBmSpaCHi}{0.61}  % SPA p-value (consistent), block length 20
\providecommand{\rhmcsOosMBmNRejTwoHi}{}\renewcommand{\rhmcsOosMBmNRejTwoHi}{0}  % ... two-sided, block length 20
\providecommand{\rhmcsOosMBmPMinHi}{}\renewcommand{\rhmcsOosMBmPMinHi}{0.61}  % smallest adjusted p, block length 20
\providecommand{\rhmcsOosMBmPTwoMinHi}{}\renewcommand{\rhmcsOosMBmPTwoMinHi}{0.58}  % ... two-sided, block length 20
\providecommand{\rhmcsOosMBmAllSpaC}{}\renewcommand{\rhmcsOosMBmAllSpaC}{0.21}  % SPA p-value (consistent), family OosMBmAll
\providecommand{\rhmcsOosMBmAllSpaL}{}\renewcommand{\rhmcsOosMBmAllSpaL}{0.20}  % ... lower
\providecommand{\rhmcsOosMBmAllSpaU}{}\renewcommand{\rhmcsOosMBmAllSpaU}{0.21}  % ... upper
\providecommand{\rhmcsOosMBmAllNRej}{}\renewcommand{\rhmcsOosMBmAllNRej}{0}  % StepM rejections at FWER 5 % (of 10)
\providecommand{\rhmcsOosMBmAllNRejWord}{}\renewcommand{\rhmcsOosMBmAllNRejWord}{none}  % ... as a word
\providecommand{\rhmcsOosMBmAllNRaw}{}\renewcommand{\rhmcsOosMBmAllNRaw}{0}  % single-test bootstrap p < 0.05 (no adjustment)
\providecommand{\rhmcsOosMBmAllNRawWord}{}\renewcommand{\rhmcsOosMBmAllNRawWord}{none}  % ... as a word
\providecommand{\rhmcsOosMBmAllNPos}{}\renewcommand{\rhmcsOosMBmAllNPos}{5}  % alternatives with a positive estimate (better than the benchmark)
\providecommand{\rhmcsOosMBmAllPMin}{}\renewcommand{\rhmcsOosMBmAllPMin}{0.21}  % smallest StepM-adjusted p-value
\providecommand{\rhmcsOosMBmAllPMax}{}\renewcommand{\rhmcsOosMBmAllPMax}{0.88}  % largest StepM-adjusted p-value
\providecommand{\rhmcsOosMBmAllPTwoMin}{}\renewcommand{\rhmcsOosMBmAllPTwoMin}{0.40}  % smallest two-sided StepM-adjusted p-value
\providecommand{\rhmcsOosMBmAllPTwoMax}{}\renewcommand{\rhmcsOosMBmAllPTwoMax}{0.88}  % largest two-sided StepM-adjusted p-value
\providecommand{\rhmcsOosMBmAllRawMin}{}\renewcommand{\rhmcsOosMBmAllRawMin}{0.094}  % smallest single-test bootstrap p-value
\providecommand{\rhmcsOosMBmAllRej}{}\renewcommand{\rhmcsOosMBmAllRej}{none}  % the alternatives StepM rejects at 5 % (none: none)
\providecommand{\rhmcsOosMBmAllNRejTwo}{}\renewcommand{\rhmcsOosMBmAllNRejTwo}{0}  % StepM rejections at FWER 5 %, two-sided
\providecommand{\rhmcsOosMBmAllNRejTwoWord}{}\renewcommand{\rhmcsOosMBmAllNRejTwoWord}{none}  % ... as a word
\providecommand{\rhmcsOosMBmAllNRawTwo}{}\renewcommand{\rhmcsOosMBmAllNRawTwo}{0}  % single-test bootstrap p < 0.05, two-sided
\providecommand{\rhmcsOosMBmAllRejTwo}{}\renewcommand{\rhmcsOosMBmAllRejTwo}{none}  % the alternatives the two-sided StepM rejects at 5 % (none: none)
\providecommand{\rhmcsOosMBmAllPLinHar}{}\renewcommand{\rhmcsOosMBmAllPLinHar}{0.73}  % StepM-adjusted p-value (one-sided), Linear, HAR
\providecommand{\rhmcsOosMBmAllRawLinHar}{}\renewcommand{\rhmcsOosMBmAllRawLinHar}{0.41}  % single-test bootstrap p-value (one-sided), Linear, HAR
\providecommand{\rhmcsOosMBmAllPTwoLinHar}{}\renewcommand{\rhmcsOosMBmAllPTwoLinHar}{0.88}  % StepM-adjusted p-value (two-sided), Linear, HAR
\providecommand{\rhmcsOosMBmAllRawTwoLinHar}{}\renewcommand{\rhmcsOosMBmAllRawTwoLinHar}{0.88}  % single-test bootstrap p-value (two-sided), Linear, HAR
\providecommand{\rhmcsOosMBmAllPRidgeFull}{}\renewcommand{\rhmcsOosMBmAllPRidgeFull}{0.88}  % StepM-adjusted p-value (one-sided), Ridge, Full
\providecommand{\rhmcsOosMBmAllRawRidgeFull}{}\renewcommand{\rhmcsOosMBmAllRawRidgeFull}{0.88}  % single-test bootstrap p-value (one-sided), Ridge, Full
\providecommand{\rhmcsOosMBmAllPTwoRidgeFull}{}\renewcommand{\rhmcsOosMBmAllPTwoRidgeFull}{0.56}  % StepM-adjusted p-value (two-sided), Ridge, Full
\providecommand{\rhmcsOosMBmAllRawTwoRidgeFull}{}\renewcommand{\rhmcsOosMBmAllRawTwoRidgeFull}{0.26}  % single-test bootstrap p-value (two-sided), Ridge, Full
\providecommand{\rhmcsOosMBmAllPLassoFull}{}\renewcommand{\rhmcsOosMBmAllPLassoFull}{0.88}  % StepM-adjusted p-value (one-sided), Lasso, Full
\providecommand{\rhmcsOosMBmAllRawLassoFull}{}\renewcommand{\rhmcsOosMBmAllRawLassoFull}{0.88}  % single-test bootstrap p-value (one-sided), Lasso, Full
\providecommand{\rhmcsOosMBmAllPTwoLassoFull}{}\renewcommand{\rhmcsOosMBmAllPTwoLassoFull}{0.54}  % StepM-adjusted p-value (two-sided), Lasso, Full
\providecommand{\rhmcsOosMBmAllRawTwoLassoFull}{}\renewcommand{\rhmcsOosMBmAllRawTwoLassoFull}{0.253}  % single-test bootstrap p-value (two-sided), Lasso, Full
\providecommand{\rhmcsOosMBmAllPEnetFull}{}\renewcommand{\rhmcsOosMBmAllPEnetFull}{0.88}  % StepM-adjusted p-value (one-sided), Elastic net, Full
\providecommand{\rhmcsOosMBmAllRawEnetFull}{}\renewcommand{\rhmcsOosMBmAllRawEnetFull}{0.88}  % single-test bootstrap p-value (one-sided), Elastic net, Full
\providecommand{\rhmcsOosMBmAllPTwoEnetFull}{}\renewcommand{\rhmcsOosMBmAllPTwoEnetFull}{0.57}  % StepM-adjusted p-value (two-sided), Elastic net, Full
\providecommand{\rhmcsOosMBmAllRawTwoEnetFull}{}\renewcommand{\rhmcsOosMBmAllRawTwoEnetFull}{0.26}  % single-test bootstrap p-value (two-sided), Elastic net, Full
\providecommand{\rhmcsOosMBmAllPLgbmHar}{}\renewcommand{\rhmcsOosMBmAllPLgbmHar}{0.30}  % StepM-adjusted p-value (one-sided), LightGBM, HAR
\providecommand{\rhmcsOosMBmAllRawLgbmHar}{}\renewcommand{\rhmcsOosMBmAllRawLgbmHar}{0.12}  % single-test bootstrap p-value (one-sided), LightGBM, HAR
\providecommand{\rhmcsOosMBmAllPTwoLgbmHar}{}\renewcommand{\rhmcsOosMBmAllPTwoLgbmHar}{0.54}  % StepM-adjusted p-value (two-sided), LightGBM, HAR
\providecommand{\rhmcsOosMBmAllRawTwoLgbmHar}{}\renewcommand{\rhmcsOosMBmAllRawTwoLgbmHar}{0.24}  % single-test bootstrap p-value (two-sided), LightGBM, HAR
\providecommand{\rhmcsOosMBmAllPLgbmFull}{}\renewcommand{\rhmcsOosMBmAllPLgbmFull}{0.87}  % StepM-adjusted p-value (one-sided), LightGBM, Full
\providecommand{\rhmcsOosMBmAllRawLgbmFull}{}\renewcommand{\rhmcsOosMBmAllRawLgbmFull}{0.81}  % single-test bootstrap p-value (one-sided), LightGBM, Full
\providecommand{\rhmcsOosMBmAllPTwoLgbmFull}{}\renewcommand{\rhmcsOosMBmAllPTwoLgbmFull}{0.70}  % StepM-adjusted p-value (two-sided), LightGBM, Full
\providecommand{\rhmcsOosMBmAllRawTwoLgbmFull}{}\renewcommand{\rhmcsOosMBmAllRawTwoLgbmFull}{0.40}  % single-test bootstrap p-value (two-sided), LightGBM, Full
\providecommand{\rhmcsOosMBmAllPXgbHar}{}\renewcommand{\rhmcsOosMBmAllPXgbHar}{0.21}  % StepM-adjusted p-value (one-sided), XGBoost, HAR
\providecommand{\rhmcsOosMBmAllRawXgbHar}{}\renewcommand{\rhmcsOosMBmAllRawXgbHar}{0.094}  % single-test bootstrap p-value (one-sided), XGBoost, HAR
\providecommand{\rhmcsOosMBmAllPTwoXgbHar}{}\renewcommand{\rhmcsOosMBmAllPTwoXgbHar}{0.40}  % StepM-adjusted p-value (two-sided), XGBoost, HAR
\providecommand{\rhmcsOosMBmAllRawTwoXgbHar}{}\renewcommand{\rhmcsOosMBmAllRawTwoXgbHar}{0.16}  % single-test bootstrap p-value (two-sided), XGBoost, HAR
\providecommand{\rhmcsOosMBmAllPXgbFull}{}\renewcommand{\rhmcsOosMBmAllPXgbFull}{0.86}  % StepM-adjusted p-value (one-sided), XGBoost, Full
\providecommand{\rhmcsOosMBmAllRawXgbFull}{}\renewcommand{\rhmcsOosMBmAllRawXgbFull}{0.76}  % single-test bootstrap p-value (one-sided), XGBoost, Full
\providecommand{\rhmcsOosMBmAllPTwoXgbFull}{}\renewcommand{\rhmcsOosMBmAllPTwoXgbFull}{0.82}  % StepM-adjusted p-value (two-sided), XGBoost, Full
\providecommand{\rhmcsOosMBmAllRawTwoXgbFull}{}\renewcommand{\rhmcsOosMBmAllRawTwoXgbFull}{0.55}  % single-test bootstrap p-value (two-sided), XGBoost, Full
\providecommand{\rhmcsOosMBmAllPRfHar}{}\renewcommand{\rhmcsOosMBmAllPRfHar}{0.21}  % StepM-adjusted p-value (one-sided), Random forest, HAR
\providecommand{\rhmcsOosMBmAllRawRfHar}{}\renewcommand{\rhmcsOosMBmAllRawRfHar}{0.100}  % single-test bootstrap p-value (one-sided), Random forest, HAR
\providecommand{\rhmcsOosMBmAllPTwoRfHar}{}\renewcommand{\rhmcsOosMBmAllPTwoRfHar}{0.42}  % StepM-adjusted p-value (two-sided), Random forest, HAR
\providecommand{\rhmcsOosMBmAllRawTwoRfHar}{}\renewcommand{\rhmcsOosMBmAllRawTwoRfHar}{0.18}  % single-test bootstrap p-value (two-sided), Random forest, HAR
\providecommand{\rhmcsOosMBmAllPRfFull}{}\renewcommand{\rhmcsOosMBmAllPRfFull}{0.69}  % StepM-adjusted p-value (one-sided), Random forest, Full
\providecommand{\rhmcsOosMBmAllRawRfFull}{}\renewcommand{\rhmcsOosMBmAllRawRfFull}{0.33}  % single-test bootstrap p-value (one-sided), Random forest, Full
\providecommand{\rhmcsOosMBmAllPTwoRfFull}{}\renewcommand{\rhmcsOosMBmAllPTwoRfFull}{0.85}  % StepM-adjusted p-value (two-sided), Random forest, Full
\providecommand{\rhmcsOosMBmAllRawTwoRfFull}{}\renewcommand{\rhmcsOosMBmAllRawTwoRfFull}{0.72}  % single-test bootstrap p-value (two-sided), Random forest, Full
\providecommand{\rhmcsOosMBmAllNRejLo}{}\renewcommand{\rhmcsOosMBmAllNRejLo}{0}  % StepM rejections at 5 %, block length 5
\providecommand{\rhmcsOosMBmAllSpaCLo}{}\renewcommand{\rhmcsOosMBmAllSpaCLo}{0.14}  % SPA p-value (consistent), block length 5
\providecommand{\rhmcsOosMBmAllNRejTwoLo}{}\renewcommand{\rhmcsOosMBmAllNRejTwoLo}{0}  % ... two-sided, block length 5
\providecommand{\rhmcsOosMBmAllPMinLo}{}\renewcommand{\rhmcsOosMBmAllPMinLo}{0.14}  % smallest adjusted p, block length 5
\providecommand{\rhmcsOosMBmAllPTwoMinLo}{}\renewcommand{\rhmcsOosMBmAllPTwoMinLo}{0.30}  % ... two-sided, block length 5
\providecommand{\rhmcsOosMBmAllNRejHi}{}\renewcommand{\rhmcsOosMBmAllNRejHi}{0}  % StepM rejections at 5 %, block length 20
\providecommand{\rhmcsOosMBmAllSpaCHi}{}\renewcommand{\rhmcsOosMBmAllSpaCHi}{0.31}  % SPA p-value (consistent), block length 20
\providecommand{\rhmcsOosMBmAllNRejTwoHi}{}\renewcommand{\rhmcsOosMBmAllNRejTwoHi}{0}  % ... two-sided, block length 20
\providecommand{\rhmcsOosMBmAllPMinHi}{}\renewcommand{\rhmcsOosMBmAllPMinHi}{0.31}  % smallest adjusted p, block length 20
\providecommand{\rhmcsOosMBmAllPTwoMinHi}{}\renewcommand{\rhmcsOosMBmAllPTwoMinHi}{0.63}  % ... two-sided, block length 20
\providecommand{\rhmcsChQBmSpaC}{}\renewcommand{\rhmcsChQBmSpaC}{0.54}  % SPA p-value (consistent), family ChQBm
\providecommand{\rhmcsChQBmSpaL}{}\renewcommand{\rhmcsChQBmSpaL}{0.11}  % ... lower
\providecommand{\rhmcsChQBmSpaU}{}\renewcommand{\rhmcsChQBmSpaU}{0.67}  % ... upper
\providecommand{\rhmcsChQBmNRej}{}\renewcommand{\rhmcsChQBmNRej}{0}  % StepM rejections at FWER 5 % (of 6)
\providecommand{\rhmcsChQBmNRejWord}{}\renewcommand{\rhmcsChQBmNRejWord}{none}  % ... as a word
\providecommand{\rhmcsChQBmNRaw}{}\renewcommand{\rhmcsChQBmNRaw}{0}  % single-test bootstrap p < 0.05 (no adjustment)
\providecommand{\rhmcsChQBmNRawWord}{}\renewcommand{\rhmcsChQBmNRawWord}{none}  % ... as a word
\providecommand{\rhmcsChQBmNPos}{}\renewcommand{\rhmcsChQBmNPos}{0}  % alternatives with a positive estimate (better than the benchmark)
\providecommand{\rhmcsChQBmPMin}{}\renewcommand{\rhmcsChQBmPMin}{0.95}  % smallest StepM-adjusted p-value
\providecommand{\rhmcsChQBmPMax}{}\renewcommand{\rhmcsChQBmPMax}{1.00}  % largest StepM-adjusted p-value
\providecommand{\rhmcsChQBmPTwoMin}{}\renewcommand{\rhmcsChQBmPTwoMin}{0.004}  % smallest two-sided StepM-adjusted p-value
\providecommand{\rhmcsChQBmPTwoMax}{}\renewcommand{\rhmcsChQBmPTwoMax}{0.22}  % largest two-sided StepM-adjusted p-value
\providecommand{\rhmcsChQBmRawMin}{}\renewcommand{\rhmcsChQBmRawMin}{0.89}  % smallest single-test bootstrap p-value
\providecommand{\rhmcsChQBmRej}{}\renewcommand{\rhmcsChQBmRej}{none}  % the alternatives StepM rejects at 5 % (none: none)
\providecommand{\rhmcsChQBmNRejTwo}{}\renewcommand{\rhmcsChQBmNRejTwo}{4}  % StepM rejections at FWER 5 %, two-sided
\providecommand{\rhmcsChQBmNRejTwoWord}{}\renewcommand{\rhmcsChQBmNRejTwoWord}{four}  % ... as a word
\providecommand{\rhmcsChQBmNRawTwo}{}\renewcommand{\rhmcsChQBmNRawTwo}{4}  % single-test bootstrap p < 0.05, two-sided
\providecommand{\rhmcsChQBmRejTwo}{}\renewcommand{\rhmcsChQBmRejTwo}{ridge regression on the full predictor set, the lasso on the full predictor set, the elastic net on the full predictor set, and the random forest on the full predictor set}  % the alternatives the two-sided StepM rejects at 5 % (none: none)
\providecommand{\rhmcsChQBmPRidge}{}\renewcommand{\rhmcsChQBmPRidge}{1.00}  % StepM-adjusted p-value (one-sided), Ridge
\providecommand{\rhmcsChQBmRawRidge}{}\renewcommand{\rhmcsChQBmRawRidge}{1.00}  % single-test bootstrap p-value (one-sided), Ridge
\providecommand{\rhmcsChQBmPTwoRidge}{}\renewcommand{\rhmcsChQBmPTwoRidge}{0.008}  % StepM-adjusted p-value (two-sided), Ridge
\providecommand{\rhmcsChQBmRawTwoRidge}{}\renewcommand{\rhmcsChQBmRawTwoRidge}{0.003}  % single-test bootstrap p-value (two-sided), Ridge
\providecommand{\rhmcsChQBmPLasso}{}\renewcommand{\rhmcsChQBmPLasso}{1.00}  % StepM-adjusted p-value (one-sided), Lasso
\providecommand{\rhmcsChQBmRawLasso}{}\renewcommand{\rhmcsChQBmRawLasso}{1.00}  % single-test bootstrap p-value (one-sided), Lasso
\providecommand{\rhmcsChQBmPTwoLasso}{}\renewcommand{\rhmcsChQBmPTwoLasso}{0.004}  % StepM-adjusted p-value (two-sided), Lasso
\providecommand{\rhmcsChQBmRawTwoLasso}{}\renewcommand{\rhmcsChQBmRawTwoLasso}{0.001}  % single-test bootstrap p-value (two-sided), Lasso
\providecommand{\rhmcsChQBmPEnet}{}\renewcommand{\rhmcsChQBmPEnet}{1.00}  % StepM-adjusted p-value (one-sided), Elastic net
\providecommand{\rhmcsChQBmRawEnet}{}\renewcommand{\rhmcsChQBmRawEnet}{1.00}  % single-test bootstrap p-value (one-sided), Elastic net
\providecommand{\rhmcsChQBmPTwoEnet}{}\renewcommand{\rhmcsChQBmPTwoEnet}{0.013}  % StepM-adjusted p-value (two-sided), Elastic net
\providecommand{\rhmcsChQBmRawTwoEnet}{}\renewcommand{\rhmcsChQBmRawTwoEnet}{0.005}  % single-test bootstrap p-value (two-sided), Elastic net
\providecommand{\rhmcsChQBmPLgbm}{}\renewcommand{\rhmcsChQBmPLgbm}{0.97}  % StepM-adjusted p-value (one-sided), LightGBM
\providecommand{\rhmcsChQBmRawLgbm}{}\renewcommand{\rhmcsChQBmRawLgbm}{0.91}  % single-test bootstrap p-value (one-sided), LightGBM
\providecommand{\rhmcsChQBmPTwoLgbm}{}\renewcommand{\rhmcsChQBmPTwoLgbm}{0.19}  % StepM-adjusted p-value (two-sided), LightGBM
\providecommand{\rhmcsChQBmRawTwoLgbm}{}\renewcommand{\rhmcsChQBmRawTwoLgbm}{0.17}  % single-test bootstrap p-value (two-sided), LightGBM
\providecommand{\rhmcsChQBmPXgb}{}\renewcommand{\rhmcsChQBmPXgb}{0.95}  % StepM-adjusted p-value (one-sided), XGBoost
\providecommand{\rhmcsChQBmRawXgb}{}\renewcommand{\rhmcsChQBmRawXgb}{0.89}  % single-test bootstrap p-value (one-sided), XGBoost
\providecommand{\rhmcsChQBmPTwoXgb}{}\renewcommand{\rhmcsChQBmPTwoXgb}{0.22}  % StepM-adjusted p-value (two-sided), XGBoost
\providecommand{\rhmcsChQBmRawTwoXgb}{}\renewcommand{\rhmcsChQBmRawTwoXgb}{0.22}  % single-test bootstrap p-value (two-sided), XGBoost
\providecommand{\rhmcsChQBmPRf}{}\renewcommand{\rhmcsChQBmPRf}{1.00}  % StepM-adjusted p-value (one-sided), Random forest
\providecommand{\rhmcsChQBmRawRf}{}\renewcommand{\rhmcsChQBmRawRf}{0.99}  % single-test bootstrap p-value (one-sided), Random forest
\providecommand{\rhmcsChQBmPTwoRf}{}\renewcommand{\rhmcsChQBmPTwoRf}{0.013}  % StepM-adjusted p-value (two-sided), Random forest
\providecommand{\rhmcsChQBmRawTwoRf}{}\renewcommand{\rhmcsChQBmRawTwoRf}{0.008}  % single-test bootstrap p-value (two-sided), Random forest
\providecommand{\rhmcsChQBmNRejLo}{}\renewcommand{\rhmcsChQBmNRejLo}{0}  % StepM rejections at 5 %, block length 5
\providecommand{\rhmcsChQBmSpaCLo}{}\renewcommand{\rhmcsChQBmSpaCLo}{0.54}  % SPA p-value (consistent), block length 5
\providecommand{\rhmcsChQBmNRejTwoLo}{}\renewcommand{\rhmcsChQBmNRejTwoLo}{4}  % ... two-sided, block length 5
\providecommand{\rhmcsChQBmPMinLo}{}\renewcommand{\rhmcsChQBmPMinLo}{0.96}  % smallest adjusted p, block length 5
\providecommand{\rhmcsChQBmPTwoMinLo}{}\renewcommand{\rhmcsChQBmPTwoMinLo}{0.003}  % ... two-sided, block length 5
\providecommand{\rhmcsChQBmNRejHi}{}\renewcommand{\rhmcsChQBmNRejHi}{0}  % StepM rejections at 5 %, block length 20
\providecommand{\rhmcsChQBmSpaCHi}{}\renewcommand{\rhmcsChQBmSpaCHi}{0.54}  % SPA p-value (consistent), block length 20
\providecommand{\rhmcsChQBmNRejTwoHi}{}\renewcommand{\rhmcsChQBmNRejTwoHi}{4}  % ... two-sided, block length 20
\providecommand{\rhmcsChQBmPMinHi}{}\renewcommand{\rhmcsChQBmPMinHi}{0.96}  % smallest adjusted p, block length 20
\providecommand{\rhmcsChQBmPTwoMinHi}{}\renewcommand{\rhmcsChQBmPTwoMinHi}{0.004}  % ... two-sided, block length 20
\providecommand{\rhmcsChQBmAllSpaC}{}\renewcommand{\rhmcsChQBmAllSpaC}{0.54}  % SPA p-value (consistent), family ChQBmAll
\providecommand{\rhmcsChQBmAllSpaL}{}\renewcommand{\rhmcsChQBmAllSpaL}{0.11}  % ... lower
\providecommand{\rhmcsChQBmAllSpaU}{}\renewcommand{\rhmcsChQBmAllSpaU}{0.78}  % ... upper
\providecommand{\rhmcsChQBmAllNRej}{}\renewcommand{\rhmcsChQBmAllNRej}{0}  % StepM rejections at FWER 5 % (of 10)
\providecommand{\rhmcsChQBmAllNRejWord}{}\renewcommand{\rhmcsChQBmAllNRejWord}{none}  % ... as a word
\providecommand{\rhmcsChQBmAllNRaw}{}\renewcommand{\rhmcsChQBmAllNRaw}{0}  % single-test bootstrap p < 0.05 (no adjustment)
\providecommand{\rhmcsChQBmAllNRawWord}{}\renewcommand{\rhmcsChQBmAllNRawWord}{none}  % ... as a word
\providecommand{\rhmcsChQBmAllNPos}{}\renewcommand{\rhmcsChQBmAllNPos}{0}  % alternatives with a positive estimate (better than the benchmark)
\providecommand{\rhmcsChQBmAllPMin}{}\renewcommand{\rhmcsChQBmAllPMin}{0.98}  % smallest StepM-adjusted p-value
\providecommand{\rhmcsChQBmAllPMax}{}\renewcommand{\rhmcsChQBmAllPMax}{1.00}  % largest StepM-adjusted p-value
\providecommand{\rhmcsChQBmAllPTwoMin}{}\renewcommand{\rhmcsChQBmAllPTwoMin}{\ensuremath{<}0.001}  % smallest two-sided StepM-adjusted p-value
\providecommand{\rhmcsChQBmAllPTwoMax}{}\renewcommand{\rhmcsChQBmAllPTwoMax}{0.22}  % largest two-sided StepM-adjusted p-value
\providecommand{\rhmcsChQBmAllRawMin}{}\renewcommand{\rhmcsChQBmAllRawMin}{0.89}  % smallest single-test bootstrap p-value
\providecommand{\rhmcsChQBmAllRej}{}\renewcommand{\rhmcsChQBmAllRej}{none}  % the alternatives StepM rejects at 5 % (none: none)
\providecommand{\rhmcsChQBmAllNRejTwo}{}\renewcommand{\rhmcsChQBmAllNRejTwo}{8}  % StepM rejections at FWER 5 %, two-sided
\providecommand{\rhmcsChQBmAllNRejTwoWord}{}\renewcommand{\rhmcsChQBmAllNRejTwoWord}{eight}  % ... as a word
\providecommand{\rhmcsChQBmAllNRawTwo}{}\renewcommand{\rhmcsChQBmAllNRawTwo}{8}  % single-test bootstrap p < 0.05, two-sided
\providecommand{\rhmcsChQBmAllRejTwo}{}\renewcommand{\rhmcsChQBmAllRejTwo}{the linear models on the HAR benchmark, ridge regression on the full predictor set, the lasso on the full predictor set, the elastic net on the full predictor set, LightGBM on the HAR benchmark, XGBoost on the HAR benchmark, the random forest on the HAR benchmark, and the random forest on the full predictor set}  % the alternatives the two-sided StepM rejects at 5 % (none: none)
\providecommand{\rhmcsChQBmAllPLinHar}{}\renewcommand{\rhmcsChQBmAllPLinHar}{1.00}  % StepM-adjusted p-value (one-sided), Linear, HAR
\providecommand{\rhmcsChQBmAllRawLinHar}{}\renewcommand{\rhmcsChQBmAllRawLinHar}{1.00}  % single-test bootstrap p-value (one-sided), Linear, HAR
\providecommand{\rhmcsChQBmAllPTwoLinHar}{}\renewcommand{\rhmcsChQBmAllPTwoLinHar}{\ensuremath{<}0.001}  % StepM-adjusted p-value (two-sided), Linear, HAR
\providecommand{\rhmcsChQBmAllRawTwoLinHar}{}\renewcommand{\rhmcsChQBmAllRawTwoLinHar}{\ensuremath{<}0.001}  % single-test bootstrap p-value (two-sided), Linear, HAR
\providecommand{\rhmcsChQBmAllPRidgeFull}{}\renewcommand{\rhmcsChQBmAllPRidgeFull}{1.00}  % StepM-adjusted p-value (one-sided), Ridge, Full
\providecommand{\rhmcsChQBmAllRawRidgeFull}{}\renewcommand{\rhmcsChQBmAllRawRidgeFull}{1.00}  % single-test bootstrap p-value (one-sided), Ridge, Full
\providecommand{\rhmcsChQBmAllPTwoRidgeFull}{}\renewcommand{\rhmcsChQBmAllPTwoRidgeFull}{0.008}  % StepM-adjusted p-value (two-sided), Ridge, Full
\providecommand{\rhmcsChQBmAllRawTwoRidgeFull}{}\renewcommand{\rhmcsChQBmAllRawTwoRidgeFull}{0.003}  % single-test bootstrap p-value (two-sided), Ridge, Full
\providecommand{\rhmcsChQBmAllPLassoFull}{}\renewcommand{\rhmcsChQBmAllPLassoFull}{1.00}  % StepM-adjusted p-value (one-sided), Lasso, Full
\providecommand{\rhmcsChQBmAllRawLassoFull}{}\renewcommand{\rhmcsChQBmAllRawLassoFull}{1.00}  % single-test bootstrap p-value (one-sided), Lasso, Full
\providecommand{\rhmcsChQBmAllPTwoLassoFull}{}\renewcommand{\rhmcsChQBmAllPTwoLassoFull}{0.004}  % StepM-adjusted p-value (two-sided), Lasso, Full
\providecommand{\rhmcsChQBmAllRawTwoLassoFull}{}\renewcommand{\rhmcsChQBmAllRawTwoLassoFull}{0.001}  % single-test bootstrap p-value (two-sided), Lasso, Full
\providecommand{\rhmcsChQBmAllPEnetFull}{}\renewcommand{\rhmcsChQBmAllPEnetFull}{1.00}  % StepM-adjusted p-value (one-sided), Elastic net, Full
\providecommand{\rhmcsChQBmAllRawEnetFull}{}\renewcommand{\rhmcsChQBmAllRawEnetFull}{1.00}  % single-test bootstrap p-value (one-sided), Elastic net, Full
\providecommand{\rhmcsChQBmAllPTwoEnetFull}{}\renewcommand{\rhmcsChQBmAllPTwoEnetFull}{0.013}  % StepM-adjusted p-value (two-sided), Elastic net, Full
\providecommand{\rhmcsChQBmAllRawTwoEnetFull}{}\renewcommand{\rhmcsChQBmAllRawTwoEnetFull}{0.005}  % single-test bootstrap p-value (two-sided), Elastic net, Full
\providecommand{\rhmcsChQBmAllPLgbmHar}{}\renewcommand{\rhmcsChQBmAllPLgbmHar}{1.00}  % StepM-adjusted p-value (one-sided), LightGBM, HAR
\providecommand{\rhmcsChQBmAllRawLgbmHar}{}\renewcommand{\rhmcsChQBmAllRawLgbmHar}{1.00}  % single-test bootstrap p-value (one-sided), LightGBM, HAR
\providecommand{\rhmcsChQBmAllPTwoLgbmHar}{}\renewcommand{\rhmcsChQBmAllPTwoLgbmHar}{0.003}  % StepM-adjusted p-value (two-sided), LightGBM, HAR
\providecommand{\rhmcsChQBmAllRawTwoLgbmHar}{}\renewcommand{\rhmcsChQBmAllRawTwoLgbmHar}{\ensuremath{<}0.001}  % single-test bootstrap p-value (two-sided), LightGBM, HAR
\providecommand{\rhmcsChQBmAllPLgbmFull}{}\renewcommand{\rhmcsChQBmAllPLgbmFull}{0.99}  % StepM-adjusted p-value (one-sided), LightGBM, Full
\providecommand{\rhmcsChQBmAllRawLgbmFull}{}\renewcommand{\rhmcsChQBmAllRawLgbmFull}{0.91}  % single-test bootstrap p-value (one-sided), LightGBM, Full
\providecommand{\rhmcsChQBmAllPTwoLgbmFull}{}\renewcommand{\rhmcsChQBmAllPTwoLgbmFull}{0.19}  % StepM-adjusted p-value (two-sided), LightGBM, Full
\providecommand{\rhmcsChQBmAllRawTwoLgbmFull}{}\renewcommand{\rhmcsChQBmAllRawTwoLgbmFull}{0.17}  % single-test bootstrap p-value (two-sided), LightGBM, Full
\providecommand{\rhmcsChQBmAllPXgbHar}{}\renewcommand{\rhmcsChQBmAllPXgbHar}{1.00}  % StepM-adjusted p-value (one-sided), XGBoost, HAR
\providecommand{\rhmcsChQBmAllRawXgbHar}{}\renewcommand{\rhmcsChQBmAllRawXgbHar}{1.00}  % single-test bootstrap p-value (one-sided), XGBoost, HAR
\providecommand{\rhmcsChQBmAllPTwoXgbHar}{}\renewcommand{\rhmcsChQBmAllPTwoXgbHar}{0.002}  % StepM-adjusted p-value (two-sided), XGBoost, HAR
\providecommand{\rhmcsChQBmAllRawTwoXgbHar}{}\renewcommand{\rhmcsChQBmAllRawTwoXgbHar}{\ensuremath{<}0.001}  % single-test bootstrap p-value (two-sided), XGBoost, HAR
\providecommand{\rhmcsChQBmAllPXgbFull}{}\renewcommand{\rhmcsChQBmAllPXgbFull}{0.98}  % StepM-adjusted p-value (one-sided), XGBoost, Full
\providecommand{\rhmcsChQBmAllRawXgbFull}{}\renewcommand{\rhmcsChQBmAllRawXgbFull}{0.89}  % single-test bootstrap p-value (one-sided), XGBoost, Full
\providecommand{\rhmcsChQBmAllPTwoXgbFull}{}\renewcommand{\rhmcsChQBmAllPTwoXgbFull}{0.22}  % StepM-adjusted p-value (two-sided), XGBoost, Full
\providecommand{\rhmcsChQBmAllRawTwoXgbFull}{}\renewcommand{\rhmcsChQBmAllRawTwoXgbFull}{0.22}  % single-test bootstrap p-value (two-sided), XGBoost, Full
\providecommand{\rhmcsChQBmAllPRfHar}{}\renewcommand{\rhmcsChQBmAllPRfHar}{1.00}  % StepM-adjusted p-value (one-sided), Random forest, HAR
\providecommand{\rhmcsChQBmAllRawRfHar}{}\renewcommand{\rhmcsChQBmAllRawRfHar}{1.00}  % single-test bootstrap p-value (one-sided), Random forest, HAR
\providecommand{\rhmcsChQBmAllPTwoRfHar}{}\renewcommand{\rhmcsChQBmAllPTwoRfHar}{\ensuremath{<}0.001}  % StepM-adjusted p-value (two-sided), Random forest, HAR
\providecommand{\rhmcsChQBmAllRawTwoRfHar}{}\renewcommand{\rhmcsChQBmAllRawTwoRfHar}{\ensuremath{<}0.001}  % single-test bootstrap p-value (two-sided), Random forest, HAR
\providecommand{\rhmcsChQBmAllPRfFull}{}\renewcommand{\rhmcsChQBmAllPRfFull}{1.00}  % StepM-adjusted p-value (one-sided), Random forest, Full
\providecommand{\rhmcsChQBmAllRawRfFull}{}\renewcommand{\rhmcsChQBmAllRawRfFull}{0.99}  % single-test bootstrap p-value (one-sided), Random forest, Full
\providecommand{\rhmcsChQBmAllPTwoRfFull}{}\renewcommand{\rhmcsChQBmAllPTwoRfFull}{0.013}  % StepM-adjusted p-value (two-sided), Random forest, Full
\providecommand{\rhmcsChQBmAllRawTwoRfFull}{}\renewcommand{\rhmcsChQBmAllRawTwoRfFull}{0.008}  % single-test bootstrap p-value (two-sided), Random forest, Full
\providecommand{\rhmcsChQBmAllNRejLo}{}\renewcommand{\rhmcsChQBmAllNRejLo}{0}  % StepM rejections at 5 %, block length 5
\providecommand{\rhmcsChQBmAllSpaCLo}{}\renewcommand{\rhmcsChQBmAllSpaCLo}{0.54}  % SPA p-value (consistent), block length 5
\providecommand{\rhmcsChQBmAllNRejTwoLo}{}\renewcommand{\rhmcsChQBmAllNRejTwoLo}{8}  % ... two-sided, block length 5
\providecommand{\rhmcsChQBmAllPMinLo}{}\renewcommand{\rhmcsChQBmAllPMinLo}{0.98}  % smallest adjusted p, block length 5
\providecommand{\rhmcsChQBmAllPTwoMinLo}{}\renewcommand{\rhmcsChQBmAllPTwoMinLo}{\ensuremath{<}0.001}  % ... two-sided, block length 5
\providecommand{\rhmcsChQBmAllNRejHi}{}\renewcommand{\rhmcsChQBmAllNRejHi}{0}  % StepM rejections at 5 %, block length 20
\providecommand{\rhmcsChQBmAllSpaCHi}{}\renewcommand{\rhmcsChQBmAllSpaCHi}{0.54}  % SPA p-value (consistent), block length 20
\providecommand{\rhmcsChQBmAllNRejTwoHi}{}\renewcommand{\rhmcsChQBmAllNRejTwoHi}{8}  % ... two-sided, block length 20
\providecommand{\rhmcsChQBmAllPMinHi}{}\renewcommand{\rhmcsChQBmAllPMinHi}{0.99}  % smallest adjusted p, block length 20
\providecommand{\rhmcsChQBmAllPTwoMinHi}{}\renewcommand{\rhmcsChQBmAllPTwoMinHi}{\ensuremath{<}0.001}  % ... two-sided, block length 20
\providecommand{\rhmcsChMBmSpaC}{}\renewcommand{\rhmcsChMBmSpaC}{0.66}  % SPA p-value (consistent), family ChMBm
\providecommand{\rhmcsChMBmSpaL}{}\renewcommand{\rhmcsChMBmSpaL}{0.36}  % ... lower
\providecommand{\rhmcsChMBmSpaU}{}\renewcommand{\rhmcsChMBmSpaU}{0.66}  % ... upper
\providecommand{\rhmcsChMBmNRej}{}\renewcommand{\rhmcsChMBmNRej}{0}  % StepM rejections at FWER 5 % (of 6)
\providecommand{\rhmcsChMBmNRejWord}{}\renewcommand{\rhmcsChMBmNRejWord}{none}  % ... as a word
\providecommand{\rhmcsChMBmNRaw}{}\renewcommand{\rhmcsChMBmNRaw}{0}  % single-test bootstrap p < 0.05 (no adjustment)
\providecommand{\rhmcsChMBmNRawWord}{}\renewcommand{\rhmcsChMBmNRawWord}{none}  % ... as a word
\providecommand{\rhmcsChMBmNPos}{}\renewcommand{\rhmcsChMBmNPos}{1}  % alternatives with a positive estimate (better than the benchmark)
\providecommand{\rhmcsChMBmPMin}{}\renewcommand{\rhmcsChMBmPMin}{0.66}  % smallest StepM-adjusted p-value
\providecommand{\rhmcsChMBmPMax}{}\renewcommand{\rhmcsChMBmPMax}{0.88}  % largest StepM-adjusted p-value
\providecommand{\rhmcsChMBmPTwoMin}{}\renewcommand{\rhmcsChMBmPTwoMin}{0.51}  % smallest two-sided StepM-adjusted p-value
\providecommand{\rhmcsChMBmPTwoMax}{}\renewcommand{\rhmcsChMBmPTwoMax}{0.77}  % largest two-sided StepM-adjusted p-value
\providecommand{\rhmcsChMBmRawMin}{}\renewcommand{\rhmcsChMBmRawMin}{0.36}  % smallest single-test bootstrap p-value
\providecommand{\rhmcsChMBmRej}{}\renewcommand{\rhmcsChMBmRej}{none}  % the alternatives StepM rejects at 5 % (none: none)
\providecommand{\rhmcsChMBmNRejTwo}{}\renewcommand{\rhmcsChMBmNRejTwo}{0}  % StepM rejections at FWER 5 %, two-sided
\providecommand{\rhmcsChMBmNRejTwoWord}{}\renewcommand{\rhmcsChMBmNRejTwoWord}{none}  % ... as a word
\providecommand{\rhmcsChMBmNRawTwo}{}\renewcommand{\rhmcsChMBmNRawTwo}{0}  % single-test bootstrap p < 0.05, two-sided
\providecommand{\rhmcsChMBmRejTwo}{}\renewcommand{\rhmcsChMBmRejTwo}{none}  % the alternatives the two-sided StepM rejects at 5 % (none: none)
\providecommand{\rhmcsChMBmPRidge}{}\renewcommand{\rhmcsChMBmPRidge}{0.88}  % StepM-adjusted p-value (one-sided), Ridge
\providecommand{\rhmcsChMBmRawRidge}{}\renewcommand{\rhmcsChMBmRawRidge}{0.87}  % single-test bootstrap p-value (one-sided), Ridge
\providecommand{\rhmcsChMBmPTwoRidge}{}\renewcommand{\rhmcsChMBmPTwoRidge}{0.53}  % StepM-adjusted p-value (two-sided), Ridge
\providecommand{\rhmcsChMBmRawTwoRidge}{}\renewcommand{\rhmcsChMBmRawTwoRidge}{0.29}  % single-test bootstrap p-value (two-sided), Ridge
\providecommand{\rhmcsChMBmPLasso}{}\renewcommand{\rhmcsChMBmPLasso}{0.88}  % StepM-adjusted p-value (one-sided), Lasso
\providecommand{\rhmcsChMBmRawLasso}{}\renewcommand{\rhmcsChMBmRawLasso}{0.87}  % single-test bootstrap p-value (one-sided), Lasso
\providecommand{\rhmcsChMBmPTwoLasso}{}\renewcommand{\rhmcsChMBmPTwoLasso}{0.51}  % StepM-adjusted p-value (two-sided), Lasso
\providecommand{\rhmcsChMBmRawTwoLasso}{}\renewcommand{\rhmcsChMBmRawTwoLasso}{0.28}  % single-test bootstrap p-value (two-sided), Lasso
\providecommand{\rhmcsChMBmPEnet}{}\renewcommand{\rhmcsChMBmPEnet}{0.88}  % StepM-adjusted p-value (one-sided), Elastic net
\providecommand{\rhmcsChMBmRawEnet}{}\renewcommand{\rhmcsChMBmRawEnet}{0.87}  % single-test bootstrap p-value (one-sided), Elastic net
\providecommand{\rhmcsChMBmPTwoEnet}{}\renewcommand{\rhmcsChMBmPTwoEnet}{0.53}  % StepM-adjusted p-value (two-sided), Elastic net
\providecommand{\rhmcsChMBmRawTwoEnet}{}\renewcommand{\rhmcsChMBmRawTwoEnet}{0.28}  % single-test bootstrap p-value (two-sided), Elastic net
\providecommand{\rhmcsChMBmPLgbm}{}\renewcommand{\rhmcsChMBmPLgbm}{0.88}  % StepM-adjusted p-value (one-sided), LightGBM
\providecommand{\rhmcsChMBmRawLgbm}{}\renewcommand{\rhmcsChMBmRawLgbm}{0.82}  % single-test bootstrap p-value (one-sided), LightGBM
\providecommand{\rhmcsChMBmPTwoLgbm}{}\renewcommand{\rhmcsChMBmPTwoLgbm}{0.57}  % StepM-adjusted p-value (two-sided), LightGBM
\providecommand{\rhmcsChMBmRawTwoLgbm}{}\renewcommand{\rhmcsChMBmRawTwoLgbm}{0.36}  % single-test bootstrap p-value (two-sided), LightGBM
\providecommand{\rhmcsChMBmPXgb}{}\renewcommand{\rhmcsChMBmPXgb}{0.87}  % StepM-adjusted p-value (one-sided), XGBoost
\providecommand{\rhmcsChMBmRawXgb}{}\renewcommand{\rhmcsChMBmRawXgb}{0.76}  % single-test bootstrap p-value (one-sided), XGBoost
\providecommand{\rhmcsChMBmPTwoXgb}{}\renewcommand{\rhmcsChMBmPTwoXgb}{0.75}  % StepM-adjusted p-value (two-sided), XGBoost
\providecommand{\rhmcsChMBmRawTwoXgb}{}\renewcommand{\rhmcsChMBmRawTwoXgb}{0.55}  % single-test bootstrap p-value (two-sided), XGBoost
\providecommand{\rhmcsChMBmPRf}{}\renewcommand{\rhmcsChMBmPRf}{0.66}  % StepM-adjusted p-value (one-sided), Random forest
\providecommand{\rhmcsChMBmRawRf}{}\renewcommand{\rhmcsChMBmRawRf}{0.36}  % single-test bootstrap p-value (one-sided), Random forest
\providecommand{\rhmcsChMBmPTwoRf}{}\renewcommand{\rhmcsChMBmPTwoRf}{0.77}  % StepM-adjusted p-value (two-sided), Random forest
\providecommand{\rhmcsChMBmRawTwoRf}{}\renewcommand{\rhmcsChMBmRawTwoRf}{0.77}  % single-test bootstrap p-value (two-sided), Random forest
\providecommand{\rhmcsChMBmNRejLo}{}\renewcommand{\rhmcsChMBmNRejLo}{0}  % StepM rejections at 5 %, block length 5
\providecommand{\rhmcsChMBmSpaCLo}{}\renewcommand{\rhmcsChMBmSpaCLo}{0.74}  % SPA p-value (consistent), block length 5
\providecommand{\rhmcsChMBmNRejTwoLo}{}\renewcommand{\rhmcsChMBmNRejTwoLo}{0}  % ... two-sided, block length 5
\providecommand{\rhmcsChMBmPMinLo}{}\renewcommand{\rhmcsChMBmPMinLo}{0.74}  % smallest adjusted p, block length 5
\providecommand{\rhmcsChMBmPTwoMinLo}{}\renewcommand{\rhmcsChMBmPTwoMinLo}{0.34}  % ... two-sided, block length 5
\providecommand{\rhmcsChMBmNRejHi}{}\renewcommand{\rhmcsChMBmNRejHi}{0}  % StepM rejections at 5 %, block length 20
\providecommand{\rhmcsChMBmSpaCHi}{}\renewcommand{\rhmcsChMBmSpaCHi}{0.62}  % SPA p-value (consistent), block length 20
\providecommand{\rhmcsChMBmNRejTwoHi}{}\renewcommand{\rhmcsChMBmNRejTwoHi}{0}  % ... two-sided, block length 20
\providecommand{\rhmcsChMBmPMinHi}{}\renewcommand{\rhmcsChMBmPMinHi}{0.62}  % smallest adjusted p, block length 20
\providecommand{\rhmcsChMBmPTwoMinHi}{}\renewcommand{\rhmcsChMBmPTwoMinHi}{0.65}  % ... two-sided, block length 20
\providecommand{\rhmcsChMBmAllSpaC}{}\renewcommand{\rhmcsChMBmAllSpaC}{0.253}  % SPA p-value (consistent), family ChMBmAll
\providecommand{\rhmcsChMBmAllSpaL}{}\renewcommand{\rhmcsChMBmAllSpaL}{0.23}  % ... lower
\providecommand{\rhmcsChMBmAllSpaU}{}\renewcommand{\rhmcsChMBmAllSpaU}{0.253}  % ... upper
\providecommand{\rhmcsChMBmAllNRej}{}\renewcommand{\rhmcsChMBmAllNRej}{0}  % StepM rejections at FWER 5 % (of 10)
\providecommand{\rhmcsChMBmAllNRejWord}{}\renewcommand{\rhmcsChMBmAllNRejWord}{none}  % ... as a word
\providecommand{\rhmcsChMBmAllNRaw}{}\renewcommand{\rhmcsChMBmAllNRaw}{0}  % single-test bootstrap p < 0.05 (no adjustment)
\providecommand{\rhmcsChMBmAllNRawWord}{}\renewcommand{\rhmcsChMBmAllNRawWord}{none}  % ... as a word
\providecommand{\rhmcsChMBmAllNPos}{}\renewcommand{\rhmcsChMBmAllNPos}{5}  % alternatives with a positive estimate (better than the benchmark)
\providecommand{\rhmcsChMBmAllPMin}{}\renewcommand{\rhmcsChMBmAllPMin}{0.253}  % smallest StepM-adjusted p-value
\providecommand{\rhmcsChMBmAllPMax}{}\renewcommand{\rhmcsChMBmAllPMax}{0.88}  % largest StepM-adjusted p-value
\providecommand{\rhmcsChMBmAllPTwoMin}{}\renewcommand{\rhmcsChMBmAllPTwoMin}{0.48}  % smallest two-sided StepM-adjusted p-value
\providecommand{\rhmcsChMBmAllPTwoMax}{}\renewcommand{\rhmcsChMBmAllPTwoMax}{0.90}  % largest two-sided StepM-adjusted p-value
\providecommand{\rhmcsChMBmAllRawMin}{}\renewcommand{\rhmcsChMBmAllRawMin}{0.11}  % smallest single-test bootstrap p-value
\providecommand{\rhmcsChMBmAllRej}{}\renewcommand{\rhmcsChMBmAllRej}{none}  % the alternatives StepM rejects at 5 % (none: none)
\providecommand{\rhmcsChMBmAllNRejTwo}{}\renewcommand{\rhmcsChMBmAllNRejTwo}{0}  % StepM rejections at FWER 5 %, two-sided
\providecommand{\rhmcsChMBmAllNRejTwoWord}{}\renewcommand{\rhmcsChMBmAllNRejTwoWord}{none}  % ... as a word
\providecommand{\rhmcsChMBmAllNRawTwo}{}\renewcommand{\rhmcsChMBmAllNRawTwo}{0}  % single-test bootstrap p < 0.05, two-sided
\providecommand{\rhmcsChMBmAllRejTwo}{}\renewcommand{\rhmcsChMBmAllRejTwo}{none}  % the alternatives the two-sided StepM rejects at 5 % (none: none)
\providecommand{\rhmcsChMBmAllPLinHar}{}\renewcommand{\rhmcsChMBmAllPLinHar}{0.72}  % StepM-adjusted p-value (one-sided), Linear, HAR
\providecommand{\rhmcsChMBmAllRawLinHar}{}\renewcommand{\rhmcsChMBmAllRawLinHar}{0.41}  % single-test bootstrap p-value (one-sided), Linear, HAR
\providecommand{\rhmcsChMBmAllPTwoLinHar}{}\renewcommand{\rhmcsChMBmAllPTwoLinHar}{0.90}  % StepM-adjusted p-value (two-sided), Linear, HAR
\providecommand{\rhmcsChMBmAllRawTwoLinHar}{}\renewcommand{\rhmcsChMBmAllRawTwoLinHar}{0.86}  % single-test bootstrap p-value (two-sided), Linear, HAR
\providecommand{\rhmcsChMBmAllPRidgeFull}{}\renewcommand{\rhmcsChMBmAllPRidgeFull}{0.88}  % StepM-adjusted p-value (one-sided), Ridge, Full
\providecommand{\rhmcsChMBmAllRawRidgeFull}{}\renewcommand{\rhmcsChMBmAllRawRidgeFull}{0.87}  % single-test bootstrap p-value (one-sided), Ridge, Full
\providecommand{\rhmcsChMBmAllPTwoRidgeFull}{}\renewcommand{\rhmcsChMBmAllPTwoRidgeFull}{0.63}  % StepM-adjusted p-value (two-sided), Ridge, Full
\providecommand{\rhmcsChMBmAllRawTwoRidgeFull}{}\renewcommand{\rhmcsChMBmAllRawTwoRidgeFull}{0.29}  % single-test bootstrap p-value (two-sided), Ridge, Full
\providecommand{\rhmcsChMBmAllPLassoFull}{}\renewcommand{\rhmcsChMBmAllPLassoFull}{0.88}  % StepM-adjusted p-value (one-sided), Lasso, Full
\providecommand{\rhmcsChMBmAllRawLassoFull}{}\renewcommand{\rhmcsChMBmAllRawLassoFull}{0.87}  % single-test bootstrap p-value (one-sided), Lasso, Full
\providecommand{\rhmcsChMBmAllPTwoLassoFull}{}\renewcommand{\rhmcsChMBmAllPTwoLassoFull}{0.61}  % StepM-adjusted p-value (two-sided), Lasso, Full
\providecommand{\rhmcsChMBmAllRawTwoLassoFull}{}\renewcommand{\rhmcsChMBmAllRawTwoLassoFull}{0.28}  % single-test bootstrap p-value (two-sided), Lasso, Full
\providecommand{\rhmcsChMBmAllPEnetFull}{}\renewcommand{\rhmcsChMBmAllPEnetFull}{0.88}  % StepM-adjusted p-value (one-sided), Elastic net, Full
\providecommand{\rhmcsChMBmAllRawEnetFull}{}\renewcommand{\rhmcsChMBmAllRawEnetFull}{0.87}  % single-test bootstrap p-value (one-sided), Elastic net, Full
\providecommand{\rhmcsChMBmAllPTwoEnetFull}{}\renewcommand{\rhmcsChMBmAllPTwoEnetFull}{0.64}  % StepM-adjusted p-value (two-sided), Elastic net, Full
\providecommand{\rhmcsChMBmAllRawTwoEnetFull}{}\renewcommand{\rhmcsChMBmAllRawTwoEnetFull}{0.28}  % single-test bootstrap p-value (two-sided), Elastic net, Full
\providecommand{\rhmcsChMBmAllPLgbmHar}{}\renewcommand{\rhmcsChMBmAllPLgbmHar}{0.48}  % StepM-adjusted p-value (one-sided), LightGBM, HAR
\providecommand{\rhmcsChMBmAllRawLgbmHar}{}\renewcommand{\rhmcsChMBmAllRawLgbmHar}{0.15}  % single-test bootstrap p-value (one-sided), LightGBM, HAR
\providecommand{\rhmcsChMBmAllPTwoLgbmHar}{}\renewcommand{\rhmcsChMBmAllPTwoLgbmHar}{0.68}  % StepM-adjusted p-value (two-sided), LightGBM, HAR
\providecommand{\rhmcsChMBmAllRawTwoLgbmHar}{}\renewcommand{\rhmcsChMBmAllRawTwoLgbmHar}{0.30}  % single-test bootstrap p-value (two-sided), LightGBM, HAR
\providecommand{\rhmcsChMBmAllPLgbmFull}{}\renewcommand{\rhmcsChMBmAllPLgbmFull}{0.88}  % StepM-adjusted p-value (one-sided), LightGBM, Full
\providecommand{\rhmcsChMBmAllRawLgbmFull}{}\renewcommand{\rhmcsChMBmAllRawLgbmFull}{0.82}  % single-test bootstrap p-value (one-sided), LightGBM, Full
\providecommand{\rhmcsChMBmAllPTwoLgbmFull}{}\renewcommand{\rhmcsChMBmAllPTwoLgbmFull}{0.68}  % StepM-adjusted p-value (two-sided), LightGBM, Full
\providecommand{\rhmcsChMBmAllRawTwoLgbmFull}{}\renewcommand{\rhmcsChMBmAllRawTwoLgbmFull}{0.36}  % single-test bootstrap p-value (two-sided), LightGBM, Full
\providecommand{\rhmcsChMBmAllPXgbHar}{}\renewcommand{\rhmcsChMBmAllPXgbHar}{0.255}  % StepM-adjusted p-value (one-sided), XGBoost, HAR
\providecommand{\rhmcsChMBmAllRawXgbHar}{}\renewcommand{\rhmcsChMBmAllRawXgbHar}{0.11}  % single-test bootstrap p-value (one-sided), XGBoost, HAR
\providecommand{\rhmcsChMBmAllPTwoXgbHar}{}\renewcommand{\rhmcsChMBmAllPTwoXgbHar}{0.49}  % StepM-adjusted p-value (two-sided), XGBoost, HAR
\providecommand{\rhmcsChMBmAllRawTwoXgbHar}{}\renewcommand{\rhmcsChMBmAllRawTwoXgbHar}{0.21}  % single-test bootstrap p-value (two-sided), XGBoost, HAR
\providecommand{\rhmcsChMBmAllPXgbFull}{}\renewcommand{\rhmcsChMBmAllPXgbFull}{0.87}  % StepM-adjusted p-value (one-sided), XGBoost, Full
\providecommand{\rhmcsChMBmAllRawXgbFull}{}\renewcommand{\rhmcsChMBmAllRawXgbFull}{0.76}  % single-test bootstrap p-value (one-sided), XGBoost, Full
\providecommand{\rhmcsChMBmAllPTwoXgbFull}{}\renewcommand{\rhmcsChMBmAllPTwoXgbFull}{0.82}  % StepM-adjusted p-value (two-sided), XGBoost, Full
\providecommand{\rhmcsChMBmAllRawTwoXgbFull}{}\renewcommand{\rhmcsChMBmAllRawTwoXgbFull}{0.55}  % single-test bootstrap p-value (two-sided), XGBoost, Full
\providecommand{\rhmcsChMBmAllPRfHar}{}\renewcommand{\rhmcsChMBmAllPRfHar}{0.253}  % StepM-adjusted p-value (one-sided), Random forest, HAR
\providecommand{\rhmcsChMBmAllRawRfHar}{}\renewcommand{\rhmcsChMBmAllRawRfHar}{0.11}  % single-test bootstrap p-value (one-sided), Random forest, HAR
\providecommand{\rhmcsChMBmAllPTwoRfHar}{}\renewcommand{\rhmcsChMBmAllPTwoRfHar}{0.48}  % StepM-adjusted p-value (two-sided), Random forest, HAR
\providecommand{\rhmcsChMBmAllRawTwoRfHar}{}\renewcommand{\rhmcsChMBmAllRawTwoRfHar}{0.21}  % single-test bootstrap p-value (two-sided), Random forest, HAR
\providecommand{\rhmcsChMBmAllPRfFull}{}\renewcommand{\rhmcsChMBmAllPRfFull}{0.71}  % StepM-adjusted p-value (one-sided), Random forest, Full
\providecommand{\rhmcsChMBmAllRawRfFull}{}\renewcommand{\rhmcsChMBmAllRawRfFull}{0.36}  % single-test bootstrap p-value (one-sided), Random forest, Full
\providecommand{\rhmcsChMBmAllPTwoRfFull}{}\renewcommand{\rhmcsChMBmAllPTwoRfFull}{0.90}  % StepM-adjusted p-value (two-sided), Random forest, Full
\providecommand{\rhmcsChMBmAllRawTwoRfFull}{}\renewcommand{\rhmcsChMBmAllRawTwoRfFull}{0.77}  % single-test bootstrap p-value (two-sided), Random forest, Full
\providecommand{\rhmcsChMBmAllNRejLo}{}\renewcommand{\rhmcsChMBmAllNRejLo}{0}  % StepM rejections at 5 %, block length 5
\providecommand{\rhmcsChMBmAllSpaCLo}{}\renewcommand{\rhmcsChMBmAllSpaCLo}{0.21}  % SPA p-value (consistent), block length 5
\providecommand{\rhmcsChMBmAllNRejTwoLo}{}\renewcommand{\rhmcsChMBmAllNRejTwoLo}{0}  % ... two-sided, block length 5
\providecommand{\rhmcsChMBmAllPMinLo}{}\renewcommand{\rhmcsChMBmAllPMinLo}{0.21}  % smallest adjusted p, block length 5
\providecommand{\rhmcsChMBmAllPTwoMinLo}{}\renewcommand{\rhmcsChMBmAllPTwoMinLo}{0.40}  % ... two-sided, block length 5
\providecommand{\rhmcsChMBmAllNRejHi}{}\renewcommand{\rhmcsChMBmAllNRejHi}{0}  % StepM rejections at 5 %, block length 20
\providecommand{\rhmcsChMBmAllSpaCHi}{}\renewcommand{\rhmcsChMBmAllSpaCHi}{0.39}  % SPA p-value (consistent), block length 20
\providecommand{\rhmcsChMBmAllNRejTwoHi}{}\renewcommand{\rhmcsChMBmAllNRejTwoHi}{0}  % ... two-sided, block length 20
\providecommand{\rhmcsChMBmAllPMinHi}{}\renewcommand{\rhmcsChMBmAllPMinHi}{0.39}  % smallest adjusted p, block length 20
\providecommand{\rhmcsChMBmAllPTwoMinHi}{}\renewcommand{\rhmcsChMBmAllPTwoMinHi}{0.69}  % ... two-sided, block length 20
\providecommand{\rhmcsMidShortMeanSpaC}{}\renewcommand{\rhmcsMidShortMeanSpaC}{0.026}  % SPA p-value (consistent), family MidShortMean
\providecommand{\rhmcsMidShortMeanSpaL}{}\renewcommand{\rhmcsMidShortMeanSpaL}{0.026}  % ... lower
\providecommand{\rhmcsMidShortMeanSpaU}{}\renewcommand{\rhmcsMidShortMeanSpaU}{0.026}  % ... upper
\providecommand{\rhmcsMidShortMeanNRej}{}\renewcommand{\rhmcsMidShortMeanNRej}{3}  % StepM rejections at FWER 5 % (of 6)
\providecommand{\rhmcsMidShortMeanNRejWord}{}\renewcommand{\rhmcsMidShortMeanNRejWord}{three}  % ... as a word
\providecommand{\rhmcsMidShortMeanNRaw}{}\renewcommand{\rhmcsMidShortMeanNRaw}{6}  % single-test bootstrap p < 0.05 (no adjustment)
\providecommand{\rhmcsMidShortMeanNRawWord}{}\renewcommand{\rhmcsMidShortMeanNRawWord}{six}  % ... as a word
\providecommand{\rhmcsMidShortMeanNPos}{}\renewcommand{\rhmcsMidShortMeanNPos}{6}  % alternatives with a positive estimate (better than the benchmark)
\providecommand{\rhmcsMidShortMeanPMin}{}\renewcommand{\rhmcsMidShortMeanPMin}{0.026}  % smallest StepM-adjusted p-value
\providecommand{\rhmcsMidShortMeanPMax}{}\renewcommand{\rhmcsMidShortMeanPMax}{0.051}  % largest StepM-adjusted p-value
\providecommand{\rhmcsMidShortMeanPTwoMin}{}\renewcommand{\rhmcsMidShortMeanPTwoMin}{0.034}  % smallest two-sided StepM-adjusted p-value
\providecommand{\rhmcsMidShortMeanPTwoMax}{}\renewcommand{\rhmcsMidShortMeanPTwoMax}{0.083}  % largest two-sided StepM-adjusted p-value
\providecommand{\rhmcsMidShortMeanRawMin}{}\renewcommand{\rhmcsMidShortMeanRawMin}{0.015}  % smallest single-test bootstrap p-value
\providecommand{\rhmcsMidShortMeanRej}{}\renewcommand{\rhmcsMidShortMeanRej}{the lasso on the full predictor set, the elastic net on the full predictor set, and the random forest on the full predictor set}  % the alternatives StepM rejects at 5 % (none: none)
\providecommand{\rhmcsMidShortMeanNRejTwo}{}\renewcommand{\rhmcsMidShortMeanNRejTwo}{2}  % StepM rejections at FWER 5 %, two-sided
\providecommand{\rhmcsMidShortMeanNRejTwoWord}{}\renewcommand{\rhmcsMidShortMeanNRejTwoWord}{two}  % ... as a word
\providecommand{\rhmcsMidShortMeanNRawTwo}{}\renewcommand{\rhmcsMidShortMeanNRawTwo}{3}  % single-test bootstrap p < 0.05, two-sided
\providecommand{\rhmcsMidShortMeanRejTwo}{}\renewcommand{\rhmcsMidShortMeanRejTwo}{the elastic net on the full predictor set and the random forest on the full predictor set}  % the alternatives the two-sided StepM rejects at 5 % (none: none)
\providecommand{\rhmcsMidShortMeanPRidge}{}\renewcommand{\rhmcsMidShortMeanPRidge}{0.051}  % StepM-adjusted p-value (one-sided), Ridge
\providecommand{\rhmcsMidShortMeanRawRidge}{}\renewcommand{\rhmcsMidShortMeanRawRidge}{0.037}  % single-test bootstrap p-value (one-sided), Ridge
\providecommand{\rhmcsMidShortMeanPTwoRidge}{}\renewcommand{\rhmcsMidShortMeanPTwoRidge}{0.083}  % StepM-adjusted p-value (two-sided), Ridge
\providecommand{\rhmcsMidShortMeanRawTwoRidge}{}\renewcommand{\rhmcsMidShortMeanRawTwoRidge}{0.057}  % single-test bootstrap p-value (two-sided), Ridge
\providecommand{\rhmcsMidShortMeanPLasso}{}\renewcommand{\rhmcsMidShortMeanPLasso}{0.038}  % StepM-adjusted p-value (one-sided), Lasso
\providecommand{\rhmcsMidShortMeanRawLasso}{}\renewcommand{\rhmcsMidShortMeanRawLasso}{0.025}  % single-test bootstrap p-value (one-sided), Lasso
\providecommand{\rhmcsMidShortMeanPTwoLasso}{}\renewcommand{\rhmcsMidShortMeanPTwoLasso}{0.059}  % StepM-adjusted p-value (two-sided), Lasso
\providecommand{\rhmcsMidShortMeanRawTwoLasso}{}\renewcommand{\rhmcsMidShortMeanRawTwoLasso}{0.036}  % single-test bootstrap p-value (two-sided), Lasso
\providecommand{\rhmcsMidShortMeanPEnet}{}\renewcommand{\rhmcsMidShortMeanPEnet}{0.031}  % StepM-adjusted p-value (one-sided), Elastic net
\providecommand{\rhmcsMidShortMeanRawEnet}{}\renewcommand{\rhmcsMidShortMeanRawEnet}{0.020}  % single-test bootstrap p-value (one-sided), Elastic net
\providecommand{\rhmcsMidShortMeanPTwoEnet}{}\renewcommand{\rhmcsMidShortMeanPTwoEnet}{0.046}  % StepM-adjusted p-value (two-sided), Elastic net
\providecommand{\rhmcsMidShortMeanRawTwoEnet}{}\renewcommand{\rhmcsMidShortMeanRawTwoEnet}{0.025}  % single-test bootstrap p-value (two-sided), Elastic net
\providecommand{\rhmcsMidShortMeanPLgbm}{}\renewcommand{\rhmcsMidShortMeanPLgbm}{0.051}  % StepM-adjusted p-value (one-sided), LightGBM
\providecommand{\rhmcsMidShortMeanRawLgbm}{}\renewcommand{\rhmcsMidShortMeanRawLgbm}{0.041}  % single-test bootstrap p-value (one-sided), LightGBM
\providecommand{\rhmcsMidShortMeanPTwoLgbm}{}\renewcommand{\rhmcsMidShortMeanPTwoLgbm}{0.083}  % StepM-adjusted p-value (two-sided), LightGBM
\providecommand{\rhmcsMidShortMeanRawTwoLgbm}{}\renewcommand{\rhmcsMidShortMeanRawTwoLgbm}{0.067}  % single-test bootstrap p-value (two-sided), LightGBM
\providecommand{\rhmcsMidShortMeanPXgb}{}\renewcommand{\rhmcsMidShortMeanPXgb}{0.051}  % StepM-adjusted p-value (one-sided), XGBoost
\providecommand{\rhmcsMidShortMeanRawXgb}{}\renewcommand{\rhmcsMidShortMeanRawXgb}{0.039}  % single-test bootstrap p-value (one-sided), XGBoost
\providecommand{\rhmcsMidShortMeanPTwoXgb}{}\renewcommand{\rhmcsMidShortMeanPTwoXgb}{0.083}  % StepM-adjusted p-value (two-sided), XGBoost
\providecommand{\rhmcsMidShortMeanRawTwoXgb}{}\renewcommand{\rhmcsMidShortMeanRawTwoXgb}{0.065}  % single-test bootstrap p-value (two-sided), XGBoost
\providecommand{\rhmcsMidShortMeanPRf}{}\renewcommand{\rhmcsMidShortMeanPRf}{0.026}  % StepM-adjusted p-value (one-sided), Random forest
\providecommand{\rhmcsMidShortMeanRawRf}{}\renewcommand{\rhmcsMidShortMeanRawRf}{0.015}  % single-test bootstrap p-value (one-sided), Random forest
\providecommand{\rhmcsMidShortMeanPTwoRf}{}\renewcommand{\rhmcsMidShortMeanPTwoRf}{0.034}  % StepM-adjusted p-value (two-sided), Random forest
\providecommand{\rhmcsMidShortMeanRawTwoRf}{}\renewcommand{\rhmcsMidShortMeanRawTwoRf}{0.020}  % single-test bootstrap p-value (two-sided), Random forest
\providecommand{\rhmcsMidShortMeanNRejLo}{}\renewcommand{\rhmcsMidShortMeanNRejLo}{3}  % StepM rejections at 5 %, block length 5
\providecommand{\rhmcsMidShortMeanSpaCLo}{}\renewcommand{\rhmcsMidShortMeanSpaCLo}{0.037}  % SPA p-value (consistent), block length 5
\providecommand{\rhmcsMidShortMeanNRejTwoLo}{}\renewcommand{\rhmcsMidShortMeanNRejTwoLo}{0}  % ... two-sided, block length 5
\providecommand{\rhmcsMidShortMeanPMinLo}{}\renewcommand{\rhmcsMidShortMeanPMinLo}{0.037}  % smallest adjusted p, block length 5
\providecommand{\rhmcsMidShortMeanPTwoMinLo}{}\renewcommand{\rhmcsMidShortMeanPTwoMinLo}{0.055}  % ... two-sided, block length 5
\providecommand{\rhmcsMidShortMeanNRejHi}{}\renewcommand{\rhmcsMidShortMeanNRejHi}{3}  % StepM rejections at 5 %, block length 20
\providecommand{\rhmcsMidShortMeanSpaCHi}{}\renewcommand{\rhmcsMidShortMeanSpaCHi}{0.028}  % SPA p-value (consistent), block length 20
\providecommand{\rhmcsMidShortMeanNRejTwoHi}{}\renewcommand{\rhmcsMidShortMeanNRejTwoHi}{1}  % ... two-sided, block length 20
\providecommand{\rhmcsMidShortMeanPMinHi}{}\renewcommand{\rhmcsMidShortMeanPMinHi}{0.028}  % smallest adjusted p, block length 20
\providecommand{\rhmcsMidShortMeanPTwoMinHi}{}\renewcommand{\rhmcsMidShortMeanPTwoMinHi}{0.036}  % ... two-sided, block length 20
\providecommand{\rhmcsMidShortSrSpaC}{}\renewcommand{\rhmcsMidShortSrSpaC}{0.008}  % SPA p-value (consistent), family MidShortSr
\providecommand{\rhmcsMidShortSrSpaL}{}\renewcommand{\rhmcsMidShortSrSpaL}{0.008}  % ... lower
\providecommand{\rhmcsMidShortSrSpaU}{}\renewcommand{\rhmcsMidShortSrSpaU}{0.008}  % ... upper
\providecommand{\rhmcsMidShortSrNRej}{}\renewcommand{\rhmcsMidShortSrNRej}{6}  % StepM rejections at FWER 5 % (of 6)
\providecommand{\rhmcsMidShortSrNRejWord}{}\renewcommand{\rhmcsMidShortSrNRejWord}{six}  % ... as a word
\providecommand{\rhmcsMidShortSrNRaw}{}\renewcommand{\rhmcsMidShortSrNRaw}{6}  % single-test bootstrap p < 0.05 (no adjustment)
\providecommand{\rhmcsMidShortSrNRawWord}{}\renewcommand{\rhmcsMidShortSrNRawWord}{six}  % ... as a word
\providecommand{\rhmcsMidShortSrNPos}{}\renewcommand{\rhmcsMidShortSrNPos}{6}  % alternatives with a positive estimate (better than the benchmark)
\providecommand{\rhmcsMidShortSrPMin}{}\renewcommand{\rhmcsMidShortSrPMin}{0.008}  % smallest StepM-adjusted p-value
\providecommand{\rhmcsMidShortSrPMax}{}\renewcommand{\rhmcsMidShortSrPMax}{0.028}  % largest StepM-adjusted p-value
\providecommand{\rhmcsMidShortSrPTwoMin}{}\renewcommand{\rhmcsMidShortSrPTwoMin}{0.016}  % smallest two-sided StepM-adjusted p-value
\providecommand{\rhmcsMidShortSrPTwoMax}{}\renewcommand{\rhmcsMidShortSrPTwoMax}{0.060}  % largest two-sided StepM-adjusted p-value
\providecommand{\rhmcsMidShortSrRawMin}{}\renewcommand{\rhmcsMidShortSrRawMin}{0.003}  % smallest single-test bootstrap p-value
\providecommand{\rhmcsMidShortSrRej}{}\renewcommand{\rhmcsMidShortSrRej}{ridge regression on the full predictor set, the lasso on the full predictor set, the elastic net on the full predictor set, LightGBM on the full predictor set, XGBoost on the full predictor set, and the random forest on the full predictor set}  % the alternatives StepM rejects at 5 % (none: none)
\providecommand{\rhmcsMidShortSrNRejTwo}{}\renewcommand{\rhmcsMidShortSrNRejTwo}{3}  % StepM rejections at FWER 5 %, two-sided
\providecommand{\rhmcsMidShortSrNRejTwoWord}{}\renewcommand{\rhmcsMidShortSrNRejTwoWord}{three}  % ... as a word
\providecommand{\rhmcsMidShortSrNRawTwo}{}\renewcommand{\rhmcsMidShortSrNRawTwo}{6}  % single-test bootstrap p < 0.05, two-sided
\providecommand{\rhmcsMidShortSrRejTwo}{}\renewcommand{\rhmcsMidShortSrRejTwo}{the lasso on the full predictor set, the elastic net on the full predictor set, and the random forest on the full predictor set}  % the alternatives the two-sided StepM rejects at 5 % (none: none)
\providecommand{\rhmcsMidShortSrPRidge}{}\renewcommand{\rhmcsMidShortSrPRidge}{0.028}  % StepM-adjusted p-value (one-sided), Ridge
\providecommand{\rhmcsMidShortSrRawRidge}{}\renewcommand{\rhmcsMidShortSrRawRidge}{0.019}  % single-test bootstrap p-value (one-sided), Ridge
\providecommand{\rhmcsMidShortSrPTwoRidge}{}\renewcommand{\rhmcsMidShortSrPTwoRidge}{0.060}  % StepM-adjusted p-value (two-sided), Ridge
\providecommand{\rhmcsMidShortSrRawTwoRidge}{}\renewcommand{\rhmcsMidShortSrRawTwoRidge}{0.039}  % single-test bootstrap p-value (two-sided), Ridge
\providecommand{\rhmcsMidShortSrPLasso}{}\renewcommand{\rhmcsMidShortSrPLasso}{0.017}  % StepM-adjusted p-value (one-sided), Lasso
\providecommand{\rhmcsMidShortSrRawLasso}{}\renewcommand{\rhmcsMidShortSrRawLasso}{0.010}  % single-test bootstrap p-value (one-sided), Lasso
\providecommand{\rhmcsMidShortSrPTwoLasso}{}\renewcommand{\rhmcsMidShortSrPTwoLasso}{0.040}  % StepM-adjusted p-value (two-sided), Lasso
\providecommand{\rhmcsMidShortSrRawTwoLasso}{}\renewcommand{\rhmcsMidShortSrRawTwoLasso}{0.021}  % single-test bootstrap p-value (two-sided), Lasso
\providecommand{\rhmcsMidShortSrPEnet}{}\renewcommand{\rhmcsMidShortSrPEnet}{0.012}  % StepM-adjusted p-value (one-sided), Elastic net
\providecommand{\rhmcsMidShortSrRawEnet}{}\renewcommand{\rhmcsMidShortSrRawEnet}{0.007}  % single-test bootstrap p-value (one-sided), Elastic net
\providecommand{\rhmcsMidShortSrPTwoEnet}{}\renewcommand{\rhmcsMidShortSrPTwoEnet}{0.027}  % StepM-adjusted p-value (two-sided), Elastic net
\providecommand{\rhmcsMidShortSrRawTwoEnet}{}\renewcommand{\rhmcsMidShortSrRawTwoEnet}{0.012}  % single-test bootstrap p-value (two-sided), Elastic net
\providecommand{\rhmcsMidShortSrPLgbm}{}\renewcommand{\rhmcsMidShortSrPLgbm}{0.028}  % StepM-adjusted p-value (one-sided), LightGBM
\providecommand{\rhmcsMidShortSrRawLgbm}{}\renewcommand{\rhmcsMidShortSrRawLgbm}{0.021}  % single-test bootstrap p-value (one-sided), LightGBM
\providecommand{\rhmcsMidShortSrPTwoLgbm}{}\renewcommand{\rhmcsMidShortSrPTwoLgbm}{0.060}  % StepM-adjusted p-value (two-sided), LightGBM
\providecommand{\rhmcsMidShortSrRawTwoLgbm}{}\renewcommand{\rhmcsMidShortSrRawTwoLgbm}{0.047}  % single-test bootstrap p-value (two-sided), LightGBM
\providecommand{\rhmcsMidShortSrPXgb}{}\renewcommand{\rhmcsMidShortSrPXgb}{0.028}  % StepM-adjusted p-value (one-sided), XGBoost
\providecommand{\rhmcsMidShortSrRawXgb}{}\renewcommand{\rhmcsMidShortSrRawXgb}{0.020}  % single-test bootstrap p-value (one-sided), XGBoost
\providecommand{\rhmcsMidShortSrPTwoXgb}{}\renewcommand{\rhmcsMidShortSrPTwoXgb}{0.060}  % StepM-adjusted p-value (two-sided), XGBoost
\providecommand{\rhmcsMidShortSrRawTwoXgb}{}\renewcommand{\rhmcsMidShortSrRawTwoXgb}{0.046}  % single-test bootstrap p-value (two-sided), XGBoost
\providecommand{\rhmcsMidShortSrPRf}{}\renewcommand{\rhmcsMidShortSrPRf}{0.008}  % StepM-adjusted p-value (one-sided), Random forest
\providecommand{\rhmcsMidShortSrRawRf}{}\renewcommand{\rhmcsMidShortSrRawRf}{0.003}  % single-test bootstrap p-value (one-sided), Random forest
\providecommand{\rhmcsMidShortSrPTwoRf}{}\renewcommand{\rhmcsMidShortSrPTwoRf}{0.016}  % StepM-adjusted p-value (two-sided), Random forest
\providecommand{\rhmcsMidShortSrRawTwoRf}{}\renewcommand{\rhmcsMidShortSrRawTwoRf}{0.007}  % single-test bootstrap p-value (two-sided), Random forest
\providecommand{\rhmcsMidShortSrNRejLo}{}\renewcommand{\rhmcsMidShortSrNRejLo}{6}  % StepM rejections at 5 %, block length 5
\providecommand{\rhmcsMidShortSrSpaCLo}{}\renewcommand{\rhmcsMidShortSrSpaCLo}{0.010}  % SPA p-value (consistent), block length 5
\providecommand{\rhmcsMidShortSrNRejTwoLo}{}\renewcommand{\rhmcsMidShortSrNRejTwoLo}{3}  % ... two-sided, block length 5
\providecommand{\rhmcsMidShortSrPMinLo}{}\renewcommand{\rhmcsMidShortSrPMinLo}{0.010}  % smallest adjusted p, block length 5
\providecommand{\rhmcsMidShortSrPTwoMinLo}{}\renewcommand{\rhmcsMidShortSrPTwoMinLo}{0.027}  % ... two-sided, block length 5
\providecommand{\rhmcsMidShortSrNRejHi}{}\renewcommand{\rhmcsMidShortSrNRejHi}{6}  % StepM rejections at 5 %, block length 20
\providecommand{\rhmcsMidShortSrSpaCHi}{}\renewcommand{\rhmcsMidShortSrSpaCHi}{0.008}  % SPA p-value (consistent), block length 20
\providecommand{\rhmcsMidShortSrNRejTwoHi}{}\renewcommand{\rhmcsMidShortSrNRejTwoHi}{3}  % ... two-sided, block length 20
\providecommand{\rhmcsMidShortSrPMinHi}{}\renewcommand{\rhmcsMidShortSrPMinHi}{0.008}  % smallest adjusted p, block length 20
\providecommand{\rhmcsMidShortSrPTwoMinHi}{}\renewcommand{\rhmcsMidShortSrPTwoMinHi}{0.015}  % ... two-sided, block length 20
\providecommand{\rhmcsMidTwinMeanSpaC}{}\renewcommand{\rhmcsMidTwinMeanSpaC}{0.027}  % SPA p-value (consistent), family MidTwinMean
\providecommand{\rhmcsMidTwinMeanSpaL}{}\renewcommand{\rhmcsMidTwinMeanSpaL}{0.027}  % ... lower
\providecommand{\rhmcsMidTwinMeanSpaU}{}\renewcommand{\rhmcsMidTwinMeanSpaU}{0.027}  % ... upper
\providecommand{\rhmcsMidTwinMeanNRej}{}\renewcommand{\rhmcsMidTwinMeanNRej}{1}  % StepM rejections at FWER 5 % (of 6)
\providecommand{\rhmcsMidTwinMeanNRejWord}{}\renewcommand{\rhmcsMidTwinMeanNRejWord}{one}  % ... as a word
\providecommand{\rhmcsMidTwinMeanNRaw}{}\renewcommand{\rhmcsMidTwinMeanNRaw}{3}  % single-test bootstrap p < 0.05 (no adjustment)
\providecommand{\rhmcsMidTwinMeanNRawWord}{}\renewcommand{\rhmcsMidTwinMeanNRawWord}{three}  % ... as a word
\providecommand{\rhmcsMidTwinMeanNPos}{}\renewcommand{\rhmcsMidTwinMeanNPos}{6}  % alternatives with a positive estimate (better than the benchmark)
\providecommand{\rhmcsMidTwinMeanPMin}{}\renewcommand{\rhmcsMidTwinMeanPMin}{0.027}  % smallest StepM-adjusted p-value
\providecommand{\rhmcsMidTwinMeanPMax}{}\renewcommand{\rhmcsMidTwinMeanPMax}{0.13}  % largest StepM-adjusted p-value
\providecommand{\rhmcsMidTwinMeanPTwoMin}{}\renewcommand{\rhmcsMidTwinMeanPTwoMin}{0.042}  % smallest two-sided StepM-adjusted p-value
\providecommand{\rhmcsMidTwinMeanPTwoMax}{}\renewcommand{\rhmcsMidTwinMeanPTwoMax}{0.26}  % largest two-sided StepM-adjusted p-value
\providecommand{\rhmcsMidTwinMeanRawMin}{}\renewcommand{\rhmcsMidTwinMeanRawMin}{0.013}  % smallest single-test bootstrap p-value
\providecommand{\rhmcsMidTwinMeanRej}{}\renewcommand{\rhmcsMidTwinMeanRej}{the random forest on the full predictor set}  % the alternatives StepM rejects at 5 % (none: none)
\providecommand{\rhmcsMidTwinMeanNRejTwo}{}\renewcommand{\rhmcsMidTwinMeanNRejTwo}{1}  % StepM rejections at FWER 5 %, two-sided
\providecommand{\rhmcsMidTwinMeanNRejTwoWord}{}\renewcommand{\rhmcsMidTwinMeanNRejTwoWord}{one}  % ... as a word
\providecommand{\rhmcsMidTwinMeanNRawTwo}{}\renewcommand{\rhmcsMidTwinMeanNRawTwo}{2}  % single-test bootstrap p < 0.05, two-sided
\providecommand{\rhmcsMidTwinMeanRejTwo}{}\renewcommand{\rhmcsMidTwinMeanRejTwo}{the random forest on the full predictor set}  % the alternatives the two-sided StepM rejects at 5 % (none: none)
\providecommand{\rhmcsMidTwinMeanPRidge}{}\renewcommand{\rhmcsMidTwinMeanPRidge}{0.13}  % StepM-adjusted p-value (one-sided), Ridge
\providecommand{\rhmcsMidTwinMeanRawRidge}{}\renewcommand{\rhmcsMidTwinMeanRawRidge}{0.082}  % single-test bootstrap p-value (one-sided), Ridge
\providecommand{\rhmcsMidTwinMeanPTwoRidge}{}\renewcommand{\rhmcsMidTwinMeanPTwoRidge}{0.26}  % StepM-adjusted p-value (two-sided), Ridge
\providecommand{\rhmcsMidTwinMeanRawTwoRidge}{}\renewcommand{\rhmcsMidTwinMeanRawTwoRidge}{0.16}  % single-test bootstrap p-value (two-sided), Ridge
\providecommand{\rhmcsMidTwinMeanPLasso}{}\renewcommand{\rhmcsMidTwinMeanPLasso}{0.11}  % StepM-adjusted p-value (one-sided), Lasso
\providecommand{\rhmcsMidTwinMeanRawLasso}{}\renewcommand{\rhmcsMidTwinMeanRawLasso}{0.063}  % single-test bootstrap p-value (one-sided), Lasso
\providecommand{\rhmcsMidTwinMeanPTwoLasso}{}\renewcommand{\rhmcsMidTwinMeanPTwoLasso}{0.21}  % StepM-adjusted p-value (two-sided), Lasso
\providecommand{\rhmcsMidTwinMeanRawTwoLasso}{}\renewcommand{\rhmcsMidTwinMeanRawTwoLasso}{0.12}  % single-test bootstrap p-value (two-sided), Lasso
\providecommand{\rhmcsMidTwinMeanPEnet}{}\renewcommand{\rhmcsMidTwinMeanPEnet}{0.069}  % StepM-adjusted p-value (one-sided), Elastic net
\providecommand{\rhmcsMidTwinMeanRawEnet}{}\renewcommand{\rhmcsMidTwinMeanRawEnet}{0.036}  % single-test bootstrap p-value (one-sided), Elastic net
\providecommand{\rhmcsMidTwinMeanPTwoEnet}{}\renewcommand{\rhmcsMidTwinMeanPTwoEnet}{0.13}  % StepM-adjusted p-value (two-sided), Elastic net
\providecommand{\rhmcsMidTwinMeanRawTwoEnet}{}\renewcommand{\rhmcsMidTwinMeanRawTwoEnet}{0.063}  % single-test bootstrap p-value (two-sided), Elastic net
\providecommand{\rhmcsMidTwinMeanPLgbm}{}\renewcommand{\rhmcsMidTwinMeanPLgbm}{0.060}  % StepM-adjusted p-value (one-sided), LightGBM
\providecommand{\rhmcsMidTwinMeanRawLgbm}{}\renewcommand{\rhmcsMidTwinMeanRawLgbm}{0.025}  % single-test bootstrap p-value (one-sided), LightGBM
\providecommand{\rhmcsMidTwinMeanPTwoLgbm}{}\renewcommand{\rhmcsMidTwinMeanPTwoLgbm}{0.11}  % StepM-adjusted p-value (two-sided), LightGBM
\providecommand{\rhmcsMidTwinMeanRawTwoLgbm}{}\renewcommand{\rhmcsMidTwinMeanRawTwoLgbm}{0.047}  % single-test bootstrap p-value (two-sided), LightGBM
\providecommand{\rhmcsMidTwinMeanPXgb}{}\renewcommand{\rhmcsMidTwinMeanPXgb}{0.13}  % StepM-adjusted p-value (one-sided), XGBoost
\providecommand{\rhmcsMidTwinMeanRawXgb}{}\renewcommand{\rhmcsMidTwinMeanRawXgb}{0.101}  % single-test bootstrap p-value (one-sided), XGBoost
\providecommand{\rhmcsMidTwinMeanPTwoXgb}{}\renewcommand{\rhmcsMidTwinMeanPTwoXgb}{0.26}  % StepM-adjusted p-value (two-sided), XGBoost
\providecommand{\rhmcsMidTwinMeanRawTwoXgb}{}\renewcommand{\rhmcsMidTwinMeanRawTwoXgb}{0.20}  % single-test bootstrap p-value (two-sided), XGBoost
\providecommand{\rhmcsMidTwinMeanPRf}{}\renewcommand{\rhmcsMidTwinMeanPRf}{0.027}  % StepM-adjusted p-value (one-sided), Random forest
\providecommand{\rhmcsMidTwinMeanRawRf}{}\renewcommand{\rhmcsMidTwinMeanRawRf}{0.013}  % single-test bootstrap p-value (one-sided), Random forest
\providecommand{\rhmcsMidTwinMeanPTwoRf}{}\renewcommand{\rhmcsMidTwinMeanPTwoRf}{0.042}  % StepM-adjusted p-value (two-sided), Random forest
\providecommand{\rhmcsMidTwinMeanRawTwoRf}{}\renewcommand{\rhmcsMidTwinMeanRawTwoRf}{0.013}  % single-test bootstrap p-value (two-sided), Random forest
\providecommand{\rhmcsMidTwinMeanNRejLo}{}\renewcommand{\rhmcsMidTwinMeanNRejLo}{1}  % StepM rejections at 5 %, block length 5
\providecommand{\rhmcsMidTwinMeanSpaCLo}{}\renewcommand{\rhmcsMidTwinMeanSpaCLo}{0.036}  % SPA p-value (consistent), block length 5
\providecommand{\rhmcsMidTwinMeanNRejTwoLo}{}\renewcommand{\rhmcsMidTwinMeanNRejTwoLo}{0}  % ... two-sided, block length 5
\providecommand{\rhmcsMidTwinMeanPMinLo}{}\renewcommand{\rhmcsMidTwinMeanPMinLo}{0.036}  % smallest adjusted p, block length 5
\providecommand{\rhmcsMidTwinMeanPTwoMinLo}{}\renewcommand{\rhmcsMidTwinMeanPTwoMinLo}{0.053}  % ... two-sided, block length 5
\providecommand{\rhmcsMidTwinMeanNRejHi}{}\renewcommand{\rhmcsMidTwinMeanNRejHi}{1}  % StepM rejections at 5 %, block length 20
\providecommand{\rhmcsMidTwinMeanSpaCHi}{}\renewcommand{\rhmcsMidTwinMeanSpaCHi}{0.033}  % SPA p-value (consistent), block length 20
\providecommand{\rhmcsMidTwinMeanNRejTwoHi}{}\renewcommand{\rhmcsMidTwinMeanNRejTwoHi}{1}  % ... two-sided, block length 20
\providecommand{\rhmcsMidTwinMeanPMinHi}{}\renewcommand{\rhmcsMidTwinMeanPMinHi}{0.033}  % smallest adjusted p, block length 20
\providecommand{\rhmcsMidTwinMeanPTwoMinHi}{}\renewcommand{\rhmcsMidTwinMeanPTwoMinHi}{0.047}  % ... two-sided, block length 20
\providecommand{\rhmcsMidTwinSrSpaC}{}\renewcommand{\rhmcsMidTwinSrSpaC}{0.007}  % SPA p-value (consistent), family MidTwinSr
\providecommand{\rhmcsMidTwinSrSpaL}{}\renewcommand{\rhmcsMidTwinSrSpaL}{0.007}  % ... lower
\providecommand{\rhmcsMidTwinSrSpaU}{}\renewcommand{\rhmcsMidTwinSrSpaU}{0.007}  % ... upper
\providecommand{\rhmcsMidTwinSrNRej}{}\renewcommand{\rhmcsMidTwinSrNRej}{1}  % StepM rejections at FWER 5 % (of 6)
\providecommand{\rhmcsMidTwinSrNRejWord}{}\renewcommand{\rhmcsMidTwinSrNRejWord}{one}  % ... as a word
\providecommand{\rhmcsMidTwinSrNRaw}{}\renewcommand{\rhmcsMidTwinSrNRaw}{3}  % single-test bootstrap p < 0.05 (no adjustment)
\providecommand{\rhmcsMidTwinSrNRawWord}{}\renewcommand{\rhmcsMidTwinSrNRawWord}{three}  % ... as a word
\providecommand{\rhmcsMidTwinSrNPos}{}\renewcommand{\rhmcsMidTwinSrNPos}{6}  % alternatives with a positive estimate (better than the benchmark)
\providecommand{\rhmcsMidTwinSrPMin}{}\renewcommand{\rhmcsMidTwinSrPMin}{0.007}  % smallest StepM-adjusted p-value
\providecommand{\rhmcsMidTwinSrPMax}{}\renewcommand{\rhmcsMidTwinSrPMax}{0.14}  % largest StepM-adjusted p-value
\providecommand{\rhmcsMidTwinSrPTwoMin}{}\renewcommand{\rhmcsMidTwinSrPTwoMin}{0.012}  % smallest two-sided StepM-adjusted p-value
\providecommand{\rhmcsMidTwinSrPTwoMax}{}\renewcommand{\rhmcsMidTwinSrPTwoMax}{0.27}  % largest two-sided StepM-adjusted p-value
\providecommand{\rhmcsMidTwinSrRawMin}{}\renewcommand{\rhmcsMidTwinSrRawMin}{0.002}  % smallest single-test bootstrap p-value
\providecommand{\rhmcsMidTwinSrRej}{}\renewcommand{\rhmcsMidTwinSrRej}{the random forest on the full predictor set}  % the alternatives StepM rejects at 5 % (none: none)
\providecommand{\rhmcsMidTwinSrNRejTwo}{}\renewcommand{\rhmcsMidTwinSrNRejTwo}{1}  % StepM rejections at FWER 5 %, two-sided
\providecommand{\rhmcsMidTwinSrNRejTwoWord}{}\renewcommand{\rhmcsMidTwinSrNRejTwoWord}{one}  % ... as a word
\providecommand{\rhmcsMidTwinSrNRawTwo}{}\renewcommand{\rhmcsMidTwinSrNRawTwo}{2}  % single-test bootstrap p < 0.05, two-sided
\providecommand{\rhmcsMidTwinSrRejTwo}{}\renewcommand{\rhmcsMidTwinSrRejTwo}{the random forest on the full predictor set}  % the alternatives the two-sided StepM rejects at 5 % (none: none)
\providecommand{\rhmcsMidTwinSrPRidge}{}\renewcommand{\rhmcsMidTwinSrPRidge}{0.14}  % StepM-adjusted p-value (one-sided), Ridge
\providecommand{\rhmcsMidTwinSrRawRidge}{}\renewcommand{\rhmcsMidTwinSrRawRidge}{0.089}  % single-test bootstrap p-value (one-sided), Ridge
\providecommand{\rhmcsMidTwinSrPTwoRidge}{}\renewcommand{\rhmcsMidTwinSrPTwoRidge}{0.27}  % StepM-adjusted p-value (two-sided), Ridge
\providecommand{\rhmcsMidTwinSrRawTwoRidge}{}\renewcommand{\rhmcsMidTwinSrRawTwoRidge}{0.17}  % single-test bootstrap p-value (two-sided), Ridge
\providecommand{\rhmcsMidTwinSrPLasso}{}\renewcommand{\rhmcsMidTwinSrPLasso}{0.12}  % StepM-adjusted p-value (one-sided), Lasso
\providecommand{\rhmcsMidTwinSrRawLasso}{}\renewcommand{\rhmcsMidTwinSrRawLasso}{0.071}  % single-test bootstrap p-value (one-sided), Lasso
\providecommand{\rhmcsMidTwinSrPTwoLasso}{}\renewcommand{\rhmcsMidTwinSrPTwoLasso}{0.22}  % StepM-adjusted p-value (two-sided), Lasso
\providecommand{\rhmcsMidTwinSrRawTwoLasso}{}\renewcommand{\rhmcsMidTwinSrRawTwoLasso}{0.12}  % single-test bootstrap p-value (two-sided), Lasso
\providecommand{\rhmcsMidTwinSrPEnet}{}\renewcommand{\rhmcsMidTwinSrPEnet}{0.082}  % StepM-adjusted p-value (one-sided), Elastic net
\providecommand{\rhmcsMidTwinSrRawEnet}{}\renewcommand{\rhmcsMidTwinSrRawEnet}{0.045}  % single-test bootstrap p-value (one-sided), Elastic net
\providecommand{\rhmcsMidTwinSrPTwoEnet}{}\renewcommand{\rhmcsMidTwinSrPTwoEnet}{0.15}  % StepM-adjusted p-value (two-sided), Elastic net
\providecommand{\rhmcsMidTwinSrRawTwoEnet}{}\renewcommand{\rhmcsMidTwinSrRawTwoEnet}{0.072}  % single-test bootstrap p-value (two-sided), Elastic net
\providecommand{\rhmcsMidTwinSrPLgbm}{}\renewcommand{\rhmcsMidTwinSrPLgbm}{0.065}  % StepM-adjusted p-value (one-sided), LightGBM
\providecommand{\rhmcsMidTwinSrRawLgbm}{}\renewcommand{\rhmcsMidTwinSrRawLgbm}{0.027}  % single-test bootstrap p-value (one-sided), LightGBM
\providecommand{\rhmcsMidTwinSrPTwoLgbm}{}\renewcommand{\rhmcsMidTwinSrPTwoLgbm}{0.12}  % StepM-adjusted p-value (two-sided), LightGBM
\providecommand{\rhmcsMidTwinSrRawTwoLgbm}{}\renewcommand{\rhmcsMidTwinSrRawTwoLgbm}{0.050}  % single-test bootstrap p-value (two-sided), LightGBM
\providecommand{\rhmcsMidTwinSrPXgb}{}\renewcommand{\rhmcsMidTwinSrPXgb}{0.14}  % StepM-adjusted p-value (one-sided), XGBoost
\providecommand{\rhmcsMidTwinSrRawXgb}{}\renewcommand{\rhmcsMidTwinSrRawXgb}{0.11}  % single-test bootstrap p-value (one-sided), XGBoost
\providecommand{\rhmcsMidTwinSrPTwoXgb}{}\renewcommand{\rhmcsMidTwinSrPTwoXgb}{0.27}  % StepM-adjusted p-value (two-sided), XGBoost
\providecommand{\rhmcsMidTwinSrRawTwoXgb}{}\renewcommand{\rhmcsMidTwinSrRawTwoXgb}{0.20}  % single-test bootstrap p-value (two-sided), XGBoost
\providecommand{\rhmcsMidTwinSrPRf}{}\renewcommand{\rhmcsMidTwinSrPRf}{0.007}  % StepM-adjusted p-value (one-sided), Random forest
\providecommand{\rhmcsMidTwinSrRawRf}{}\renewcommand{\rhmcsMidTwinSrRawRf}{0.002}  % single-test bootstrap p-value (one-sided), Random forest
\providecommand{\rhmcsMidTwinSrPTwoRf}{}\renewcommand{\rhmcsMidTwinSrPTwoRf}{0.012}  % StepM-adjusted p-value (two-sided), Random forest
\providecommand{\rhmcsMidTwinSrRawTwoRf}{}\renewcommand{\rhmcsMidTwinSrRawTwoRf}{0.003}  % single-test bootstrap p-value (two-sided), Random forest
\providecommand{\rhmcsMidTwinSrNRejLo}{}\renewcommand{\rhmcsMidTwinSrNRejLo}{1}  % StepM rejections at 5 %, block length 5
\providecommand{\rhmcsMidTwinSrSpaCLo}{}\renewcommand{\rhmcsMidTwinSrSpaCLo}{0.011}  % SPA p-value (consistent), block length 5
\providecommand{\rhmcsMidTwinSrNRejTwoLo}{}\renewcommand{\rhmcsMidTwinSrNRejTwoLo}{1}  % ... two-sided, block length 5
\providecommand{\rhmcsMidTwinSrPMinLo}{}\renewcommand{\rhmcsMidTwinSrPMinLo}{0.011}  % smallest adjusted p, block length 5
\providecommand{\rhmcsMidTwinSrPTwoMinLo}{}\renewcommand{\rhmcsMidTwinSrPTwoMinLo}{0.017}  % ... two-sided, block length 5
\providecommand{\rhmcsMidTwinSrNRejHi}{}\renewcommand{\rhmcsMidTwinSrNRejHi}{1}  % StepM rejections at 5 %, block length 20
\providecommand{\rhmcsMidTwinSrSpaCHi}{}\renewcommand{\rhmcsMidTwinSrSpaCHi}{0.010}  % SPA p-value (consistent), block length 20
\providecommand{\rhmcsMidTwinSrNRejTwoHi}{}\renewcommand{\rhmcsMidTwinSrNRejTwoHi}{1}  % ... two-sided, block length 20
\providecommand{\rhmcsMidTwinSrPMinHi}{}\renewcommand{\rhmcsMidTwinSrPMinHi}{0.010}  % smallest adjusted p, block length 20
\providecommand{\rhmcsMidTwinSrPTwoMinHi}{}\renewcommand{\rhmcsMidTwinSrPTwoMinHi}{0.016}  % ... two-sided, block length 20
\providecommand{\rhmcsBaShortMeanSpaC}{}\renewcommand{\rhmcsBaShortMeanSpaC}{0.025}  % SPA p-value (consistent), family BaShortMean
\providecommand{\rhmcsBaShortMeanSpaL}{}\renewcommand{\rhmcsBaShortMeanSpaL}{0.025}  % ... lower
\providecommand{\rhmcsBaShortMeanSpaU}{}\renewcommand{\rhmcsBaShortMeanSpaU}{0.025}  % ... upper
\providecommand{\rhmcsBaShortMeanNRej}{}\renewcommand{\rhmcsBaShortMeanNRej}{3}  % StepM rejections at FWER 5 % (of 6)
\providecommand{\rhmcsBaShortMeanNRejWord}{}\renewcommand{\rhmcsBaShortMeanNRejWord}{three}  % ... as a word
\providecommand{\rhmcsBaShortMeanNRaw}{}\renewcommand{\rhmcsBaShortMeanNRaw}{6}  % single-test bootstrap p < 0.05 (no adjustment)
\providecommand{\rhmcsBaShortMeanNRawWord}{}\renewcommand{\rhmcsBaShortMeanNRawWord}{six}  % ... as a word
\providecommand{\rhmcsBaShortMeanNPos}{}\renewcommand{\rhmcsBaShortMeanNPos}{6}  % alternatives with a positive estimate (better than the benchmark)
\providecommand{\rhmcsBaShortMeanPMin}{}\renewcommand{\rhmcsBaShortMeanPMin}{0.025}  % smallest StepM-adjusted p-value
\providecommand{\rhmcsBaShortMeanPMax}{}\renewcommand{\rhmcsBaShortMeanPMax}{0.050}  % largest StepM-adjusted p-value
\providecommand{\rhmcsBaShortMeanPTwoMin}{}\renewcommand{\rhmcsBaShortMeanPTwoMin}{0.034}  % smallest two-sided StepM-adjusted p-value
\providecommand{\rhmcsBaShortMeanPTwoMax}{}\renewcommand{\rhmcsBaShortMeanPTwoMax}{0.082}  % largest two-sided StepM-adjusted p-value
\providecommand{\rhmcsBaShortMeanRawMin}{}\renewcommand{\rhmcsBaShortMeanRawMin}{0.015}  % smallest single-test bootstrap p-value
\providecommand{\rhmcsBaShortMeanRej}{}\renewcommand{\rhmcsBaShortMeanRej}{the lasso on the full predictor set, the elastic net on the full predictor set, and the random forest on the full predictor set}  % the alternatives StepM rejects at 5 % (none: none)
\providecommand{\rhmcsBaShortMeanNRejTwo}{}\renewcommand{\rhmcsBaShortMeanNRejTwo}{2}  % StepM rejections at FWER 5 %, two-sided
\providecommand{\rhmcsBaShortMeanNRejTwoWord}{}\renewcommand{\rhmcsBaShortMeanNRejTwoWord}{two}  % ... as a word
\providecommand{\rhmcsBaShortMeanNRawTwo}{}\renewcommand{\rhmcsBaShortMeanNRawTwo}{3}  % single-test bootstrap p < 0.05, two-sided
\providecommand{\rhmcsBaShortMeanRejTwo}{}\renewcommand{\rhmcsBaShortMeanRejTwo}{the elastic net on the full predictor set and the random forest on the full predictor set}  % the alternatives the two-sided StepM rejects at 5 % (none: none)
\providecommand{\rhmcsBaShortMeanPRidge}{}\renewcommand{\rhmcsBaShortMeanPRidge}{0.050}  % StepM-adjusted p-value (one-sided), Ridge
\providecommand{\rhmcsBaShortMeanRawRidge}{}\renewcommand{\rhmcsBaShortMeanRawRidge}{0.036}  % single-test bootstrap p-value (one-sided), Ridge
\providecommand{\rhmcsBaShortMeanPTwoRidge}{}\renewcommand{\rhmcsBaShortMeanPTwoRidge}{0.082}  % StepM-adjusted p-value (two-sided), Ridge
\providecommand{\rhmcsBaShortMeanRawTwoRidge}{}\renewcommand{\rhmcsBaShortMeanRawTwoRidge}{0.056}  % single-test bootstrap p-value (two-sided), Ridge
\providecommand{\rhmcsBaShortMeanPLasso}{}\renewcommand{\rhmcsBaShortMeanPLasso}{0.037}  % StepM-adjusted p-value (one-sided), Lasso
\providecommand{\rhmcsBaShortMeanRawLasso}{}\renewcommand{\rhmcsBaShortMeanRawLasso}{0.025}  % single-test bootstrap p-value (one-sided), Lasso
\providecommand{\rhmcsBaShortMeanPTwoLasso}{}\renewcommand{\rhmcsBaShortMeanPTwoLasso}{0.058}  % StepM-adjusted p-value (two-sided), Lasso
\providecommand{\rhmcsBaShortMeanRawTwoLasso}{}\renewcommand{\rhmcsBaShortMeanRawTwoLasso}{0.035}  % single-test bootstrap p-value (two-sided), Lasso
\providecommand{\rhmcsBaShortMeanPEnet}{}\renewcommand{\rhmcsBaShortMeanPEnet}{0.031}  % StepM-adjusted p-value (one-sided), Elastic net
\providecommand{\rhmcsBaShortMeanRawEnet}{}\renewcommand{\rhmcsBaShortMeanRawEnet}{0.020}  % single-test bootstrap p-value (one-sided), Elastic net
\providecommand{\rhmcsBaShortMeanPTwoEnet}{}\renewcommand{\rhmcsBaShortMeanPTwoEnet}{0.045}  % StepM-adjusted p-value (two-sided), Elastic net
\providecommand{\rhmcsBaShortMeanRawTwoEnet}{}\renewcommand{\rhmcsBaShortMeanRawTwoEnet}{0.025}  % single-test bootstrap p-value (two-sided), Elastic net
\providecommand{\rhmcsBaShortMeanPLgbm}{}\renewcommand{\rhmcsBaShortMeanPLgbm}{0.050}  % StepM-adjusted p-value (one-sided), LightGBM
\providecommand{\rhmcsBaShortMeanRawLgbm}{}\renewcommand{\rhmcsBaShortMeanRawLgbm}{0.041}  % single-test bootstrap p-value (one-sided), LightGBM
\providecommand{\rhmcsBaShortMeanPTwoLgbm}{}\renewcommand{\rhmcsBaShortMeanPTwoLgbm}{0.082}  % StepM-adjusted p-value (two-sided), LightGBM
\providecommand{\rhmcsBaShortMeanRawTwoLgbm}{}\renewcommand{\rhmcsBaShortMeanRawTwoLgbm}{0.067}  % single-test bootstrap p-value (two-sided), LightGBM
\providecommand{\rhmcsBaShortMeanPXgb}{}\renewcommand{\rhmcsBaShortMeanPXgb}{0.050}  % StepM-adjusted p-value (one-sided), XGBoost
\providecommand{\rhmcsBaShortMeanRawXgb}{}\renewcommand{\rhmcsBaShortMeanRawXgb}{0.038}  % single-test bootstrap p-value (one-sided), XGBoost
\providecommand{\rhmcsBaShortMeanPTwoXgb}{}\renewcommand{\rhmcsBaShortMeanPTwoXgb}{0.082}  % StepM-adjusted p-value (two-sided), XGBoost
\providecommand{\rhmcsBaShortMeanRawTwoXgb}{}\renewcommand{\rhmcsBaShortMeanRawTwoXgb}{0.063}  % single-test bootstrap p-value (two-sided), XGBoost
\providecommand{\rhmcsBaShortMeanPRf}{}\renewcommand{\rhmcsBaShortMeanPRf}{0.025}  % StepM-adjusted p-value (one-sided), Random forest
\providecommand{\rhmcsBaShortMeanRawRf}{}\renewcommand{\rhmcsBaShortMeanRawRf}{0.015}  % single-test bootstrap p-value (one-sided), Random forest
\providecommand{\rhmcsBaShortMeanPTwoRf}{}\renewcommand{\rhmcsBaShortMeanPTwoRf}{0.034}  % StepM-adjusted p-value (two-sided), Random forest
\providecommand{\rhmcsBaShortMeanRawTwoRf}{}\renewcommand{\rhmcsBaShortMeanRawTwoRf}{0.019}  % single-test bootstrap p-value (two-sided), Random forest
\providecommand{\rhmcsBaShortMeanNRejLo}{}\renewcommand{\rhmcsBaShortMeanNRejLo}{3}  % StepM rejections at 5 %, block length 5
\providecommand{\rhmcsBaShortMeanSpaCLo}{}\renewcommand{\rhmcsBaShortMeanSpaCLo}{0.037}  % SPA p-value (consistent), block length 5
\providecommand{\rhmcsBaShortMeanNRejTwoLo}{}\renewcommand{\rhmcsBaShortMeanNRejTwoLo}{0}  % ... two-sided, block length 5
\providecommand{\rhmcsBaShortMeanPMinLo}{}\renewcommand{\rhmcsBaShortMeanPMinLo}{0.037}  % smallest adjusted p, block length 5
\providecommand{\rhmcsBaShortMeanPTwoMinLo}{}\renewcommand{\rhmcsBaShortMeanPTwoMinLo}{0.054}  % ... two-sided, block length 5
\providecommand{\rhmcsBaShortMeanNRejHi}{}\renewcommand{\rhmcsBaShortMeanNRejHi}{3}  % StepM rejections at 5 %, block length 20
\providecommand{\rhmcsBaShortMeanSpaCHi}{}\renewcommand{\rhmcsBaShortMeanSpaCHi}{0.027}  % SPA p-value (consistent), block length 20
\providecommand{\rhmcsBaShortMeanNRejTwoHi}{}\renewcommand{\rhmcsBaShortMeanNRejTwoHi}{1}  % ... two-sided, block length 20
\providecommand{\rhmcsBaShortMeanPMinHi}{}\renewcommand{\rhmcsBaShortMeanPMinHi}{0.027}  % smallest adjusted p, block length 20
\providecommand{\rhmcsBaShortMeanPTwoMinHi}{}\renewcommand{\rhmcsBaShortMeanPTwoMinHi}{0.035}  % ... two-sided, block length 20
\providecommand{\rhmcsBaShortSrSpaC}{}\renewcommand{\rhmcsBaShortSrSpaC}{0.007}  % SPA p-value (consistent), family BaShortSr
\providecommand{\rhmcsBaShortSrSpaL}{}\renewcommand{\rhmcsBaShortSrSpaL}{0.007}  % ... lower
\providecommand{\rhmcsBaShortSrSpaU}{}\renewcommand{\rhmcsBaShortSrSpaU}{0.007}  % ... upper
\providecommand{\rhmcsBaShortSrNRej}{}\renewcommand{\rhmcsBaShortSrNRej}{6}  % StepM rejections at FWER 5 % (of 6)
\providecommand{\rhmcsBaShortSrNRejWord}{}\renewcommand{\rhmcsBaShortSrNRejWord}{six}  % ... as a word
\providecommand{\rhmcsBaShortSrNRaw}{}\renewcommand{\rhmcsBaShortSrNRaw}{6}  % single-test bootstrap p < 0.05 (no adjustment)
\providecommand{\rhmcsBaShortSrNRawWord}{}\renewcommand{\rhmcsBaShortSrNRawWord}{six}  % ... as a word
\providecommand{\rhmcsBaShortSrNPos}{}\renewcommand{\rhmcsBaShortSrNPos}{6}  % alternatives with a positive estimate (better than the benchmark)
\providecommand{\rhmcsBaShortSrPMin}{}\renewcommand{\rhmcsBaShortSrPMin}{0.007}  % smallest StepM-adjusted p-value
\providecommand{\rhmcsBaShortSrPMax}{}\renewcommand{\rhmcsBaShortSrPMax}{0.027}  % largest StepM-adjusted p-value
\providecommand{\rhmcsBaShortSrPTwoMin}{}\renewcommand{\rhmcsBaShortSrPTwoMin}{0.015}  % smallest two-sided StepM-adjusted p-value
\providecommand{\rhmcsBaShortSrPTwoMax}{}\renewcommand{\rhmcsBaShortSrPTwoMax}{0.058}  % largest two-sided StepM-adjusted p-value
\providecommand{\rhmcsBaShortSrRawMin}{}\renewcommand{\rhmcsBaShortSrRawMin}{0.002}  % smallest single-test bootstrap p-value
\providecommand{\rhmcsBaShortSrRej}{}\renewcommand{\rhmcsBaShortSrRej}{ridge regression on the full predictor set, the lasso on the full predictor set, the elastic net on the full predictor set, LightGBM on the full predictor set, XGBoost on the full predictor set, and the random forest on the full predictor set}  % the alternatives StepM rejects at 5 % (none: none)
\providecommand{\rhmcsBaShortSrNRejTwo}{}\renewcommand{\rhmcsBaShortSrNRejTwo}{3}  % StepM rejections at FWER 5 %, two-sided
\providecommand{\rhmcsBaShortSrNRejTwoWord}{}\renewcommand{\rhmcsBaShortSrNRejTwoWord}{three}  % ... as a word
\providecommand{\rhmcsBaShortSrNRawTwo}{}\renewcommand{\rhmcsBaShortSrNRawTwo}{6}  % single-test bootstrap p < 0.05, two-sided
\providecommand{\rhmcsBaShortSrRejTwo}{}\renewcommand{\rhmcsBaShortSrRejTwo}{the lasso on the full predictor set, the elastic net on the full predictor set, and the random forest on the full predictor set}  % the alternatives the two-sided StepM rejects at 5 % (none: none)
\providecommand{\rhmcsBaShortSrPRidge}{}\renewcommand{\rhmcsBaShortSrPRidge}{0.027}  % StepM-adjusted p-value (one-sided), Ridge
\providecommand{\rhmcsBaShortSrRawRidge}{}\renewcommand{\rhmcsBaShortSrRawRidge}{0.018}  % single-test bootstrap p-value (one-sided), Ridge
\providecommand{\rhmcsBaShortSrPTwoRidge}{}\renewcommand{\rhmcsBaShortSrPTwoRidge}{0.058}  % StepM-adjusted p-value (two-sided), Ridge
\providecommand{\rhmcsBaShortSrRawTwoRidge}{}\renewcommand{\rhmcsBaShortSrRawTwoRidge}{0.037}  % single-test bootstrap p-value (two-sided), Ridge
\providecommand{\rhmcsBaShortSrPLasso}{}\renewcommand{\rhmcsBaShortSrPLasso}{0.017}  % StepM-adjusted p-value (one-sided), Lasso
\providecommand{\rhmcsBaShortSrRawLasso}{}\renewcommand{\rhmcsBaShortSrRawLasso}{0.010}  % single-test bootstrap p-value (one-sided), Lasso
\providecommand{\rhmcsBaShortSrPTwoLasso}{}\renewcommand{\rhmcsBaShortSrPTwoLasso}{0.039}  % StepM-adjusted p-value (two-sided), Lasso
\providecommand{\rhmcsBaShortSrRawTwoLasso}{}\renewcommand{\rhmcsBaShortSrRawTwoLasso}{0.021}  % single-test bootstrap p-value (two-sided), Lasso
\providecommand{\rhmcsBaShortSrPEnet}{}\renewcommand{\rhmcsBaShortSrPEnet}{0.011}  % StepM-adjusted p-value (one-sided), Elastic net
\providecommand{\rhmcsBaShortSrRawEnet}{}\renewcommand{\rhmcsBaShortSrRawEnet}{0.007}  % single-test bootstrap p-value (one-sided), Elastic net
\providecommand{\rhmcsBaShortSrPTwoEnet}{}\renewcommand{\rhmcsBaShortSrPTwoEnet}{0.026}  % StepM-adjusted p-value (two-sided), Elastic net
\providecommand{\rhmcsBaShortSrRawTwoEnet}{}\renewcommand{\rhmcsBaShortSrRawTwoEnet}{0.012}  % single-test bootstrap p-value (two-sided), Elastic net
\providecommand{\rhmcsBaShortSrPLgbm}{}\renewcommand{\rhmcsBaShortSrPLgbm}{0.027}  % StepM-adjusted p-value (one-sided), LightGBM
\providecommand{\rhmcsBaShortSrRawLgbm}{}\renewcommand{\rhmcsBaShortSrRawLgbm}{0.020}  % single-test bootstrap p-value (one-sided), LightGBM
\providecommand{\rhmcsBaShortSrPTwoLgbm}{}\renewcommand{\rhmcsBaShortSrPTwoLgbm}{0.058}  % StepM-adjusted p-value (two-sided), LightGBM
\providecommand{\rhmcsBaShortSrRawTwoLgbm}{}\renewcommand{\rhmcsBaShortSrRawTwoLgbm}{0.045}  % single-test bootstrap p-value (two-sided), LightGBM
\providecommand{\rhmcsBaShortSrPXgb}{}\renewcommand{\rhmcsBaShortSrPXgb}{0.027}  % StepM-adjusted p-value (one-sided), XGBoost
\providecommand{\rhmcsBaShortSrRawXgb}{}\renewcommand{\rhmcsBaShortSrRawXgb}{0.019}  % single-test bootstrap p-value (one-sided), XGBoost
\providecommand{\rhmcsBaShortSrPTwoXgb}{}\renewcommand{\rhmcsBaShortSrPTwoXgb}{0.058}  % StepM-adjusted p-value (two-sided), XGBoost
\providecommand{\rhmcsBaShortSrRawTwoXgb}{}\renewcommand{\rhmcsBaShortSrRawTwoXgb}{0.043}  % single-test bootstrap p-value (two-sided), XGBoost
\providecommand{\rhmcsBaShortSrPRf}{}\renewcommand{\rhmcsBaShortSrPRf}{0.007}  % StepM-adjusted p-value (one-sided), Random forest
\providecommand{\rhmcsBaShortSrRawRf}{}\renewcommand{\rhmcsBaShortSrRawRf}{0.002}  % single-test bootstrap p-value (one-sided), Random forest
\providecommand{\rhmcsBaShortSrPTwoRf}{}\renewcommand{\rhmcsBaShortSrPTwoRf}{0.015}  % StepM-adjusted p-value (two-sided), Random forest
\providecommand{\rhmcsBaShortSrRawTwoRf}{}\renewcommand{\rhmcsBaShortSrRawTwoRf}{0.006}  % single-test bootstrap p-value (two-sided), Random forest
\providecommand{\rhmcsBaShortSrNRejLo}{}\renewcommand{\rhmcsBaShortSrNRejLo}{6}  % StepM rejections at 5 %, block length 5
\providecommand{\rhmcsBaShortSrSpaCLo}{}\renewcommand{\rhmcsBaShortSrSpaCLo}{0.010}  % SPA p-value (consistent), block length 5
\providecommand{\rhmcsBaShortSrNRejTwoLo}{}\renewcommand{\rhmcsBaShortSrNRejTwoLo}{3}  % ... two-sided, block length 5
\providecommand{\rhmcsBaShortSrPMinLo}{}\renewcommand{\rhmcsBaShortSrPMinLo}{0.010}  % smallest adjusted p, block length 5
\providecommand{\rhmcsBaShortSrPTwoMinLo}{}\renewcommand{\rhmcsBaShortSrPTwoMinLo}{0.027}  % ... two-sided, block length 5
\providecommand{\rhmcsBaShortSrNRejHi}{}\renewcommand{\rhmcsBaShortSrNRejHi}{6}  % StepM rejections at 5 %, block length 20
\providecommand{\rhmcsBaShortSrSpaCHi}{}\renewcommand{\rhmcsBaShortSrSpaCHi}{0.008}  % SPA p-value (consistent), block length 20
\providecommand{\rhmcsBaShortSrNRejTwoHi}{}\renewcommand{\rhmcsBaShortSrNRejTwoHi}{3}  % ... two-sided, block length 20
\providecommand{\rhmcsBaShortSrPMinHi}{}\renewcommand{\rhmcsBaShortSrPMinHi}{0.008}  % smallest adjusted p, block length 20
\providecommand{\rhmcsBaShortSrPTwoMinHi}{}\renewcommand{\rhmcsBaShortSrPTwoMinHi}{0.014}  % ... two-sided, block length 20
\providecommand{\rhmcsBaTwinMeanSpaC}{}\renewcommand{\rhmcsBaTwinMeanSpaC}{0.028}  % SPA p-value (consistent), family BaTwinMean
\providecommand{\rhmcsBaTwinMeanSpaL}{}\renewcommand{\rhmcsBaTwinMeanSpaL}{0.028}  % ... lower
\providecommand{\rhmcsBaTwinMeanSpaU}{}\renewcommand{\rhmcsBaTwinMeanSpaU}{0.028}  % ... upper
\providecommand{\rhmcsBaTwinMeanNRej}{}\renewcommand{\rhmcsBaTwinMeanNRej}{1}  % StepM rejections at FWER 5 % (of 6)
\providecommand{\rhmcsBaTwinMeanNRejWord}{}\renewcommand{\rhmcsBaTwinMeanNRejWord}{one}  % ... as a word
\providecommand{\rhmcsBaTwinMeanNRaw}{}\renewcommand{\rhmcsBaTwinMeanNRaw}{3}  % single-test bootstrap p < 0.05 (no adjustment)
\providecommand{\rhmcsBaTwinMeanNRawWord}{}\renewcommand{\rhmcsBaTwinMeanNRawWord}{three}  % ... as a word
\providecommand{\rhmcsBaTwinMeanNPos}{}\renewcommand{\rhmcsBaTwinMeanNPos}{6}  % alternatives with a positive estimate (better than the benchmark)
\providecommand{\rhmcsBaTwinMeanPMin}{}\renewcommand{\rhmcsBaTwinMeanPMin}{0.028}  % smallest StepM-adjusted p-value
\providecommand{\rhmcsBaTwinMeanPMax}{}\renewcommand{\rhmcsBaTwinMeanPMax}{0.13}  % largest StepM-adjusted p-value
\providecommand{\rhmcsBaTwinMeanPTwoMin}{}\renewcommand{\rhmcsBaTwinMeanPTwoMin}{0.042}  % smallest two-sided StepM-adjusted p-value
\providecommand{\rhmcsBaTwinMeanPTwoMax}{}\renewcommand{\rhmcsBaTwinMeanPTwoMax}{0.26}  % largest two-sided StepM-adjusted p-value
\providecommand{\rhmcsBaTwinMeanRawMin}{}\renewcommand{\rhmcsBaTwinMeanRawMin}{0.013}  % smallest single-test bootstrap p-value
\providecommand{\rhmcsBaTwinMeanRej}{}\renewcommand{\rhmcsBaTwinMeanRej}{the random forest on the full predictor set}  % the alternatives StepM rejects at 5 % (none: none)
\providecommand{\rhmcsBaTwinMeanNRejTwo}{}\renewcommand{\rhmcsBaTwinMeanNRejTwo}{1}  % StepM rejections at FWER 5 %, two-sided
\providecommand{\rhmcsBaTwinMeanNRejTwoWord}{}\renewcommand{\rhmcsBaTwinMeanNRejTwoWord}{one}  % ... as a word
\providecommand{\rhmcsBaTwinMeanNRawTwo}{}\renewcommand{\rhmcsBaTwinMeanNRawTwo}{2}  % single-test bootstrap p < 0.05, two-sided
\providecommand{\rhmcsBaTwinMeanRejTwo}{}\renewcommand{\rhmcsBaTwinMeanRejTwo}{the random forest on the full predictor set}  % the alternatives the two-sided StepM rejects at 5 % (none: none)
\providecommand{\rhmcsBaTwinMeanPRidge}{}\renewcommand{\rhmcsBaTwinMeanPRidge}{0.13}  % StepM-adjusted p-value (one-sided), Ridge
\providecommand{\rhmcsBaTwinMeanRawRidge}{}\renewcommand{\rhmcsBaTwinMeanRawRidge}{0.083}  % single-test bootstrap p-value (one-sided), Ridge
\providecommand{\rhmcsBaTwinMeanPTwoRidge}{}\renewcommand{\rhmcsBaTwinMeanPTwoRidge}{0.26}  % StepM-adjusted p-value (two-sided), Ridge
\providecommand{\rhmcsBaTwinMeanRawTwoRidge}{}\renewcommand{\rhmcsBaTwinMeanRawTwoRidge}{0.16}  % single-test bootstrap p-value (two-sided), Ridge
\providecommand{\rhmcsBaTwinMeanPLasso}{}\renewcommand{\rhmcsBaTwinMeanPLasso}{0.11}  % StepM-adjusted p-value (one-sided), Lasso
\providecommand{\rhmcsBaTwinMeanRawLasso}{}\renewcommand{\rhmcsBaTwinMeanRawLasso}{0.063}  % single-test bootstrap p-value (one-sided), Lasso
\providecommand{\rhmcsBaTwinMeanPTwoLasso}{}\renewcommand{\rhmcsBaTwinMeanPTwoLasso}{0.21}  % StepM-adjusted p-value (two-sided), Lasso
\providecommand{\rhmcsBaTwinMeanRawTwoLasso}{}\renewcommand{\rhmcsBaTwinMeanRawTwoLasso}{0.12}  % single-test bootstrap p-value (two-sided), Lasso
\providecommand{\rhmcsBaTwinMeanPEnet}{}\renewcommand{\rhmcsBaTwinMeanPEnet}{0.069}  % StepM-adjusted p-value (one-sided), Elastic net
\providecommand{\rhmcsBaTwinMeanRawEnet}{}\renewcommand{\rhmcsBaTwinMeanRawEnet}{0.036}  % single-test bootstrap p-value (one-sided), Elastic net
\providecommand{\rhmcsBaTwinMeanPTwoEnet}{}\renewcommand{\rhmcsBaTwinMeanPTwoEnet}{0.13}  % StepM-adjusted p-value (two-sided), Elastic net
\providecommand{\rhmcsBaTwinMeanRawTwoEnet}{}\renewcommand{\rhmcsBaTwinMeanRawTwoEnet}{0.063}  % single-test bootstrap p-value (two-sided), Elastic net
\providecommand{\rhmcsBaTwinMeanPLgbm}{}\renewcommand{\rhmcsBaTwinMeanPLgbm}{0.061}  % StepM-adjusted p-value (one-sided), LightGBM
\providecommand{\rhmcsBaTwinMeanRawLgbm}{}\renewcommand{\rhmcsBaTwinMeanRawLgbm}{0.025}  % single-test bootstrap p-value (one-sided), LightGBM
\providecommand{\rhmcsBaTwinMeanPTwoLgbm}{}\renewcommand{\rhmcsBaTwinMeanPTwoLgbm}{0.12}  % StepM-adjusted p-value (two-sided), LightGBM
\providecommand{\rhmcsBaTwinMeanRawTwoLgbm}{}\renewcommand{\rhmcsBaTwinMeanRawTwoLgbm}{0.048}  % single-test bootstrap p-value (two-sided), LightGBM
\providecommand{\rhmcsBaTwinMeanPXgb}{}\renewcommand{\rhmcsBaTwinMeanPXgb}{0.13}  % StepM-adjusted p-value (one-sided), XGBoost
\providecommand{\rhmcsBaTwinMeanRawXgb}{}\renewcommand{\rhmcsBaTwinMeanRawXgb}{0.103}  % single-test bootstrap p-value (one-sided), XGBoost
\providecommand{\rhmcsBaTwinMeanPTwoXgb}{}\renewcommand{\rhmcsBaTwinMeanPTwoXgb}{0.26}  % StepM-adjusted p-value (two-sided), XGBoost
\providecommand{\rhmcsBaTwinMeanRawTwoXgb}{}\renewcommand{\rhmcsBaTwinMeanRawTwoXgb}{0.20}  % single-test bootstrap p-value (two-sided), XGBoost
\providecommand{\rhmcsBaTwinMeanPRf}{}\renewcommand{\rhmcsBaTwinMeanPRf}{0.028}  % StepM-adjusted p-value (one-sided), Random forest
\providecommand{\rhmcsBaTwinMeanRawRf}{}\renewcommand{\rhmcsBaTwinMeanRawRf}{0.013}  % single-test bootstrap p-value (one-sided), Random forest
\providecommand{\rhmcsBaTwinMeanPTwoRf}{}\renewcommand{\rhmcsBaTwinMeanPTwoRf}{0.042}  % StepM-adjusted p-value (two-sided), Random forest
\providecommand{\rhmcsBaTwinMeanRawTwoRf}{}\renewcommand{\rhmcsBaTwinMeanRawTwoRf}{0.014}  % single-test bootstrap p-value (two-sided), Random forest
\providecommand{\rhmcsBaTwinMeanNRejLo}{}\renewcommand{\rhmcsBaTwinMeanNRejLo}{1}  % StepM rejections at 5 %, block length 5
\providecommand{\rhmcsBaTwinMeanSpaCLo}{}\renewcommand{\rhmcsBaTwinMeanSpaCLo}{0.036}  % SPA p-value (consistent), block length 5
\providecommand{\rhmcsBaTwinMeanNRejTwoLo}{}\renewcommand{\rhmcsBaTwinMeanNRejTwoLo}{0}  % ... two-sided, block length 5
\providecommand{\rhmcsBaTwinMeanPMinLo}{}\renewcommand{\rhmcsBaTwinMeanPMinLo}{0.036}  % smallest adjusted p, block length 5
\providecommand{\rhmcsBaTwinMeanPTwoMinLo}{}\renewcommand{\rhmcsBaTwinMeanPTwoMinLo}{0.054}  % ... two-sided, block length 5
\providecommand{\rhmcsBaTwinMeanNRejHi}{}\renewcommand{\rhmcsBaTwinMeanNRejHi}{1}  % StepM rejections at 5 %, block length 20
\providecommand{\rhmcsBaTwinMeanSpaCHi}{}\renewcommand{\rhmcsBaTwinMeanSpaCHi}{0.034}  % SPA p-value (consistent), block length 20
\providecommand{\rhmcsBaTwinMeanNRejTwoHi}{}\renewcommand{\rhmcsBaTwinMeanNRejTwoHi}{1}  % ... two-sided, block length 20
\providecommand{\rhmcsBaTwinMeanPMinHi}{}\renewcommand{\rhmcsBaTwinMeanPMinHi}{0.034}  % smallest adjusted p, block length 20
\providecommand{\rhmcsBaTwinMeanPTwoMinHi}{}\renewcommand{\rhmcsBaTwinMeanPTwoMinHi}{0.048}  % ... two-sided, block length 20
\providecommand{\rhmcsBaTwinSrSpaC}{}\renewcommand{\rhmcsBaTwinSrSpaC}{0.008}  % SPA p-value (consistent), family BaTwinSr
\providecommand{\rhmcsBaTwinSrSpaL}{}\renewcommand{\rhmcsBaTwinSrSpaL}{0.008}  % ... lower
\providecommand{\rhmcsBaTwinSrSpaU}{}\renewcommand{\rhmcsBaTwinSrSpaU}{0.008}  % ... upper
\providecommand{\rhmcsBaTwinSrNRej}{}\renewcommand{\rhmcsBaTwinSrNRej}{1}  % StepM rejections at FWER 5 % (of 6)
\providecommand{\rhmcsBaTwinSrNRejWord}{}\renewcommand{\rhmcsBaTwinSrNRejWord}{one}  % ... as a word
\providecommand{\rhmcsBaTwinSrNRaw}{}\renewcommand{\rhmcsBaTwinSrNRaw}{3}  % single-test bootstrap p < 0.05 (no adjustment)
\providecommand{\rhmcsBaTwinSrNRawWord}{}\renewcommand{\rhmcsBaTwinSrNRawWord}{three}  % ... as a word
\providecommand{\rhmcsBaTwinSrNPos}{}\renewcommand{\rhmcsBaTwinSrNPos}{6}  % alternatives with a positive estimate (better than the benchmark)
\providecommand{\rhmcsBaTwinSrPMin}{}\renewcommand{\rhmcsBaTwinSrPMin}{0.008}  % smallest StepM-adjusted p-value
\providecommand{\rhmcsBaTwinSrPMax}{}\renewcommand{\rhmcsBaTwinSrPMax}{0.14}  % largest StepM-adjusted p-value
\providecommand{\rhmcsBaTwinSrPTwoMin}{}\renewcommand{\rhmcsBaTwinSrPTwoMin}{0.013}  % smallest two-sided StepM-adjusted p-value
\providecommand{\rhmcsBaTwinSrPTwoMax}{}\renewcommand{\rhmcsBaTwinSrPTwoMax}{0.27}  % largest two-sided StepM-adjusted p-value
\providecommand{\rhmcsBaTwinSrRawMin}{}\renewcommand{\rhmcsBaTwinSrRawMin}{0.002}  % smallest single-test bootstrap p-value
\providecommand{\rhmcsBaTwinSrRej}{}\renewcommand{\rhmcsBaTwinSrRej}{the random forest on the full predictor set}  % the alternatives StepM rejects at 5 % (none: none)
\providecommand{\rhmcsBaTwinSrNRejTwo}{}\renewcommand{\rhmcsBaTwinSrNRejTwo}{1}  % StepM rejections at FWER 5 %, two-sided
\providecommand{\rhmcsBaTwinSrNRejTwoWord}{}\renewcommand{\rhmcsBaTwinSrNRejTwoWord}{one}  % ... as a word
\providecommand{\rhmcsBaTwinSrNRawTwo}{}\renewcommand{\rhmcsBaTwinSrNRawTwo}{2}  % single-test bootstrap p < 0.05, two-sided
\providecommand{\rhmcsBaTwinSrRejTwo}{}\renewcommand{\rhmcsBaTwinSrRejTwo}{the random forest on the full predictor set}  % the alternatives the two-sided StepM rejects at 5 % (none: none)
\providecommand{\rhmcsBaTwinSrPRidge}{}\renewcommand{\rhmcsBaTwinSrPRidge}{0.14}  % StepM-adjusted p-value (one-sided), Ridge
\providecommand{\rhmcsBaTwinSrRawRidge}{}\renewcommand{\rhmcsBaTwinSrRawRidge}{0.089}  % single-test bootstrap p-value (one-sided), Ridge
\providecommand{\rhmcsBaTwinSrPTwoRidge}{}\renewcommand{\rhmcsBaTwinSrPTwoRidge}{0.27}  % StepM-adjusted p-value (two-sided), Ridge
\providecommand{\rhmcsBaTwinSrRawTwoRidge}{}\renewcommand{\rhmcsBaTwinSrRawTwoRidge}{0.17}  % single-test bootstrap p-value (two-sided), Ridge
\providecommand{\rhmcsBaTwinSrPLasso}{}\renewcommand{\rhmcsBaTwinSrPLasso}{0.12}  % StepM-adjusted p-value (one-sided), Lasso
\providecommand{\rhmcsBaTwinSrRawLasso}{}\renewcommand{\rhmcsBaTwinSrRawLasso}{0.071}  % single-test bootstrap p-value (one-sided), Lasso
\providecommand{\rhmcsBaTwinSrPTwoLasso}{}\renewcommand{\rhmcsBaTwinSrPTwoLasso}{0.22}  % StepM-adjusted p-value (two-sided), Lasso
\providecommand{\rhmcsBaTwinSrRawTwoLasso}{}\renewcommand{\rhmcsBaTwinSrRawTwoLasso}{0.12}  % single-test bootstrap p-value (two-sided), Lasso
\providecommand{\rhmcsBaTwinSrPEnet}{}\renewcommand{\rhmcsBaTwinSrPEnet}{0.080}  % StepM-adjusted p-value (one-sided), Elastic net
\providecommand{\rhmcsBaTwinSrRawEnet}{}\renewcommand{\rhmcsBaTwinSrRawEnet}{0.043}  % single-test bootstrap p-value (one-sided), Elastic net
\providecommand{\rhmcsBaTwinSrPTwoEnet}{}\renewcommand{\rhmcsBaTwinSrPTwoEnet}{0.14}  % StepM-adjusted p-value (two-sided), Elastic net
\providecommand{\rhmcsBaTwinSrRawTwoEnet}{}\renewcommand{\rhmcsBaTwinSrRawTwoEnet}{0.070}  % single-test bootstrap p-value (two-sided), Elastic net
\providecommand{\rhmcsBaTwinSrPLgbm}{}\renewcommand{\rhmcsBaTwinSrPLgbm}{0.066}  % StepM-adjusted p-value (one-sided), LightGBM
\providecommand{\rhmcsBaTwinSrRawLgbm}{}\renewcommand{\rhmcsBaTwinSrRawLgbm}{0.027}  % single-test bootstrap p-value (one-sided), LightGBM
\providecommand{\rhmcsBaTwinSrPTwoLgbm}{}\renewcommand{\rhmcsBaTwinSrPTwoLgbm}{0.12}  % StepM-adjusted p-value (two-sided), LightGBM
\providecommand{\rhmcsBaTwinSrRawTwoLgbm}{}\renewcommand{\rhmcsBaTwinSrRawTwoLgbm}{0.050}  % single-test bootstrap p-value (two-sided), LightGBM
\providecommand{\rhmcsBaTwinSrPXgb}{}\renewcommand{\rhmcsBaTwinSrPXgb}{0.14}  % StepM-adjusted p-value (one-sided), XGBoost
\providecommand{\rhmcsBaTwinSrRawXgb}{}\renewcommand{\rhmcsBaTwinSrRawXgb}{0.11}  % single-test bootstrap p-value (one-sided), XGBoost
\providecommand{\rhmcsBaTwinSrPTwoXgb}{}\renewcommand{\rhmcsBaTwinSrPTwoXgb}{0.27}  % StepM-adjusted p-value (two-sided), XGBoost
\providecommand{\rhmcsBaTwinSrRawTwoXgb}{}\renewcommand{\rhmcsBaTwinSrRawTwoXgb}{0.21}  % single-test bootstrap p-value (two-sided), XGBoost
\providecommand{\rhmcsBaTwinSrPRf}{}\renewcommand{\rhmcsBaTwinSrPRf}{0.008}  % StepM-adjusted p-value (one-sided), Random forest
\providecommand{\rhmcsBaTwinSrRawRf}{}\renewcommand{\rhmcsBaTwinSrRawRf}{0.002}  % single-test bootstrap p-value (one-sided), Random forest
\providecommand{\rhmcsBaTwinSrPTwoRf}{}\renewcommand{\rhmcsBaTwinSrPTwoRf}{0.013}  % StepM-adjusted p-value (two-sided), Random forest
\providecommand{\rhmcsBaTwinSrRawTwoRf}{}\renewcommand{\rhmcsBaTwinSrRawTwoRf}{0.003}  % single-test bootstrap p-value (two-sided), Random forest
\providecommand{\rhmcsBaTwinSrNRejLo}{}\renewcommand{\rhmcsBaTwinSrNRejLo}{1}  % StepM rejections at 5 %, block length 5
\providecommand{\rhmcsBaTwinSrSpaCLo}{}\renewcommand{\rhmcsBaTwinSrSpaCLo}{0.011}  % SPA p-value (consistent), block length 5
\providecommand{\rhmcsBaTwinSrNRejTwoLo}{}\renewcommand{\rhmcsBaTwinSrNRejTwoLo}{1}  % ... two-sided, block length 5
\providecommand{\rhmcsBaTwinSrPMinLo}{}\renewcommand{\rhmcsBaTwinSrPMinLo}{0.011}  % smallest adjusted p, block length 5
\providecommand{\rhmcsBaTwinSrPTwoMinLo}{}\renewcommand{\rhmcsBaTwinSrPTwoMinLo}{0.017}  % ... two-sided, block length 5
\providecommand{\rhmcsBaTwinSrNRejHi}{}\renewcommand{\rhmcsBaTwinSrNRejHi}{1}  % StepM rejections at 5 %, block length 20
\providecommand{\rhmcsBaTwinSrSpaCHi}{}\renewcommand{\rhmcsBaTwinSrSpaCHi}{0.012}  % SPA p-value (consistent), block length 20
\providecommand{\rhmcsBaTwinSrNRejTwoHi}{}\renewcommand{\rhmcsBaTwinSrNRejTwoHi}{1}  % ... two-sided, block length 20
\providecommand{\rhmcsBaTwinSrPMinHi}{}\renewcommand{\rhmcsBaTwinSrPMinHi}{0.012}  % smallest adjusted p, block length 20
\providecommand{\rhmcsBaTwinSrPTwoMinHi}{}\renewcommand{\rhmcsBaTwinSrPTwoMinHi}{0.018}  % ... two-sided, block length 20
\providecommand{\rhmcsTwinTwoSidedPRidge}{}\renewcommand{\rhmcsTwinTwoSidedPRidge}{0.17}  % two-sided pooled-SR bootstrap p of the Full - HAR Sharpe difference, Ridge (D19: the check p_bt; the paper's test is Ledoit-Wolf)
\providecommand{\rhmcsTwinTwoSidedPLasso}{}\renewcommand{\rhmcsTwinTwoSidedPLasso}{0.12}  % two-sided pooled-SR bootstrap p of the Full - HAR Sharpe difference, Lasso (D19: the check p_bt; the paper's test is Ledoit-Wolf)
\providecommand{\rhmcsTwinTwoSidedPEnet}{}\renewcommand{\rhmcsTwinTwoSidedPEnet}{0.072}  % two-sided pooled-SR bootstrap p of the Full - HAR Sharpe difference, Elastic net (D19: the check p_bt; the paper's test is Ledoit-Wolf)
\providecommand{\rhmcsTwinTwoSidedPLgbm}{}\renewcommand{\rhmcsTwinTwoSidedPLgbm}{0.050}  % two-sided pooled-SR bootstrap p of the Full - HAR Sharpe difference, LightGBM (D19: the check p_bt; the paper's test is Ledoit-Wolf)
\providecommand{\rhmcsTwinTwoSidedPXgb}{}\renewcommand{\rhmcsTwinTwoSidedPXgb}{0.20}  % two-sided pooled-SR bootstrap p of the Full - HAR Sharpe difference, XGBoost (D19: the check p_bt; the paper's test is Ledoit-Wolf)
\providecommand{\rhmcsTwinTwoSidedPRf}{}\renewcommand{\rhmcsTwinTwoSidedPRf}{0.003}  % two-sided pooled-SR bootstrap p of the Full - HAR Sharpe difference, Random forest (D19: the check p_bt; the paper's test is Ledoit-Wolf)
 in the preamble of paper_rehaul.tex; % AUTO-GENERATED by writeup/make_rehaul_mcs.py -- do not edit; re-run the script.
% tab:rh_mcs: MCS p-values (T_max and T_R) of the pooled intraday forecasts and the closing-interval HAR
% benchmark; circular block bootstrap, daily losses (every-bar losses averaged to the day).
\begin{table}[htbp]
\centering
\small
\setlength{\tabcolsep}{4pt}
\begin{threeparttable}
\caption{Model confidence sets of the pooled intraday forecasts}
\label{tab:rh_mcs}
\begin{tabular*}{\linewidth}{@{\extracolsep{\fill}}lrrrrrr@{}}
\toprule
 & \multicolumn{2}{c}{Every bar, 2014--2024} & \multicolumn{2}{c}{16:00 bar, 2014--2024} & \multicolumn{2}{c}{16:00 bar, $\mathcal T_{\mathrm{oos}}$} \\
\cmidrule(lr){2-3}\cmidrule(lr){4-5}\cmidrule(lr){6-7}
Forecast & $T_{\max}$ & $T_R$ & $T_{\max}$ & $T_R$ & $T_{\max}$ & $T_R$ \\
\midrule
\multicolumn{7}{@{}l}{\textit{Panel A: QLIKE}} \\
\addlinespace[2pt]
Linear, HAR$^{a}$ & \ensuremath{<}0.001 & \ensuremath{<}0.001 & \textit{0.126} & \ensuremath{<}0.001 & 0.052 & 0.061 \\
Ridge, Full & \ensuremath{<}0.001 & \ensuremath{<}0.001 & \textit{0.159} & 0.015 & 0.011 & 0.008 \\
Lasso, Full & \ensuremath{<}0.001 & \ensuremath{<}0.001 & \textit{0.126} & 0.011 & 0.011 & 0.011 \\
Elastic net, Full & \ensuremath{<}0.001 & \ensuremath{<}0.001 & \textit{0.246} & 0.030 & 0.036 & 0.016 \\
LightGBM, HAR & \ensuremath{<}0.001 & \ensuremath{<}0.001 & \textit{0.246} & \textit{0.103} & \textit{0.246} & \textit{0.116} \\
LightGBM, Full & \textbf{1.000} & \textbf{1.000} & \textit{0.246} & \textit{0.159} & \textit{0.246} & \textit{0.138} \\
XGBoost, HAR & \ensuremath{<}0.001 & \ensuremath{<}0.001 & \textit{0.246} & \textit{0.103} & \textit{0.246} & \textit{0.116} \\
XGBoost, Full & 0.003 & 0.003 & \textbf{1.000} & \textbf{1.000} & \textbf{0.256} & \textbf{0.256} \\
Random forest, HAR & \ensuremath{<}0.001 & \ensuremath{<}0.001 & \textit{0.159} & \ensuremath{<}0.001 & 0.052 & 0.052 \\
Random forest, Full & \ensuremath{<}0.001 & \ensuremath{<}0.001 & \textit{0.182} & \ensuremath{<}0.001 & \textit{0.129} & 0.016 \\
Closing-interval HAR benchmark$^{b}$ & -- & -- & -- & -- & \textbf{1.000} & \textbf{1.000} \\
\addlinespace[4pt]
\multicolumn{7}{@{}l}{\textit{Panel B: Squared error}} \\
\addlinespace[2pt]
Linear, HAR$^{a}$ & \textbf{1.000} & \textbf{1.000} & \textbf{0.572} & \textbf{0.653} & \textbf{0.537} & \textbf{0.613} \\
Ridge, Full & \textbf{0.489} & \textbf{0.415} & \textbf{0.437} & \textbf{0.653} & \textbf{0.446} & \textbf{0.613} \\
Lasso, Full & \textbf{0.334} & \textbf{0.415} & \textbf{0.371} & \textbf{0.653} & \textbf{0.383} & \textbf{0.613} \\
Elastic net, Full & \textbf{0.422} & \textbf{0.415} & \textbf{0.488} & \textbf{0.653} & \textbf{0.496} & \textbf{0.613} \\
LightGBM, HAR & \textbf{0.489} & \textbf{0.310} & \textbf{0.572} & \textbf{0.369} & \textbf{0.537} & \textbf{0.561} \\
LightGBM, Full & \textbf{0.489} & \textbf{0.415} & \textbf{0.488} & \textbf{0.653} & \textbf{0.537} & \textbf{0.613} \\
XGBoost, HAR & \textbf{0.489} & \textbf{0.415} & \textbf{0.572} & \textbf{0.653} & \textbf{0.537} & \textbf{0.613} \\
XGBoost, Full & \textbf{0.489} & \textbf{0.415} & \textbf{0.488} & \textbf{0.653} & \textbf{0.537} & \textbf{0.613} \\
Random forest, HAR & \textbf{0.489} & \textbf{0.310} & \textbf{1.000} & \textbf{1.000} & \textbf{1.000} & \textbf{1.000} \\
Random forest, Full & \textbf{0.489} & \textbf{0.415} & \textbf{0.558} & \textbf{0.653} & \textbf{0.537} & \textbf{0.613} \\
Closing-interval HAR benchmark$^{b}$ & -- & -- & -- & -- & \textbf{0.537} & \textbf{0.572} \\
\bottomrule
\end{tabular*}
\begin{tablenotes}[flushleft]
\footnotesize
\item \textit{Notes:} MCS $p$-values \citep{HansenLundeNason2011} of the pooled intraday forecasts of Table~\ref{tab:rh_oos} and, on $\mathcal T_{\mathrm{oos}}$, of the closing-interval HAR benchmark. A forecast is in the 90\% (75\%) model confidence set when its $p$-value is at least 0.10 (0.25); \textbf{bold}: in the 75\% set, \textit{italic}: in the 90\% set only. $T_{\max}$ and $T_R$: the two statistics of the elimination rule. Losses are daily (on every bar, the mean loss of the day's \rhorNBars{} bars), so the bootstrap resamples trading days: the \rhmcsDaysAll{} days from \rhorSpanStart{} to \rhorSpanEnd{} in the first four columns and the \rhmcsDaysOos{} days of $\mathcal T_{\mathrm{oos}}$ in the last two. Circular moving-block bootstrap, blocks of \rhmcsBlock{} trading days, \rhmcsB{} replications, seed \rhmcsSeed; the same resamples at every elimination step. $^{a}$The pooled ridge, lasso and elastic net on the HAR benchmark reproduce the pooled OLS (their penalties are numerically zero), so they enter as one model, the OLS forecast. $^{b}$The closing-interval HAR benchmark has forecasts from \rhmcsChStart{} only, so it is not in the 2014--2024 sets.
\end{tablenotes}
\end{threeparttable}
\end{table}

% tab:rh_mcs_spa: SPA p-values and StepM-adjusted p-values (single-test bootstrap p in parentheses).
\begin{table}[htbp]
\centering
\small
\setlength{\tabcolsep}{3pt}
\begin{threeparttable}
\caption{Superior predictive ability and StepM tests of the full predictor set}
\label{tab:rh_mcs_spa}
\begin{tabular*}{\linewidth}{@{\extracolsep{\fill}}lrrrrrrr@{}}
\toprule
 & SPA & \multicolumn{6}{c}{StepM, by model class} \\
\cmidrule(lr){3-8}
Comparison & $p$ & Ridge & Lasso & Elastic net & LightGBM & XGBoost & Random forest \\
\midrule
\multicolumn{8}{@{}l}{\textit{Panel A: QLIKE, full predictor set against the HAR benchmark of the same class}} \\
\addlinespace[2pt]
Every bar, 2014--2024 & \ensuremath{<}0.001 & \textbf{\ensuremath{<}0.001} & \textbf{\ensuremath{<}0.001} & \textbf{\ensuremath{<}0.001} & \textbf{\ensuremath{<}0.001} & \textbf{\ensuremath{<}0.001} & \textbf{\ensuremath{<}0.001} \\
 &  & (\ensuremath{<}0.001) & (\ensuremath{<}0.001) & (\ensuremath{<}0.001) & (\ensuremath{<}0.001) & (\ensuremath{<}0.001) & (\ensuremath{<}0.001) \\
16:00 bar, 2014--2024 & 0.033 & 0.675 & 0.842 & 0.499 & 0.072 & \textbf{0.033} & 0.302 \\
 &  & (0.646) & (0.842) & (0.437) & (0.034) & (0.013) & (0.177) \\
16:00 bar, $\mathcal T_{\mathrm{oos}}$ & 0.380 & 0.987 & 0.987 & 0.973 & 0.581 & 0.415 & 0.599 \\
 &  & (0.984) & (0.986) & (0.961) & (0.345) & (0.219) & (0.402) \\
\addlinespace[4pt]
\multicolumn{8}{@{}l}{\textit{Panel B: QLIKE, full predictor set against the closing-interval HAR benchmark}} \\
\addlinespace[2pt]
16:00 bar, $\mathcal T_{\mathrm{oos}}$ & 0.529 & 1.000 & 1.000 & 0.999 & 0.972 & 0.959 & 0.995 \\
 &  & (0.999) & (0.999) & (0.999) & (0.905) & (0.867) & (0.974) \\
\addlinespace[4pt]
\multicolumn{8}{@{}l}{\textit{Panel C: sign strategies against the unconditional short straddle}} \\
\addlinespace[2pt]
Mean return, mid-quotes & 0.024 & \textbf{0.035} & \textbf{0.035} & \textbf{0.030} & \textbf{0.043} & 0.050 & \textbf{0.024} \\
 &  & (0.023) & (0.025) & (0.020) & (0.038) & (0.050) & (0.014) \\
Mean return, bid/ask & 0.024 & \textbf{0.034} & \textbf{0.034} & \textbf{0.029} & \textbf{0.042} & 0.050 & \textbf{0.024} \\
 &  & (0.023) & (0.024) & (0.019) & (0.038) & (0.050) & (0.014) \\
Sharpe ratio, mid-quotes & 0.002 & \textbf{0.002} & \textbf{0.003} & \textbf{0.002} & \textbf{0.002} & \textbf{0.003} & \textbf{0.002} \\
 &  & (\ensuremath{<}0.001) & (0.003) & (0.002) & (\ensuremath{<}0.001) & (0.002) & (0.001) \\
Sharpe ratio, bid/ask & 0.002 & \textbf{0.003} & \textbf{0.004} & \textbf{0.003} & \textbf{0.002} & \textbf{0.004} & \textbf{0.002} \\
 &  & (\ensuremath{<}0.001) & (0.003) & (0.002) & (0.001) & (0.003) & (0.001) \\
\addlinespace[4pt]
\multicolumn{8}{@{}l}{\textit{Panel D: sign strategies against the HAR-benchmark strategy of the same class}} \\
\addlinespace[2pt]
Mean return, mid-quotes & 0.026 & 0.087 & 0.093 & 0.072 & \textbf{0.048} & 0.138 & \textbf{0.026} \\
 &  & (0.049) & (0.059) & (0.037) & (0.021) & (0.138) & (0.012) \\
Mean return, bid/ask & 0.026 & 0.087 & 0.093 & 0.073 & \textbf{0.048} & 0.138 & \textbf{0.026} \\
 &  & (0.049) & (0.059) & (0.037) & (0.021) & (0.138) & (0.012) \\
Sharpe ratio, mid-quotes & 0.022 & 0.175 & 0.175 & 0.175 & 0.096 & 0.175 & \textbf{0.022} \\
 &  & (0.118) & (0.150) & (0.101) & (0.047) & (0.100) & (0.003) \\
Sharpe ratio, bid/ask & 0.022 & 0.169 & 0.169 & 0.169 & 0.088 & 0.169 & \textbf{0.022} \\
 &  & (0.110) & (0.145) & (0.097) & (0.044) & (0.097) & (0.003) \\
\bottomrule
\end{tabular*}
\begin{tablenotes}[flushleft]
\footnotesize
\item \textit{Notes:} Each row is one family of \rhmcsNAlt{} one-sided hypotheses $H_k$: the full-set forecast (strategy) of class $k$ is no better than its comparison. SPA: the consistent $p$-value of the test of superior predictive ability of \citet{Hansen2005}, $H_0$: no full-set forecast in the row beats its comparison. StepM: $p$-values adjusted for the \rhmcsNAlt{} comparisons of the row by the studentized step-down procedure of \citet{RomanoWolf2005}, which controls the familywise error rate; \textbf{bold}: below 0.05. In parentheses: the one-sided bootstrap $p$-value of the comparison alone, without the adjustment. Panels A and B: daily QLIKE loss differentials (every bar: the mean over the day's \rhorNBars{} bars), circular moving-block bootstrap with blocks of \rhmcsBlock{} trading days, \rhmcsB{} replications, seed \rhmcsSeed. Panels C and D: daily returns of the sign strategy (Section~\ref{sec:rh_econ_strategy}) on the \rhmcsNPw{} post-warmup days, the mean daily return differential and the difference of the Sharpe ratios of Section~\ref{sec:rh_econ_measures}, with the paired circular block bootstrap of Appendix~\ref{sec:rh_methods_tests} (blocks of \rhmcsBlock{} days, \rhmcsB{} replications, seed \rhmcsSeedEcon), each statistic studentized by its bootstrap standard deviation. The single-test $p$-values in parentheses are those of this bootstrap, one-sided; they are not the two-sided Diebold--Mariano and fixed-$b$ $p$-values of Sections~\ref{sec:rh_results} and~\ref{sec:rh_econ}.
\end{tablenotes}
\end{threeparttable}
\end{table}

%          in rehaul/10_appendix_floats.tex (or 06_results_floats.tex). Part A replaces the item "Multiple-comparison
%          procedures" of rehaul/10_statistical_tests.tex and closes note E14; Part B goes at the end of
%          sec:rh_results_classes (06_results.tex); Part C after the item "Full predictor set versus the HAR benchmark"
%          of sec:rh_econ_results (07_economics.tex); Part D lists the sentences of the current text that the
%          procedures qualify.
% Every number below is a macro (\rhmcs..., or an existing \rhr... / \rhe... / \rhor... macro); only the levels (5%,
% 90%, 75%) are typed, as elsewhere in the paper. The words "only", "alone", "all", "every", "none", "no" and the lists
% of forecasts that do or do not survive are asserted in make_rehaul_mcs.py (draft_claims), so a rerun whose data no
% longer bear them out stops.

% =====================================================================================================================
% Part A. Appendix sec:rh_methods_tests (10_statistical_tests.tex): replaces the item "Multiple-comparison procedures"
% and its note E14.
% =====================================================================================================================
\item \textbf{Multiple-comparison procedures.} Each test above compares two forecasts, or two strategies, at a time,
  and Sections~\ref{sec:rh_results} and~\ref{sec:rh_econ} make many such comparisons. Three procedures account for
  their number; Tables~\ref{tab:rh_mcs} and~\ref{tab:rh_mcs_spa} report them.
  \begin{itemize}
  \item \emph{Model confidence set} \citep{HansenLundeNason2011}. Starting from a set $\mathcal M$ of forecasts with
    daily losses $L_{i,t}$, each step tests the equal predictive ability of the forecasts in $\mathcal M$ and, if it
    rejects, removes one of them. With $\bar d_{ij}$ the average of $L_{i,t} - L_{j,t}$ over $\mathcal O$ and
    $\bar d_{i\cdot}$ the average loss of forecast $i$ minus the average over $\mathcal M$, the two statistics are
    \[
      T_{\max} = \max_{i \in \mathcal M} t_{i\cdot}, \quad
      t_{i\cdot} = \frac{\bar d_{i\cdot}}{\sqrt{\widehat{\mathrm{var}}(\bar d_{i\cdot})}},
      \qquad
      T_R = \max_{i, j \in \mathcal M} |t_{ij}|, \quad
      t_{ij} = \frac{\bar d_{ij}}{\sqrt{\widehat{\mathrm{var}}(\bar d_{ij})}},
    \]
    and their elimination rules remove $\arg\max_i t_{i\cdot}$ and $\arg\max_i \max_j t_{ij}$, respectively. The
    variances and the null distributions of $T_{\max}$ and $T_R$ come from the bootstrap below, with the statistics
    recentred at their sample values. The elimination continues to the last forecast, so every forecast has an MCS
    $p$-value: the largest step $p$-value up to its removal, and one for the last forecast. The $(1-\alpha)$ model
    confidence set holds the forecasts whose $p$-value is at least $\alpha$; we report the 90\% and 75\% sets.
  \item \emph{Superior predictive ability} \citep{Hansen2005}. For a family of $K$ comparisons with daily
    differentials $d_{k,t}$, positive when the alternative $k$ is better than its comparison (the loss of the
    comparison minus that of the alternative, or the return of the alternative minus that of the comparison), the
    test of $H_0\colon \max_k \E[d_{k,t}] \le 0$ uses $T^{\mathrm{SPA}} = \max\bigl(0, \max_k \bar d_k/\hat\sigma_k\bigr)$,
    with $\hat\sigma_k$ the bootstrap standard deviation of $\bar d_k$. We report its consistent $p$-value.
  \item \emph{StepM} \citep{RomanoWolf2005}. The studentized step-down procedure tests the $K$ hypotheses
    $H_k\colon \E[d_{k,t}] \le 0$ and controls the familywise error rate, the probability of rejecting at least one
    true $H_k$. With $t_k = \bar d_k/\hat\sigma_k$ ordered as $t_{(1)} \ge \dots \ge t_{(K)}$ and
    $Z^{*}_k = (\bar d^{*}_k - \bar d_k)/\hat\sigma_k$ in a bootstrap sample, the adjusted $p$-value of the $j$-th
    hypothesis is
    \[
      \tilde p_{(j)} = \max_{i \le j} \Pr\nolimits^{*}\Bigl(\max_{l \ge i} Z^{*}_{(l)} \ge t_{(i)}\Bigr),
    \]
    and StepM at level $\alpha$ rejects the $H_k$ with $\tilde p_k < \alpha$. The two-sided version, for
    $H_k\colon \E[d_{k,t}] = 0$, replaces $t_k$ and $Z^{*}_k$ by their absolute values. Beside each adjusted $p$-value
    we report the bootstrap $p$-value of the comparison alone, $\Pr^{*}(Z^{*}_k \ge t_k)$.
  \item \emph{Sharpe ratios.} For a difference of the Sharpe ratios of Section~\ref{sec:rh_econ_measures},
    $\Delta\SR_k$ takes the place of $\bar d_k$ and $\Delta\SR^{*}_k$, computed in each bootstrap sample as in the
    test of equal Sharpe ratios above, the place of $\bar d^{*}_k$. For a single comparison the two-sided bootstrap
    $p$-value is then the $p$-value of that test.
  \item \emph{Bootstrap.} The three procedures resample trading days with the circular moving-block bootstrap: blocks
    of \rhmcsBlock{} consecutive days, \rhmcsB{} samples, the same samples for every forecast of a comparison and at
    every elimination step. The forecast losses use seed \rhmcsSeed; the strategy returns use the samples of the test
    of equal Sharpe ratios above (seed \rhmcsSeedEcon). The block length is fixed rather than calibrated; we also
    report the results at block lengths \rhmcsBlockLo{} and \rhmcsBlockHi.
  \item \emph{Losses and forecasts.} The losses are daily: for the forecasts of every regular-hours bar, the average
    loss of the day's \rhorNBars{} bars, as in the Diebold--Mariano tests of Table~\ref{tab:rh_oos}, so the bootstrap
    resamples days. The pooled ridge regression, lasso and elastic net on the HAR benchmark reproduce the pooled OLS
    (their penalties are numerically zero; Section~\ref{sec:rh_exp_setup}): their daily every-bar losses differ from
    those of the OLS by at most $\rhmcsLinGap$ in relative terms, so they enter the model confidence set as one
    forecast, the OLS. The set therefore holds \rhmcsNModelsWord{} distinct pooled forecasts, and
    \rhmcsNModelsBenchWord{} with the closing-interval HAR benchmark, which has forecasts from \rhmcsChStart{} only
    and so enters the sets on $\mathcal T_{\mathrm{oos}}$.
  \end{itemize}

% =====================================================================================================================
% Part B. Section 6 (06_results.tex), a new item at the end of sec:rh_results_classes.
% =====================================================================================================================
\item \textbf{Multiple comparisons.} The tests above compare two forecasts at a time. The model confidence set
  (Table~\ref{tab:rh_mcs}) and the SPA and StepM tests of the full predictor set (Table~\ref{tab:rh_mcs_spa},
  Panels~A and~B) account for the number of comparisons (Appendix~\ref{sec:rh_methods_tests}).
  \begin{itemize}
  \item \emph{Every bar.} Under QLIKE, the 90\% model confidence set holds a single forecast,
    \rhmcsBarQlikeTmaxSetNinety, under both statistics and at block lengths \rhmcsBlockLo, \rhmcsBlock{} and
    \rhmcsBlockHi; the next forecast, XGBoost on the full predictor set, has an MCS $p$-value of
    \rhmcsBarQlikeTmaxPXgbFull, and no forecast on the HAR benchmark is in the set. The reductions in QLIKE from the
    full predictor set survive the adjustment in all \rhorAfPairs{} model classes: the SPA $p$-value is
    \rhmcsBarQAfSpaC, and StepM rejects all \rhmcsNAlt{} hypotheses, with every adjusted $p$-value \rhmcsBarQAfPMax{}
    (and all \rhmcsNAlt{} rejections, one- and two-sided, at the three block lengths). Under squared error, by
    contrast, the 75\% set
    holds all \rhmcsBarMseTmaxNSevenFive{} forecasts, and StepM rejects none of the \rhmcsNAlt{} hypotheses (SPA
    $p$-value \rhmcsBarMAfSpaC).
  \item \emph{The closing half hour, \rhorSpanStart{} to \rhorSpanEnd.} On the \rhmcsDaysAll{} days of
    Table~\ref{tab:rh_oos}, the forecasts of the closing half hour are hard to tell apart: under $T_{\max}$ the 90\%
    set holds all \rhmcsAllQlikeTmaxNNinety{} pooled forecasts, and under $T_R$ it holds \rhmcsAllQlikeTrangeSetNinety;
    the 75\% set holds \rhmcsAllQlikeTmaxSetSevenFive{} alone under both statistics. Against the HAR benchmark of the
    same class, the SPA $p$-value is \rhmcsAllQAfSpaC, and StepM rejects for \rhmcsAllQAfRej{} alone (adjusted
    $p$-value \rhmcsAllQAfPXgb; LightGBM \rhmcsAllQAfPLgbm), and for no class in the two-sided test (smallest adjusted
    $p$-value \rhmcsAllQAfPTwoMin).
  \item \emph{The closing half hour on $\mathcal T_{\mathrm{oos}}$.} The closing-interval HAR benchmark is the last
    forecast left under both statistics (MCS $p$-value one). Under $T_R$ the 90\% set holds
    \rhmcsOosQlikeTrangeSetNinety; under $T_{\max}$ it also holds the random forest on the full predictor set (MCS
    $p$-value \rhmcsOosQlikeTmaxPRfFull, against \rhmcsOosQlikeTrangePRfFull{} under $T_R$). No linear forecast is in
    it: every forecast left out has an MCS $p$-value of at most \rhmcsOosQlikeTmaxPMaxOutNinety{} under $T_{\max}$ and
    \rhmcsOosQlikeTrangePMaxOutNinety{} under $T_R$. The 75\% set holds \rhmcsOosQlikeTmaxSetSevenFive{} under both
    statistics, and \rhmcsOosQlikeTmaxSetSevenFiveHi{} alone with blocks of \rhmcsBlockHi{} days.
    \begin{itemize}
    \item Against the HAR benchmark of the same class, no full-set forecast is more accurate after the adjustment
      (SPA $p$-value \rhmcsOosQAfSpaC; smallest StepM-adjusted $p$-value \rhmcsOosQAfPMin), in line with the
      Diebold--Mariano tests above. Of the increases in loss that are significant in those tests, those of ridge
      regression and the lasso, the lasso's survives the two-sided StepM (adjusted $p$-value \rhmcsOosQAfPTwoLasso)
      and that of ridge regression does not (\rhmcsOosQAfPTwoRidge).
    \item No full-set forecast is more accurate than the closing-interval HAR benchmark (SPA $p$-value
      \rhmcsOosQBmSpaC). Testing all \rhmcsNModelsWord{} distinct pooled forecasts against the benchmark at once, the
      two-sided StepM confirms a higher loss for \rhmcsOosQBmAllRejTwo, but not for LightGBM on the HAR benchmark
      (adjusted $p$-value \rhmcsOosQBmAllPTwoLgbmHar), XGBoost on the HAR benchmark (\rhmcsOosQBmAllPTwoXgbHar) or
      the random forest on the full predictor set (\rhmcsOosQBmAllPTwoRfFull), each significant in the pairwise test
      above.
    \item On the \rhmcsDaysCh{} days from \rhmcsChStart{} to \rhorSpanEnd{} on which the benchmark has forecasts, it
      is again the last forecast left, and the 90\% set under $T_R$ holds \rhmcsChQlikeTrangeSetNinety.
    \item Under squared error, every forecast is in the 75\% sets of the closing half hour
      (\rhmcsOosMseTmaxNSevenFive{} of \rhmcsNModelsBench{} on $\mathcal T_{\mathrm{oos}}$), and StepM rejects no
      hypothesis.
    \end{itemize}
  \end{itemize}

% =====================================================================================================================
% Part C. Section 7 (07_economics.tex), a new item after "Full predictor set versus the HAR benchmark" in
% sec:rh_econ_results.
% =====================================================================================================================
\item \textbf{Multiple comparisons of the strategies.} Table~\ref{tab:rh_mcs_spa} (Panels~C and~D) tests the
  \rhmcsNAlt{} sign strategies on the full predictor set jointly, on the \rheSzDays{} post-warmup days, against the
  alternative that the strategy is better, unless stated otherwise (Appendix~\ref{sec:rh_methods_tests}).
  \begin{itemize}
  \item \emph{Against the unconditional short straddle.} The Sharpe ratio gain survives the adjustment for every
    full-set strategy: the SPA $p$-value is \rhmcsMidShortSrSpaC{} at mid-quotes and \rhmcsBaShortSrSpaC{} at quoted
    bid/ask prices, and StepM rejects all \rhmcsNAlt{} hypotheses, one- and two-sided, at both prices and block
    lengths \rhmcsBlockLo, \rhmcsBlock{} and \rhmcsBlockHi{} (two-sided adjusted $p$-values at most
    \rhmcsMidShortSrPTwoMax{} at mid-quotes and \rhmcsBaShortSrPTwoMax{} at bid/ask). For the mean daily return
    differential the SPA $p$-value is \rhmcsMidShortMeanSpaC; StepM rejects \rhmcsMidShortMeanNRej{} of the
    \rhmcsNAlt{} hypotheses, all but XGBoost's (adjusted $p$-value \rhmcsMidShortMeanPXgb, at the level;
    \rhmcsMidShortMeanNRejLo{} and \rhmcsMidShortMeanNRejHi{} rejections with blocks of \rhmcsBlockLo{} and
    \rhmcsBlockHi{} days), and in the two-sided test those of the elastic net and the random forest
    (\rhmcsMidShortMeanPTwoEnet{} and \rhmcsMidShortMeanPTwoRf; \rhmcsMidShortMeanNRejTwoLo{} and
    \rhmcsMidShortMeanNRejTwoHi{} rejections with blocks of \rhmcsBlockLo{} and \rhmcsBlockHi{} days). These bootstrap
    tests reject more often than the fixed-$b$ test of the mean (Section~\ref{sec:rh_econ_measures}), which finds a
    significant gain for \rhrPoolShortSig{} of the \rhrPoolN{} pooled forecasts. The difference lies in the test, not
    in the adjustment: before any adjustment, the two-sided bootstrap $p$-value of the comparison alone is
    \rhmcsMidShortMeanRawTwoRf{} for the random forest and \rhmcsMidShortMeanRawTwoEnet{} for the elastic net, whose
    fixed-$b$ $t$-statistics are \rhrVsShortRfPoolFullT{} and \rhrVsShortEnetPoolFullT{} against a 5\% critical value
    of \rheHACCV. The bootstrap estimates the variance of the mean from blocks of \rhmcsBlock{} days and treats it as
    known, while the fixed-$b$ test uses a bandwidth of \rheHACS{} days and critical values that allow for the error
    of the variance estimate.
  \item \emph{Against the HAR-benchmark strategy of the same class.} The Sharpe ratio gain survives the adjustment
    for the random forest alone: SPA $p$-value \rhmcsMidTwinSrSpaC, StepM-adjusted $p$-value \rhmcsMidTwinSrPRf{}
    (two-sided \rhmcsMidTwinSrPTwoRf, against \rhmcsTwinTwoSidedPRf{} for the comparison alone), with the same
    conclusion at quoted bid/ask prices (two-sided \rhmcsBaTwinSrPTwoRf). It is the one significant Sharpe ratio gain
    of the item above; the next smallest single-test $p$-value, LightGBM's \rhmcsTwinTwoSidedPLgbm, becomes
    \rhmcsMidTwinSrPTwoLgbm{} after the adjustment. With blocks of \rhmcsBlockLo{} and \rhmcsBlockHi{} days the
    random forest's two-sided adjusted $p$-value is \rhmcsMidTwinSrPTwoMinLo{} and \rhmcsMidTwinSrPTwoMinHi{}
    (one-sided \rhmcsMidTwinSrPMinLo{} and \rhmcsMidTwinSrPMinHi). For the mean daily return differential, StepM
    rejects for \rhmcsMidTwinMeanRej{} (adjusted $p$-values \rhmcsMidTwinMeanPLgbm{} and \rhmcsMidTwinMeanPRf), and in
    the two-sided test for the random forest alone (\rhmcsMidTwinMeanPTwoRf; \rhmcsMidTwinMeanPTwoMinLo{} and
    \rhmcsMidTwinMeanPTwoMinHi{} with blocks of \rhmcsBlockLo{} and \rhmcsBlockHi{} days).
  \end{itemize}

% =====================================================================================================================
% Part D. Sentences of the current text that the procedures qualify (suggestions; none of them is edited here).
% =====================================================================================================================
% D1  Introduction ("Main findings", first item) and Section 6 (sec:rh_results_predictors, the penalized linear models):
%     "raises it for the three penalized linear models, significantly so for ridge regression and the lasso" / "at the
%     5% level the increase is statistically significant for ridge regression and the lasso". After the two-sided
%     StepM over the six classes only the lasso's increase survives (\rhmcsOosQAfPTwoLasso; ridge \rhmcsOosQAfPTwoRidge).
%     Suggested: "..., significantly so for the lasso, and for ridge regression before the adjustment for the six
%     comparisons (Table~\ref{tab:rh_mcs_spa})".
% D2  Section 6 (sec:rh_results_classes, "Under the HAR benchmark ... significant at the 5% level for \rhrVsBmPoolHarSig
%     of the \rhrAfPoolPairs models and at the 1% level for \rhrVsBmPoolHarSigOne"): with the ten distinct pooled
%     forecasts tested against the benchmark at once, the higher loss of LightGBM and XGBoost on the HAR benchmark does
%     not survive (\rhmcsOosQBmAllPTwoLgbmHar, \rhmcsOosQBmAllPTwoXgbHar), and both stay in the 90% model confidence
%     set with the benchmark. The linear HAR forecasts (\rhmcsOosQBmAllPTwoLinHar) and the random forest on the HAR
%     benchmark (\rhmcsOosQBmAllPTwoRfHar) survive.
% D3  Introduction, Section 6 and Conclusion ("on the full predictor set the gap is statistically significant at the
%     5% level for \rhrVsBmPoolLinFullSig of the three penalized linear models and for \rhrVsBmPoolTreeFullSig of the
%     three tree ensembles"): the linear models' gap survives the adjustment; the one tree ensemble's (the random
%     forest, pairwise \rhrVsBmRfPoolFullPval) does not (\rhmcsOosQBmAllPTwoRfFull), and it stays in the 90% set under
%     T_max (MCS p-value \rhmcsOosQlikeTmaxPRfFull; under T_R \rhmcsOosQlikeTrangePRfFull).
% D4  Abstract, Introduction, Section 6, Conclusion ("every intraday specification remains less accurate under QLIKE
%     than" the closing-interval benchmark): true of the average losses, and the benchmark is the last forecast left in
%     the model confidence set; but the boosted trees on both predictor sets are in its 90% set (and XGBoost on the full
%     set in its 75% set), so "less accurate" should stay a statement about point estimates.
% D5  Section 7 (sec:rh_econ_results, "the mean return gain over the short-straddle benchmark is statistically
%     significant at the 5% level in \rhrPoolShortSig of the \rhrPoolN specifications") and the Conclusion (comparison
%     with Maisch et al.: "its gain over that benchmark is statistically significant in \rhrPoolShortSig of the
%     \rhrPoolN pooled forecasts"): the Conclusion should say "its mean return gain"; the Sharpe ratio gain over the
%     short straddle, which the paper reports only in point estimates, is significant for all six full-set strategies
%     under the paper's own D6 bootstrap, after the StepM adjustment (Part C). The mean-return results differ by test
%     (fixed-b against the block bootstrap), not by the adjustment (Part C).
% D6  Abstract, Introduction, Section 7, Conclusion ("the gain is statistically significant at the 5% level in
%     \rhrAfPoolShHigherSig of them"): survives the adjustment for the six classes (the random forest, Part C), one-sided
%     at the three block lengths and two-sided at \rhmcsBlockLo{} and \rhmcsBlock{} days (\rhmcsMidTwinSrPTwoMinHi{}
%     with blocks of \rhmcsBlockHi{} days).
% D7  Every-bar claims (Abstract, Introduction, Section 6, Conclusion: the full predictor set lowers QLIKE in all six
%     classes, each significantly): survive (Part B).
% D8  Appendix sec:rh_methods_tests, item "Reporting convention" ("All main tables report point estimates ...
%     without bootstrap confidence intervals"): still true of the main tables; add "Tables~\ref{tab:rh_mcs}
%     and~\ref{tab:rh_mcs_spa} report the multiple-comparison procedures."
